\input{preamble}

% OK, start here.
%
\begin{document}

\title{Moduli Stacks}

\maketitle

\phantomsection
\label{section-phantom}

\tableofcontents




\section{Introduction}
\label{section-introduction}

\noindent
In this chapter we verify basic properties of moduli spaces
and moduli stacks such as
$\mathit{Hom}$, $\mathit{Isom}$, $\Cohstack_{X/B}$,
$\Quotfunctor_{\mathcal{F}/X/B}$, $\Hilbfunctor_{X/B}$,
$\Picardstack_{X/B}$, $\Picardfunctor_{X/B}$, $\mathit{Mor}_B(Z, X)$,
$\Spacesstack'_{fp, flat, proper}$, $\Polarizedstack$, and
$\Complexesstack_{X/B}$.
We have already shown these algebraic spaces or algebraic stacks
under suitable hypotheses, see Quot, Sections
\ref{quot-section-hom},
\ref{quot-section-isom},
\ref{quot-section-stack-coherent-sheaves},
\ref{quot-section-not-flat},
\ref{quot-section-quot},
\ref{quot-section-hilb},
\ref{quot-section-picard-stack},
\ref{quot-section-picard-functor},
\ref{quot-section-relative-morphisms},
\ref{quot-section-stack-of-spaces},
\ref{quot-section-polarized}, and
\ref{quot-section-moduli-complexes}.
The stack of curves, denoted $\textit{Curves}$ and introduced
in Quot, Section \ref{quot-section-curves}, is discussed in the
chapter on moduli of curves, see
Moduli of Curves, Section \ref{moduli-curves-section-stack-curves}.

\medskip\noindent
In some sense this chapter is following the footsteps of
Grothendieck's lectures \cite{Gr-I},
\cite{Gr-II},
\cite{Gr-III},
\cite{Gr-IV},
\cite{Gr-V}, and
\cite{Gr-VI}.







\section{Conventions and abuse of language}
\label{section-conventions}

\noindent
We continue to use the conventions and the abuse of language
introduced in
Properties of Stacks, Section \ref{stacks-properties-section-conventions}.
Unless otherwise mentioned our base scheme will be $\Spec(\mathbf{Z})$.






\section{Properties of Hom and Isom}
\label{section-hom-isom}

\noindent
Let $f : X \to B$ be a morphism of algebraic spaces which is
of finite presentation. Assume $\mathcal{F}$ and $\mathcal{G}$
are quasi-coherent $\mathcal{O}_X$-modules.
If $\mathcal{G}$ is of finite presentation, flat over $B$
with support proper over $B$, then the functor
$\mathit{Hom}(\mathcal{F}, \mathcal{G})$ defined by
$$
T/B \longmapsto \Hom_{\mathcal{O}_{X_T}}(\mathcal{F}_T, \mathcal{G}_T)
$$
is an algebraic space affine over $B$. If $\mathcal{F}$ is of
finite presentation, then
$\mathit{Hom}(\mathcal{F}, \mathcal{G}) \to B$
is of finite presentation. See
Quot, Proposition \ref{quot-proposition-hom}.

\medskip\noindent
If both $\mathcal{F}$ and $\mathcal{G}$ are of finite presentation,
flat over $B$ with support proper over $B$, then the subfunctor
$$
\mathit{Isom}(\mathcal{F}, \mathcal{G}) \subset
\mathit{Hom}(\mathcal{F}, \mathcal{G})
$$
is an algebraic space affine of finite presentation over $B$.
See Quot, Proposition \ref{quot-proposition-isom}.




\section{Properties of the stack of coherent sheaves}
\label{section-stack-coherent-sheaves}

\noindent
Let $f : X \to B$ be a morphism of algebraic spaces which is
separated and of finite presentation. Then the stack
$\Cohstack_{X/B}$ parametrizing flat families of coherent
modules with proper support is algebraic. See
Quot, Theorem \ref{quot-theorem-coherent-algebraic-general}.

\begin{lemma}
\label{lemma-coherent-diagonal-affine-fp}
The diagonal of $\Cohstack_{X/B}$ over $B$ is affine
and of finite presentation.
\end{lemma}

\begin{proof}
The representability of the diagonal by algebraic spaces
was shown in Quot, Lemma \ref{quot-lemma-coherent-diagonal}.
From the proof we find that we have to show
$\mathit{Isom}(\mathcal{F}, \mathcal{G}) \to T$
is affine and of finite presentation for a pair of
finitely presented $\mathcal{O}_{X_T}$-modules
$\mathcal{F}$, $\mathcal{G}$ flat over $T$ with support
proper over $T$. This was discussed in Section \ref{section-hom-isom}.
\end{proof}

\begin{lemma}
\label{lemma-coherent-qs-lfp}
The morphism $\Cohstack_{X/B} \to B$ is quasi-separated and
locally of finite presentation.
\end{lemma}

\begin{proof}
To check $\Cohstack_{X/B} \to B$ is quasi-separated we have to
show that its diagonal is quasi-compact and quasi-separated.
This is immediate from Lemma \ref{lemma-coherent-diagonal-affine-fp}.
To prove that $\Cohstack_{X/B} \to B$ is locally of finite
presentation, we have to show that $\Cohstack_{X/B} \to B$
is limit preserving, see
Limits of Stacks, Proposition
\ref{stacks-limits-proposition-characterize-locally-finite-presentation}.
This follows from Quot, Lemma \ref{quot-lemma-coherent-limits}
(small detail omitted).
\end{proof}

\begin{lemma}
\label{lemma-coherent-existence-part}
Assume $X \to B$ is proper as well as of finite presentation.
Then $\Cohstack_{X/B} \to B$ satisfies the existence part
of the valuative criterion (Morphisms of Stacks, Definition
\ref{stacks-morphisms-definition-existence}).
\end{lemma}

\begin{proof}
Taking base change, this immediately reduces to the following
problem: given a valuation ring $R$ with fraction field $K$ and
an algebraic space $X$ proper over $R$ and a coherent
$\mathcal{O}_{X_K}$-module $\mathcal{F}_K$, show there exists
a finitely presented $\mathcal{O}_X$-module $\mathcal{F}$
flat over $R$ whose generic fibre is $\mathcal{F}_K$.
Observe that by Flatness on Spaces, Theorem
\ref{spaces-flat-theorem-finite-type-flat}
any finite type quasi-coherent $\mathcal{O}_X$-module
$\mathcal{F}$ flat over $R$ is of finite presentation.
Denote $j : X_K \to X$ the embedding of the generic fibre.
As a base change of the affine morphism $\Spec(K) \to \Spec(R)$
the morphism $j$ is affine. Thus $j_*\mathcal{F}_K$ is
quasi-coherent. Write
$$
j_*\mathcal{F}_K = \colim \mathcal{F}_i
$$
as a filtered colimit of its finite type quasi-coherent
$\mathcal{O}_X$-submodules, see
Limits of Spaces, Lemma \ref{spaces-limits-lemma-directed-colimit-finite-type}.
Since $j_*\mathcal{F}_K$ is a sheaf of $K$-vector spaces over $X$,
it is flat over $\Spec(R)$. Thus each $\mathcal{F}_i$ is flat
over $R$ as flatness over a valuation ring is the same as being
torsion free
(More on Algebra, Lemma
\ref{more-algebra-lemma-valuation-ring-torsion-free-flat})
and torsion freeness is inherited by submodules.
Finally, we have to show that the map
$j^*\mathcal{F}_i \to \mathcal{F}_K$
is an isomorphism for some $i$.
Since $j^*j_*\mathcal{F}_K = \mathcal{F}_K$ (small detail omitted)
and since $j^*$ is exact, we see that $j^*\mathcal{F}_i \to \mathcal{F}_K$
is injective for all $i$.
Since $j^*$ commutes with colimits, we have
$\mathcal{F}_K = j^*j_*\mathcal{F}_K = \colim j^*\mathcal{F}_i$.
Since $\mathcal{F}_K$ is coherent (i.e., finitely presented),
there is an $i$ such that $j^*\mathcal{F}_i$ contains all the
(finitely many) generators over an affine \'etale cover of $X$.
Thus we get surjectivity of $j^*\mathcal{F}_i \to \mathcal{F}_K$
for $i$ large enough.
\end{proof}

\begin{lemma}
\label{lemma-coherent-functorial}
Let $B$ be an algebraic space. Let $\pi : X \to Y$ be a quasi-finite
morphism of algebraic spaces which are separated and of finite presentation
over $B$. Then $\pi_*$ induces a morphism
$\Cohstack_{X/B} \to \Cohstack_{Y/B}$.
\end{lemma}

\begin{proof}
Let $(T \to B, \mathcal{F})$ be an object of $\Cohstack_{X/B}$.
We claim
\begin{enumerate}
\item[(a)] $(T \to B, \pi_{T, *}\mathcal{F})$ is an object
of $\Cohstack_{Y/B}$ and
\item[(b)] for $T' \to T$ we have
$\pi_{T', *}(X_{T'} \to X_T)^*\mathcal{F} =
(Y_{T'} \to Y_T)^*\pi_{T, *}\mathcal{F}$.
\end{enumerate}
Part (b) guarantees that this construction defines a functor
$\Cohstack_{X/B} \to \Cohstack_{Y/B}$ as desired.

\medskip\noindent
Let $i : Z \to X_T$ be the closed subspace cut out by the zeroth
fitting ideal of $\mathcal{F}$
(Divisors on Spaces, Section
\ref{spaces-divisors-section-fitting-ideals}).
Then $Z \to B$ is proper by assumption (see
Derived Categories of Spaces, Section
\ref{spaces-perfect-section-proper-over-base}).
On the other hand $i$ is of finite presentation
(Divisors on Spaces, Lemma
\ref{spaces-divisors-lemma-fitting-ideal-of-finitely-presented} and
Morphisms of Spaces, Lemma
\ref{spaces-morphisms-lemma-closed-immersion-finite-presentation}).
There exists a quasi-coherent $\mathcal{O}_Z$-module
$\mathcal{G}$ of finite type with $i_*\mathcal{G} = \mathcal{F}$
(Divisors on Spaces, Lemma
\ref{spaces-divisors-lemma-on-subscheme-cut-out-by-Fit-0}).
In fact $\mathcal{G}$ is of finite presentation as an $\mathcal{O}_Z$-module
by Descent on Spaces, Lemma
\ref{spaces-descent-lemma-finite-finitely-presented-module}.
Observe that $\mathcal{G}$ is flat over $B$, for example
because the stalks of $\mathcal{G}$ and $\mathcal{F}$ agree
(Morphisms of Spaces, Lemma \ref{spaces-morphisms-lemma-stalk-push-closed}).
Observe that $\pi_T \circ i : Z \to Y_T$ is quasi-finite as a composition
of quasi-finite morphisms and that
$\pi_{T, *}\mathcal{F} = (\pi_T \circ i)_*\mathcal{G})$.
Since $i$ is affine, formation of $i_*$ commutes with base change
(Cohomology of Spaces, Lemma \ref{spaces-cohomology-lemma-affine-base-change}).
Therefore we may replace $B$ by $T$, $X$ by $Z$,
$\mathcal{F}$ by $\mathcal{G}$, and $Y$ by $Y_T$
to reduce to the case discussed in the next paragraph.

\medskip\noindent
Assume that $X \to B$ is proper. Then $\pi$ is proper
by Morphisms of Spaces, Lemma
\ref{spaces-morphisms-lemma-universally-closed-permanence}
and hence finite by
More on Morphisms of Spaces,
Lemma \ref{spaces-more-morphisms-lemma-characterize-finite}.
Since a finite morphism is affine we see that (b) holds by
Cohomology of Spaces, Lemma \ref{spaces-cohomology-lemma-affine-base-change}.
On the other hand, $\pi$ is of finite presentation by
Morphisms of Spaces, Lemma
\ref{spaces-morphisms-lemma-finite-presentation-permanence}.
Thus $\pi_{T, *}\mathcal{F}$ is of finite presentation by
Descent on Spaces, Lemma
\ref{spaces-descent-lemma-finite-finitely-presented-module}.
Finally, $\pi_{T, *}\mathcal{F} $ is flat over $B$ for example
by looking at stalks using
Cohomology of Spaces, Lemma \ref{spaces-cohomology-lemma-stalk-push-finite}.
\end{proof}

\begin{lemma}
\label{lemma-coherent-open}
Let $B$ be an algebraic space. Let $\pi : X \to Y$ be an open immersion
of algebraic spaces which are separated and of finite presentation over $B$.
Then the morphism $\Cohstack_{X/B} \to \Cohstack_{Y/B}$ of
Lemma \ref{lemma-coherent-functorial} is an open immersion.
\end{lemma}

\begin{proof}
Omitted. Hint: If $\mathcal{F}$ is an object of $\Cohstack_{Y/B}$ over $T$
and for $t \in T$ we have $\text{Supp}(\mathcal{F}_t) \subset |X_t|$,
then the same is true for $t' \in T$ in a neighbourhood of $t$.
\end{proof}

\begin{lemma}
\label{lemma-coherent-closed}
Let $B$ be an algebraic space. Let $\pi : X \to Y$ be a closed immersion
of algebraic spaces which are separated and of finite presentation over $B$.
Then the morphism $\Cohstack_{X/B} \to \Cohstack_{Y/B}$ of
Lemma \ref{lemma-coherent-functorial} is a closed immersion.
\end{lemma}

\begin{proof}
Let $\mathcal{I} \subset \mathcal{O}_Y$ be the sheaf of ideals cutting
out $X$ as a closed subspace of $Y$. Recall that $\pi_*$ induces
an equivalence between the category of quasi-coherent $\mathcal{O}_X$-modules
and the category of quasi-coherent $\mathcal{O}_Y$-modules annihilated
by $\mathcal{I}$, see Morphisms of Spaces, Lemma
\ref{spaces-morphisms-lemma-i-star-equivalence}.
The same, mutatis mutandis, is true after base by $T \to B$ with
$\mathcal{I}$ replaced by the ideal sheaf
$\mathcal{I}_T = \Im((Y_T \to Y)^*\mathcal{I} \to \mathcal{O}_{Y_T})$.
Analyzing the proof of Lemma \ref{lemma-coherent-functorial}
we find that the essential image of
$\Cohstack_{X/B} \to \Cohstack_{Y/B}$
is exactly the objects $\xi = (T \to B, \mathcal{F})$
where $\mathcal{F}$ is annihilated by $\mathcal{I}_T$.
In other words, $\xi$ is in the essential image if and only if
the multiplication map
$$
\mathcal{F} \otimes_{\mathcal{O}_{Y_T}} (Y_T \to Y)^*\mathcal{I}
\longrightarrow
\mathcal{F}
$$
is zero and similarly after any further base change $T' \to T$.
Note that
$$
(Y_{T'} \to Y_T)^*(
\mathcal{F} \otimes_{\mathcal{O}_{Y_T}} (Y_T \to Y)^*\mathcal{I}) =
(Y_{T'} \to Y_T)^*\mathcal{F} \otimes_{\mathcal{O}_{Y_{T'}}}
(Y_{T'} \to Y)^*\mathcal{I})
$$
Hence the vanishing of the multiplication map on $T'$
is representable by a closed subspace of $T$ by
Flatness on Spaces, Lemma \ref{spaces-flat-lemma-F-zero-closed-proper}.
\end{proof}

\begin{situation}[Numerical invariants]
\label{situation-numerical}
Let $f : X \to B$ be as in the introduction to this section. Let $I$
be a set and for $i \in I$ let $E_i \in D(\mathcal{O}_X)$ be perfect.
Given an object $(T \to B, \mathcal{F})$ of $\Cohstack_{X/B}$
denote $E_{i, T}$ the derived pullback of $E_i$ to $X_T$.
The object
$$
K_i = Rf_{T, *}(E_{i, T} \otimes_{\mathcal{O}_{X_T}}^\mathbf{L} \mathcal{F})
$$
of $D(\mathcal{O}_T)$ is perfect and its formation commutes with base change,
see Derived Categories of Spaces, Lemma
\ref{spaces-perfect-lemma-base-change-tensor-perfect}.
Thus the function
$$
\chi_i : |T| \longrightarrow \mathbf{Z},\quad
\chi_i(t) =
\chi(X_t, E_{i, t} \otimes_{\mathcal{O}_{X_t}}^\mathbf{L} \mathcal{F}_t) =
\chi(K_i \otimes_{\mathcal{O}_T}^\mathbf{L} \kappa(t))
$$
is locally constant by Derived Categories of Spaces, Lemma
\ref{spaces-perfect-lemma-chi-locally-constant}.
Let $P : I \to \mathbf{Z}$ be a map. Consider the substack
$$
\Cohstack^P_{X/B} \subset \Cohstack_{X/B}
$$
consisting of flat families of coherent sheaves with proper support
whose numerical invariants agree with $P$. More precisely, an object
$(T \to B, \mathcal{F})$ of $\Cohstack_{X/B}$ is in
$\Cohstack^P_{X/B}$ if and only if $\chi_i(t) = P(i)$ for all $i \in I$
and $t \in T$.
\end{situation}

\begin{lemma}
\label{lemma-open-P}
In Situation \ref{situation-numerical} the stack
$\Cohstack^P_{X/B}$ is algebraic and
$$
\Cohstack^P_{X/B} \longrightarrow \Cohstack_{X/B}
$$
is a flat closed immersion. If $I$ is finite or $B$ is locally
Noetherian, then $\Cohstack^P_{X/B}$ is an open and closed substack of
$\Cohstack_{X/B}$.
\end{lemma}

\begin{proof}
This is immediately clear if $I$ is finite, because the functions
$t \mapsto \chi_i(t)$ are locally constant. If $I$ is infinite, then
we write
$$
I = \bigcup\nolimits_{I' \subset I\text{ finite}} I'
$$
and we denote $P' = P|_{I'}$. Then we have
$$
\Cohstack^P_{X/B} = \bigcap\nolimits_{I' \subset I\text{ finite}}
\Cohstack^{P'}_{X/B}
$$
Therefore, $\Cohstack^P_{X/B}$ is always an algebraic stack and the morphism
$\Cohstack^P_{X/B} \subset \Cohstack_{X/B}$ is always a flat closed immersion,
but it may no longer be an open substack. (We leave it to the reader to
make examples). However, if $B$ is locally Noetherian, then so
is $\Cohstack_{X/B}$ by Lemma \ref{lemma-coherent-qs-lfp} and
Morphisms of Stacks, Lemma
\ref{stacks-morphisms-lemma-locally-finite-type-locally-noetherian}.
Hence if $U \to \Cohstack_{X/B}$ is a smooth surjective morphism
where $U$ is a locally Noetherian scheme, then the inverse images of
the open and closed substacks $\Cohstack^{P'}_{X/B}$
have an open intersection in $U$ (because connected components of
locally Noetherian topological spaces are open).
Thus the result in this case.
\end{proof}

\begin{lemma}
\label{lemma-finite-list-perfect-objects}
Let $f : X \to B$ be as in the introduction to this section.
Let $E_1, \ldots, E_r \in D(\mathcal{O}_X)$ be perfect.
Let $I = \mathbf{Z}^{\oplus r}$ and consider the map
$$
I \longrightarrow D(\mathcal{O}_X),\quad
(n_1, \ldots, n_r) \longmapsto
E_1^{\otimes n_1}
\otimes \ldots \otimes
E_r^{\otimes n_r}
$$
Let $P : I \to \mathbf{Z}$ be a map. Then
$\Cohstack^P_{X/B} \subset \Cohstack_{X/B}$
as defined in Situation \ref{situation-numerical}
is an open and closed substack.
\end{lemma}

\begin{proof}
We may work \'etale locally on $B$, hence we may assume that $B$ is affine.
In this case we may perform absolute Noetherian reduction; we suggest
the reader skip the proof. Namely, say $B = \Spec(\Lambda)$.
Write $\Lambda = \colim \Lambda_i$ as a filtered colimit with each $\Lambda_i$
of finite type over $\mathbf{Z}$. For some $i$ we can find
a morphism of algebraic spaces $X_i \to \Spec(\Lambda_i)$
which is separated and of finite presentation and whose base change
to $\Lambda$ is $X$. See Limits of Spaces, Lemmas
\ref{spaces-limits-lemma-descend-finite-presentation} and
\ref{spaces-limits-lemma-descend-separated-morphism}.
Then after increasing $i$ we may assume there exist
perfect objects $E_{1, i}, \ldots, E_{r, i}$
in $D(\mathcal{O}_{X_i})$ whose derived pullback to $X$
are isomorphic to $E_1, \ldots, E_r$, see
Derived Categories of Spaces, Lemma
\ref{spaces-perfect-lemma-perfect-on-limit}.
Clearly we have a cartesian square
$$
\xymatrix{
\Cohstack^P_{X/B} \ar[r] \ar[d] &
\Cohstack_{X/B} \ar[d] \\
\Cohstack^P_{X_i/\Spec(\Lambda_i)} \ar[r] &
\Cohstack_{X_i/\Spec(\Lambda_i)}
}
$$
and hence we may appeal to Lemma \ref{lemma-open-P}
to finish the proof.
\end{proof}

\begin{example}[Coherent sheaves with fixed Hilbert polynomial]
\label{example-hilbert-polynomial}
Let $f : X \to B$ be as in the introduction to this section.
Let $\mathcal{L}$ be an invertible $\mathcal{O}_X$-module.
Let $P : \mathbf{Z} \to \mathbf{Z}$ be a numerical polynomial.
Then we can consider the open and closed algebraic substack
$$
\Cohstack^P_{X/B} =
\Cohstack^{P, \mathcal{L}}_{X/B}
\subset \Cohstack_{X/B}
$$
consisting of flat families of coherent sheaves with proper support
whose numerical invariants agree with $P$: an object
$(T \to B, \mathcal{F})$ of $\Cohstack_{X/B}$ lies in
$\Cohstack^P_{X/B}$ if and only if
$$
P(n) =
\chi(X_t, \mathcal{F}_t \otimes_{\mathcal{O}_{X_t}} \mathcal{L}_t^{\otimes n})
$$
for all $n \in \mathbf{Z}$ and $t \in T$. Of course this is a
special case of Situation \ref{situation-numerical}
where $I = \mathbf{Z} \to D(\mathcal{O}_X)$ is given by
$n \mapsto \mathcal{L}^{\otimes n}$. It follows from
Lemma \ref{lemma-finite-list-perfect-objects}
that this is an open and closed substack. Since the functions
$n \mapsto
\chi(X_t, \mathcal{F}_t \otimes_{\mathcal{O}_{X_t}} \mathcal{L}_t^{\otimes n})$
are always numerical polynomials (Spaces over Fields, Lemma
\ref{spaces-over-fields-lemma-numerical-polynomial-from-euler})
we conclude that
$$
\Cohstack_{X/B} = \coprod\nolimits_{P\text{ numerical polynomial}}
\Cohstack^P_{X/B}
$$
is a disjoint union decomposition.
\end{example}







\section{Properties of Quot}
\label{section-quot}

\noindent
Let $f : X \to B$ be a morphism of algebraic spaces which is
separated and of finite presentation. Let $\mathcal{F}$ be a
quasi-coherent $\mathcal{O}_X$-module. Then
$\Quotfunctor_{\mathcal{F}/X/B}$ is an algebraic space.
If $\mathcal{F}$ is of finite presentation, then
$\Quotfunctor_{\mathcal{F}/X/B} \to B$ is locally of finite
presentation. See Quot, Proposition \ref{quot-proposition-quot}.

\begin{lemma}
\label{lemma-quot-diagonal-closed}
The diagonal of $\Quotfunctor_{\mathcal{F}/X/B} \to B$ is a closed immersion.
If $\mathcal{F}$ is of finite type, then the diagonal is a closed
immersion of finite presentation.
\end{lemma}

\begin{proof}
Suppose we have a scheme $T/B$ and two quotients
$\mathcal{F}_T \to \mathcal{Q}_i$, $i = 1, 2$ corresponding
to $T$-valued points of $\Quotfunctor_{\mathcal{F}/X/B}$ over $B$.
Denote $\mathcal{K}_1$ the kernel of the first one and set
$u : \mathcal{K}_1 \to \mathcal{Q}_2$ the composition.
By Flatness on Spaces, Lemma \ref{spaces-flat-lemma-F-zero-closed-proper}
there is a closed subspace of $T$ such that $T' \to T$
factors through it if and only if the pullback $u_{T'}$ is zero.
This proves the diagonal is a closed immersion.
Moreover, if $\mathcal{F}$ is of finite type, then
$\mathcal{K}_1$ is of finite type
(Modules on Sites, Lemma
\ref{sites-modules-lemma-kernel-surjection-finite-onto-finite-presentation})
and we see that the diagonal is of finite presentation by
the same lemma.
\end{proof}

\begin{lemma}
\label{lemma-quot-s-lfp}
The morphism $\Quotfunctor_{\mathcal{F}/X/B} \to B$ is separated.
If $\mathcal{F}$ is of finite presentation, then it is also
locally of finite presentation.
\end{lemma}

\begin{proof}
To check $\Quotfunctor_{\mathcal{F}/X/B} \to B$ is separated we have to
show that its diagonal is a closed immersion. This
is true by Lemma \ref{lemma-quot-diagonal-closed}.
The second statement is part of
Quot, Proposition \ref{quot-proposition-quot}.
\end{proof}

\begin{lemma}
\label{lemma-quotient-extension-valuation-ring}
\begin{reference}
\cite[Expos\'e 221, Lemma 3.7, pp. 264--265]{FGA}.
\end{reference}
Let $R$ be a valuation ring with fraction field $K$, let $X$ be an
algebraic space over $R$, and let $\mathcal{F}$ be a quasi-coherent
$\mathcal{O}_X$-module. Every quasi-coherent quotient
$$
\mathcal{F}_K \longrightarrow \mathcal{Q}_K
$$
extends uniquely, up to unique isomorphism compatible with the quotient,
to a quasi-coherent quotient
$$
\mathcal{F} \longrightarrow \mathcal{Q}
$$
such that $\mathcal{Q}$ is flat over $R$.
\end{lemma}

\begin{proof}
Write $j : X_K \to X$. Define $\mathcal{L}$ to be the kernel of the
composition
$$
\mathcal{F} \longrightarrow j_*\mathcal{F}_K
\longrightarrow j_*\mathcal{Q}_K
$$
and set $\mathcal{Q} = \mathcal{F}/\mathcal{L}$. These modules are
quasi-coherent because $j$ is affine. On an affine scheme over $X$ this
construction is the following. If $M$ is the module corresponding to
$\mathcal{F}$ and $M_K \to N_K$ is the given quotient, then
$$
L = \{x \in M \mid x/1 \text{ maps to }0\text{ in }N_K\}
$$
and $\mathcal{Q}$ corresponds to $M/L$. The module $M/L$ embeds into
$N_K$ and hence is torsion free as an $R$-module. It is flat over $R$ by
More on Algebra, Lemma
\ref{more-algebra-lemma-valuation-ring-torsion-free-flat}.

Suppose $\mathcal{F} \to \mathcal{Q}'$ is another extension flat over
$R$. On an affine scheme as above, flatness implies that the module
corresponding to $\mathcal{Q}'$ is $R$-torsion free. Its kernel is
therefore exactly $L$: if the image of $x$ vanishes after tensoring with
$K$, then it is annihilated by a nonzero element of $R$, hence is zero.
This proves uniqueness and completes the proof.
\end{proof}

\begin{lemma}
\label{lemma-quot-existence-part}
Assume $X \to B$ is proper as well as of finite presentation
and $\mathcal{F}$ quasi-coherent of finite type.
Then $\Quotfunctor_{\mathcal{F}/X/B} \to B$ satisfies the existence part
of the valuative criterion (Morphisms of Spaces, Definition
\ref{spaces-morphisms-definition-valuative-criterion}).
\end{lemma}

\begin{proof}
Taking base change, this immediately reduces to the following
problem: given a valuation ring $R$ with fraction field $K$,
an algebraic space $X$ proper over $R$, a finite type quasi-coherent
$\mathcal{O}_X$-module $\mathcal{F}$, and a coherent
quotient $\mathcal{F}_K \to \mathcal{Q}_K$, show there exists
a quotient $\mathcal{F} \to \mathcal{Q}$ where $\mathcal{Q}$ is a
finitely presented $\mathcal{O}_X$-module
flat over $R$ whose generic fibre is $\mathcal{Q}_K$.
Apply Lemma \ref{lemma-quotient-extension-valuation-ring}. The resulting
$\mathcal{Q}$ is of finite type because it is a quotient of $\mathcal{F}$.
By Flatness on Spaces, Theorem
\ref{spaces-flat-theorem-finite-type-flat}
it is of finite presentation. Its support is closed in $X$ and hence
proper over $R$. Thus this quotient is an object of the Quot functor
and gives the required extension.
\end{proof}

\begin{lemma}
\label{lemma-quot-functorial}
Let $B$ be an algebraic space. Let $\pi : X \to Y$ be an affine quasi-finite
morphism of algebraic spaces which are separated and of finite presentation
over $B$. Let $\mathcal{F}$ be a quasi-coherent $\mathcal{O}_X$-module.
Then $\pi_*$ induces a morphism
$\Quotfunctor_{\mathcal{F}/X/B} \to \Quotfunctor_{\pi_*\mathcal{F}/Y/B}$.
\end{lemma}

\begin{proof}
Set $\mathcal{G} = \pi_*\mathcal{F}$. Since $\pi$ is affine we see that for
any scheme $T$ over $B$ we have $\mathcal{G}_T = \pi_{T, *}\mathcal{F}_T$ by
Cohomology of Spaces, Lemma \ref{spaces-cohomology-lemma-affine-base-change}.
Moreover $\pi_T$ is affine, hence $\pi_{T, *}$ is exact and transforms
quotients into quotients. Observe that a quasi-coherent quotient
$\mathcal{F}_T \to \mathcal{Q}$ defines a point of $\Quotfunctor_{X/B}$
if and only if $\mathcal{Q}$ defines an object of $\Cohstack_{X/B}$
over $T$ (similarly for $\mathcal{G}$ and $Y$). Since we've seen in
Lemma \ref{lemma-coherent-functorial}
that $\pi_*$ induces a morphism $\Cohstack_{X/B} \to \Cohstack_{Y/B}$
we see that if $\mathcal{F}_T \to \mathcal{Q}$ is in
$\Quotfunctor_{\mathcal{F}/X/B}(T)$, then
$\mathcal{G}_T \to \pi_{T, *}\mathcal{Q}$ is
in $\Quotfunctor_{\mathcal{G}/Y/B}(T)$.
\end{proof}

\begin{lemma}
\label{lemma-quot-open}
Let $B$ be an algebraic space. Let $\pi : X \to Y$ be an affine open immersion
of algebraic spaces which are separated and of finite presentation over $B$.
Let $\mathcal{F}$ be a quasi-coherent $\mathcal{O}_X$-module. Then the morphism
$\Quotfunctor_{\mathcal{F}/X/B} \to \Quotfunctor_{\pi_*\mathcal{F}/Y/B}$ of
Lemma \ref{lemma-quot-functorial} is an open immersion.
\end{lemma}

\begin{proof}
Omitted. Hint: If $(\pi_*\mathcal{F})_T \to \mathcal{Q}$ is an element of
$\Quotfunctor_{\pi_*\mathcal{F}/Y/B}(T)$
and for $t \in T$ we have $\text{Supp}(\mathcal{Q}_t) \subset |X_t|$,
then the same is true for $t' \in T$ in a neighbourhood of $t$.
\end{proof}

\begin{lemma}
\label{lemma-quot-better-open}
Let $B$ be an algebraic space. Let $j : X \to Y$ be an open immersion
of algebraic spaces which are separated and of finite presentation over $B$.
Let $\mathcal{G}$ be a quasi-coherent $\mathcal{O}_Y$-module and set
$\mathcal{F} = j^*\mathcal{G}$. Then there is an open immersion
$$
\Quotfunctor_{\mathcal{F}/X/B}
\longrightarrow
\Quotfunctor_{\mathcal{G}/Y/B}
$$
of algebraic spaces over $B$.
\end{lemma}

\begin{proof}
If $\mathcal{F}_T \to \mathcal{Q}$ is an element of
$\Quotfunctor_{\mathcal{F}/X/B}(T)$ then we can consider
$\mathcal{G}_T \to j_{T, *}\mathcal{F}_T \to j_{T, *}\mathcal{Q}$.
Looking at stalks one finds that this is surjective.
By Lemma \ref{lemma-coherent-functorial}
we see that $j_{T, *}\mathcal{Q}$ is finitely presented, flat over $B$
with support proper over $B$. Thus we obtain a $T$-valued
point of $\Quotfunctor_{\mathcal{G}/Y/B}$.
This defines the morphism of the lemma.
We omit the proof that this is an open immersion. Hint:
If $\mathcal{G}_T \to \mathcal{Q}$ is an element of
$\Quotfunctor_{\mathcal{G}/Y/B}(T)$
and for $t \in T$ we have $\text{Supp}(\mathcal{Q}_t) \subset |X_t|$,
then the same is true for $t' \in T$ in a neighbourhood of $t$.
\end{proof}

\begin{lemma}
\label{lemma-quot-closed}
Let $B$ be an algebraic space. Let $\pi : X \to Y$ be a closed immersion
of algebraic spaces which are separated and of finite presentation over $B$.
Let $\mathcal{F}$ be a quasi-coherent $\mathcal{O}_X$-module.
Then the morphism
$\Quotfunctor_{\mathcal{F}/X/B} \to \Quotfunctor_{\pi_*\mathcal{F}/Y/B}$ of
Lemma \ref{lemma-quot-functorial} is an isomorphism.
\end{lemma}

\begin{proof}
For every scheme $T$ over $B$ the morphism $\pi_T : X_T \to Y_T$
is a closed immersion. Then $\pi_{T, *}$ is an equivalence of
categories between $\QCoh(\mathcal{O}_{X_T})$ and the full subcategory
of $\QCoh(\mathcal{O}_{Y_T})$ whose objects are those quasi-coherent
modules annihilated by the ideal sheaf of $X_T$, see
Morphisms of Spaces, Lemma \ref{spaces-morphisms-lemma-i-star-equivalence}.
Since a qotient of
$(\pi_*\mathcal{F})_T$ is annihilated by this ideal we obtain the
bijectivity of the map
$\Quotfunctor_{\mathcal{F}/X/B}(T) \to \Quotfunctor_{\pi_*\mathcal{F}/Y/B}(T)$
for all $T$ as desired.
\end{proof}

\begin{lemma}
\label{lemma-quot-quotient}
Let $X \to B$ be as in the introduction to this section. Let
$\mathcal{F} \to \mathcal{G}$ be a surjection of quasi-coherent
$\mathcal{O}_X$-modules. Then there is a canonical closed immersion
$\Quotfunctor_{\mathcal{G}/X/B} \to \Quotfunctor_{\mathcal{F}/X/B}$.
\end{lemma}

\begin{proof}
Let $\mathcal{K} = \Ker(\mathcal{F} \to \mathcal{G})$. By right
exactness of pullbacks we find that
$\mathcal{K}_T \to \mathcal{F}_T \to \mathcal{G}_T \to 0$
is an exact sequecnce for all schemes $T$ over $B$.
In particular, a quotient of $\mathcal{G}_T$
determines a quotient of $\mathcal{F}_T$ and we obtain our transformation
of functors
$\Quotfunctor_{\mathcal{G}/X/B} \to \Quotfunctor_{\mathcal{F}/X/B}$.
This transformation is a closed immersion by
Flatness on Spaces, Lemma \ref{spaces-flat-lemma-F-zero-closed-proper}.
Namely, given an element $\mathcal{F}_T \to \mathcal{Q}$ of
$\Quotfunctor_{\mathcal{F}/X/B}(T)$, then we see that the pull
back to $T'/T$ is in the image of the transformation if and
only if $\mathcal{K}_{T'} \to \mathcal{Q}_{T'}$ is zero.
\end{proof}

\begin{remark}[Numerical invariants]
\label{remark-quot-numerical}
Let $f : X \to B$ and $\mathcal{F}$ be as in the introduction to this section.
Let $I$ be a set and for $i \in I$ let $E_i \in D(\mathcal{O}_X)$ be perfect.
Let $P : I \to \mathbf{Z}$ be a function. Recall that we have a morphism
$$
\Quotfunctor_{\mathcal{F}/X/B} \longrightarrow \Cohstack_{X/B}
$$
which sends the element $\mathcal{F}_T \to \mathcal{Q}$
of $\Quotfunctor_{\mathcal{F}/X/B}(T)$ to the object $\mathcal{Q}$
of $\Cohstack_{X/B}$ over $T$, see proof of
Quot, Proposition \ref{quot-proposition-quot}. Hence we can form
the fibre product diagram
$$
\xymatrix{
\Quotfunctor^P_{\mathcal{F}/X/B} \ar[r] \ar[d] &
\Cohstack^P_{X/B} \ar[d] \\
\Quotfunctor_{\mathcal{F}/X/B} \ar[r] &
\Cohstack_{X/B}
}
$$
This is the defining diagram for the algebraic space in the
upper left corner. The left vertical arrow is a
flat closed immersion which is an open and closed immersion
for example if $I$ is finite, or $B$ is locally Noetherian, or
$I = \mathbf{Z}$ and $E_i = \mathcal{L}^{\otimes i}$ for some
invertible $\mathcal{O}_X$-module $\mathcal{L}$ (in the last
case we sometimes use the notation
$\Quotfunctor^{P, \mathcal{L}}_{\mathcal{F}/X/B}$).
See Situation \ref{situation-numerical} and
Lemmas \ref{lemma-open-P} and \ref{lemma-finite-list-perfect-objects} and
Example \ref{example-hilbert-polynomial}.
\end{remark}

\begin{lemma}
\label{lemma-quot-tensor-invertible}
Let $f : X \to B$ and $\mathcal{F}$ be as in the introduction to this section.
Let $\mathcal{L}$ be an invertible $\mathcal{O}_X$-module.
Then tensoring with $\mathcal{L}$ defines an isomorphism
$$
\Quotfunctor_{\mathcal{F}/X/B}
\longrightarrow
\Quotfunctor_{\mathcal{F} \otimes_{\mathcal{O}_X} \mathcal{L}/X/B}
$$
Given a numerical polynomial $P(t)$, then setting $P'(t) = P(t + 1)$
this map induces an isomorphism
$\Quotfunctor^P_{\mathcal{F}/X/B}
\longrightarrow
\Quotfunctor^{P'}_{\mathcal{F} \otimes_{\mathcal{O}_X} \mathcal{L}/X/B}$
of open and closed substacks.
\end{lemma}

\begin{proof}
Set $\mathcal{G} = \mathcal{F} \otimes_{\mathcal{O}_X} \mathcal{L}$.
Observe that
$\mathcal{G}_T = \mathcal{F}_T \otimes_{\mathcal{O}_{X_T}} \mathcal{L}_T$.
If $\mathcal{F}_T \to \mathcal{Q}$ is an element of
$\Quotfunctor_{\mathcal{F}/X/B}(T)$, then we send it
to the element
$\mathcal{G}_T \to \mathcal{Q} \otimes_{\mathcal{O}_{X_T}} \mathcal{L}_T$
of
$\Quotfunctor_{\mathcal{F} \otimes_{\mathcal{O}_X} \mathcal{L}/X/B}(T)$.
This is compatible with pullbacks and hence
defines a transformation of functors as desired.
Since there is an obvious inverse transformation,
it is an isomorphism. We omit the proof of the final statement.
\end{proof}

\begin{lemma}
\label{lemma-quot-power-invertible}
Let $f : X \to B$ and $\mathcal{F}$ be as in the introduction to this section.
Let $\mathcal{L}$ be an invertible $\mathcal{O}_X$-module.
Then
$$
\Quotfunctor^{P, \mathcal{L}}_{\mathcal{F}/X/B} =
\Quotfunctor^{P', \mathcal{L}^{\otimes n}}_{\mathcal{F}/X/B}
$$
where $P'(t) = P(nt)$.
\end{lemma}

\begin{proof}
Follows immediately after unwinding all the definitions.
\end{proof}


\section{Bounded families of coherent sheaves}
\label{section-bounded-coherent}

\noindent
Let $S$ be a Noetherian scheme and let $X \to S$ be a morphism of
finite type. We introduce a fibrewise notion of boundedness which is useful
when the sheaves in question do not come with a fixed parameter scheme.

\begin{definition}
\label{definition-bounded-coherent}
\begin{reference}
\cite[Expos\'e 221, D\'efinition 1.1, pp. 249--252]{FGA}.
\end{reference}
A fibrewise coherent sheaf over $X/S$ is a triple
$$
(s, K, \mathcal{F})
$$
where $s \in S$, where $K/\kappa(s)$ is a field extension, and where
$\mathcal{F}$ is a coherent $\mathcal{O}_{X_K}$-module. Two such triples
$(s, K, \mathcal{F})$ and $(s', K', \mathcal{F}')$ are equivalent if
$s=s'$ and there are $\kappa(s)$-embeddings of $K$ and $K'$ into a field
$L$ such that
$$
\mathcal{F}_L \cong \mathcal{F}'_L
$$
on $X_L$. A family of fibrewise coherent sheaves is a set of equivalence
classes of such triples.

A family $E$ is \emph{bounded} if there are a scheme $T$ of finite type
over $S$ and a coherent $\mathcal{O}_{X_T}$-module $\mathcal{G}$ with the
following property. For every $(s, K, \mathcal{F})$ whose class belongs to
$E$, there are a point $t \in T$ above $s$, a field $L$, and
$\kappa(s)$-embeddings
$$
K \longrightarrow L \longleftarrow \kappa(t)
$$
such that $\mathcal{F}_L \cong \mathcal{G}_t \otimes_{\kappa(t)} L$.
We say that $(T, \mathcal{G})$ bounds $E$.
\end{definition}

\begin{lemma}
\label{lemma-bounded-coherent-elementary}
Let $E$ and $E'$ be bounded families of fibrewise coherent sheaves over
$X/S$.
\begin{enumerate}
\item A finite union of bounded families is bounded.
\item After any base change $S' \to S$, the family obtained from $E$ by
base change is bounded over $X_{S'}/S'$.
\item The family of all tensor products
$\mathcal{F} \otimes_{\mathcal{O}_{X_K}} \mathcal{F}'$, where the two
classes belong to $E$ and $E'$ over the same field-valued point, is bounded.
\item There is a pair $(T, \mathcal{G})$ bounding $E$ for which
$\mathcal{G}$ is flat over $T$.
\end{enumerate}
\end{lemma}

\begin{proof}
For (1), take the disjoint union of finitely many parameter schemes and the
sheaf which restricts to the given sheaf on each component. For (2), if
$(T, \mathcal{G})$ bounds $E$, then
$$
(T \times_S S', \mathcal{G}_{S'})
$$
bounds the base-changed family.

Choose $(T, \mathcal{G})$ and $(T', \mathcal{G}')$ bounding $E$ and $E'$.
On $X_{T \times_S T'}$ form the tensor product of the pullbacks of
$\mathcal{G}$ and $\mathcal{G}'$. It is coherent, and its formation
commutes with arbitrary base change. Thus this sheaf, with parameter scheme
$T \times_S T'$, bounds the family in (3).

Finally, $T$ is Noetherian. More on Morphisms, Lemma
\ref{more-morphisms-lemma-generic-flatness-stratification} gives a finite
stratification of $T$ by locally closed subschemes such that the restriction
of $\mathcal{G}$ to each stratum is flat over that stratum. The disjoint
union of the strata is still of finite type over $S$, meets every point of
$T$, and carries a flat coherent sheaf whose fibre classes are the same.
This proves (4).
\end{proof}

\begin{proposition}
\label{proposition-bounded-coherent-morphisms-extensions}
\begin{reference}
\cite[Expos\'e 221, Proposition 1.2, pp. 252--253]{FGA}.
\end{reference}
Assume in addition that $X \to S$ is proper. Let $E$ and $E'$ be bounded
families of fibrewise coherent sheaves over $X/S$.
\begin{enumerate}
\item The kernels, images, and cokernels of all homomorphisms
$$
\mathcal{F} \longrightarrow \mathcal{F}'
$$
between members of $E$ and $E'$ form bounded families.
\item The middle terms of all short exact sequences
$$
0 \longrightarrow \mathcal{F}' \longrightarrow \mathcal{H}
\longrightarrow \mathcal{F} \longrightarrow 0
$$
between members of $E$ and $E'$ form a bounded family. The same holds with
$E$ and $E'$ interchanged.
\end{enumerate}
\end{proposition}

\begin{proof}
By Lemma \ref{lemma-bounded-coherent-elementary}, after replacing the two
parameter schemes by finite stratifications on which their sheaves are flat
and taking their fibre product over $S$, we may work with one scheme $T$ of
finite type over $S$ and two coherent modules $\mathcal{G}$ and
$\mathcal{G}'$ on $X_T$ which are flat over $T$ and bound $E$ and $E'$.

The functor
$$
U/T \longmapsto
\Hom_{\mathcal{O}_{X_U}}(\mathcal{G}_U, \mathcal{G}'_U)
$$
is represented by an algebraic space $H$ affine and of finite presentation
over $T$; see Quot, Proposition \ref{quot-proposition-hom}. In particular,
$H$ is a scheme of finite type over $S$. On $X_H$ there is a universal
homomorphism
$$
u : \mathcal{G}_H \longrightarrow \mathcal{G}'_H.
$$
Its kernel, image, and cokernel are coherent. Every homomorphism between a
member of $E$ and a member of $E'$ becomes a fibre of $u$ after passage to
a common field extension. There is one base-change point to address. Apply
More on Morphisms, Lemma
\ref{more-morphisms-lemma-generic-flatness-stratification} to
$\Coker(u)$ and replace $H$ by the finite disjoint union of the resulting
strata. Then $\Coker(u)$ is flat over $H$. Since $\mathcal{G}'_H$ is flat,
$\Im(u)$ is flat over $H$ as well. The two exact sequences defining
$\Im(u)$ and $\Ker(u)$ therefore remain exact after every base change.
Thus the fibres of the three coherent modules are respectively the kernel,
image, and cokernel of the fibre homomorphism, and they bound the three
families in (1).

For completeness, we recall the additional base-change step needed for
extensions. Noetherian generic flatness and cohomology and base change,
applied to the relative derived Hom of $\mathcal{G}'$ into $\mathcal{G}$,
give a finite stratification of $T$ such that on each stratum the relative
$\Ext^1$ functor commutes with arbitrary base change and is the functor of
sections of a finite locally free module. This is the Ext stratification
used in the proof of the cited result. Replacing $T$ by the finite disjoint
union of these strata, let $V \to T$ be the corresponding vector bundle.
Its tautological section is a relative extension class and hence gives a
short exact sequence
$$
0 \longrightarrow \mathcal{G}_V \longrightarrow \mathcal{H}
\longrightarrow \mathcal{G}'_V \longrightarrow 0
$$
on $X_V$. The base-change assertion says that every extension of a geometric
fibre of $\mathcal{G}'$ by the corresponding fibre of $\mathcal{G}$ occurs
as a fibre of this sequence after a field extension. Since $V$ is of finite
type over $S$ and $\mathcal{H}$ is coherent, it bounds the required middle
terms. Interchanging $\mathcal{G}$ and $\mathcal{G}'$ proves the final
assertion.
\end{proof}





\begin{lemma}
\label{lemma-bounded-fibrewise-reductions}
Let $T$ be a scheme of finite type over $S$, and let
$Z \subset X_T$ be a closed subscheme. The structure sheaves of the reduced
schemes
$$
(Z_{\overline{t}})_{red} \subset X_{\overline{t}},
$$
where $\overline{t}$ ranges over geometric points of $T$, form a bounded
family of fibrewise coherent sheaves over $X/S$.
\end{lemma}

\begin{proof}
We prove the assertion by Noetherian induction on $T$. If $T$ is reducible,
each of its finitely many irreducible components is a proper closed
subscheme and the induction hypothesis together with Lemma
\ref{lemma-bounded-coherent-elementary} applies. Thus we may assume that
$T$ is irreducible.

Apply More on Morphisms, Lemma
\ref{more-morphisms-lemma-make-generic-fibre-geometrically-reduced} to
$Z \to T$. After replacing $T$ by a nonempty open $V$, it gives a finite
universal homeomorphism $T' \to V$ and
$$
Z'=(Z \times_T T')_{red}.
$$
The scheme $T'$ is integral and the generic fibre of $Z' \to T'$ is
geometrically reduced. By More on Morphisms, Lemma
\ref{more-morphisms-lemma-geometrically-reduced-generic-fibre}, after
shrinking $T'$ and $V$ compatibly, every fibre of $Z' \to T'$ is
geometrically reduced.

The closed immersion $Z' \to Z \times_T T'$ is a universal homeomorphism.
The same is true on every geometric fibre. Since the source fibre is
reduced, it is the unique reduced closed subscheme with the underlying
topological space of the target fibre. Hence
$$
Z'_{\overline{t}'} = (Z_{\overline{t}'})_{red}
$$
for every geometric point $\overline{t}'$ of $T'$. The pushforward of
$\mathcal{O}_{Z'}$ to $X_{T'}$ is coherent, and $T'$ is of finite type over
$S$, so this one family bounds all reduced fibres over $V$. The complement
$T \setminus V$, with its reduced induced scheme structure, is a proper
closed subscheme of $T$ and is handled by the induction hypothesis. A finite
union finishes the proof.
\end{proof}

\begin{proposition}
\label{proposition-bounded-coherent-reduced-supports}
\begin{reference}
\cite[Expos\'e 221, Proposition 1.3, p. 253]{FGA}.
\end{reference}
Let $E$ be a bounded family of fibrewise coherent sheaves over $X/S$.
The structure sheaves
$$
\mathcal{O}_{(\operatorname{Supp}(\mathcal{F}))_{red}}
$$
for members $\mathcal{F}$ of $E$ over algebraically closed fields form a
bounded family.
\end{proposition}

\begin{proof}
Choose a pair $(T, \mathcal{G})$ bounding $E$. Let
$Z \subset X_T$ be the closed subscheme cut out by
$\operatorname{Fit}_0(\mathcal{G})$. Formation of this Fitting ideal
commutes with arbitrary base change by Divisors, Lemma
\ref{divisors-lemma-base-change-fitting-ideal}. Moreover, Divisors, Lemma
\ref{divisors-lemma-on-subscheme-cut-out-by-Fit-0} gives
$$
|Z_{\overline{t}}|=\operatorname{Supp}(\mathcal{G}_{\overline{t}})
$$
for every geometric point $\overline{t}$ of $T$. Consequently
$$
(Z_{\overline{t}})_{red}=
(\operatorname{Supp}(\mathcal{G}_{\overline{t}}))_{red}.
$$
Lemma \ref{lemma-bounded-fibrewise-reductions} bounds the structure sheaves
of the schemes on the left. If a member of $E$ is represented by
$\mathcal{F}$ over an algebraically closed field, then after passage to a
common extension it is isomorphic to a geometric fibre of $\mathcal{G}$.
The same is therefore true of the reduced supports. This proves the
proposition.
\end{proof}


\begin{lemma}
\label{lemma-annihilator-stratification}
Let $T$ be a scheme of finite type over $S$, and let $\mathcal{H}$ be a
coherent $\mathcal{O}_{X_T}$-module. There is a finite stratification
$$
T = \coprod T_a
$$
by locally closed subschemes such that, for every $a$ and every morphism
$T' \to T_a$, one has
$$
\mathop{\rm Ann}(\mathcal{H}_{T'}) =
\mathop{\rm Ann}(\mathcal{H}_{T_a})\mathcal{O}_{X_{T'}}.
$$
In particular, the structure sheaves
$$
\mathcal{O}_{X_{\overline{t}}}/
\mathop{\rm Ann}(\mathcal{H}_{\overline{t}})
$$
of the scheme theoretic supports of the geometric fibres form a bounded
family.
\end{lemma}

\begin{proof}
Choose a finite affine open covering $X_T = \bigcup U_j$. On $U_j$ choose
generators $h_{j1}, \ldots, h_{jr_j}$ of $\mathcal{H}|_{U_j}$ and consider
the map
$$
\mathcal{O}_{U_j} \longrightarrow
(\mathcal{H}|_{U_j})^{\oplus r_j},\qquad
a \longmapsto (ah_{j1},\ldots,ah_{jr_j}).
$$
Its kernel is $\mathop{\rm Ann}(\mathcal{H}|_{U_j})$. Apply More on
Morphisms, Lemma
\ref{more-morphisms-lemma-generic-flatness-stratification}, first to
$\mathcal{H}|_{U_j}$ and then to the cokernel of the displayed map, and
take a common finite refinement of the resulting stratifications of $T$.

On a stratum both the target and the cokernel of the displayed map are
flat over the stratum. Its image is therefore flat as well. The two short
exact sequences obtained by factoring the map through its image remain
exact after arbitrary base change. Hence its kernel commutes with arbitrary
base change. The chosen sections continue to generate after base change,
so the base-changed kernel is the annihilator of the base-changed module.
This proves the asserted equality on every $U_j$, and hence on $X_{T'}$.
The quotients by these annihilator ideals over the finitely many strata
give the bounded family in the last assertion.
\end{proof}


\begin{proposition}
\label{proposition-bounded-coherent-primary-pieces}
\begin{reference}
\cite[Expos\'e 221, the paragraph following Proposition 1.3, p. 253]{FGA};
\cite[IV, 3.2.6 and 9.8.2--9.8.4]{EGA}.
\end{reference}
Let $E$ be a bounded family of fibrewise coherent sheaves over $X/S$.
Let $(s,K,\mathcal{F})$ represent a member of $E$, with $K$ algebraically
closed. The following families are bounded.
\begin{enumerate}
\item For every generic point $x$ of an irreducible component of
$\operatorname{Supp}(\mathcal{F})$, the canonical primary quotient
$$
\mathcal{P}_x(\mathcal{F}) =
\mathop{\rm Im}\bigl(
\mathcal{F} \longrightarrow (j_x)_*j_x^*\mathcal{F}
\bigr),
$$
where $j_x : \Spec(\mathcal{O}_{X_K,x}) \to X_K$ is the canonical
morphism.
\item For every $x \in \operatorname{Ass}(\mathcal{F})$, the structure
sheaf $\mathcal{O}_{(\overline{\{x\}})_{red}}$.
\item For every $x$ as in (1), the structure sheaf
$$
\mathcal{O}_{X_K}/\mathfrak q_x,
\qquad
\mathfrak q_x =
\mathop{\rm Ann}(\mathcal{P}_x(\mathcal{F})).
$$
Here $\mathfrak q_x$ is primary with radical the prime ideal defining
$(\overline{\{x\}})_{red}$.
\end{enumerate}
\end{proposition}

\begin{remark}
\label{remark-EGA-primary-localization-typo}
In the last sentence of \cite[IV, 3.2.6]{EGA},
$\Spec(\kappa(x))$ is a typographical error for
$\Spec(\mathcal{O}_{X,x})$. The localization appearing in (1) is also
the construction used in the proof there. For example, if
$X = \Spec(k[\epsilon]/(\epsilon^2))$ and
$\mathcal{F}=\mathcal{O}_X$, then $x$ is the unique associated point and
the unique primary quotient is $\mathcal{F}$ itself, whereas the map to
the residue-field fibre has nonzero kernel.
\end{remark}

\begin{proof}
Choose a pair $(T,\mathcal{G})$ bounding $E$. It is enough to bound the
three constructions for every geometric fibre of $\mathcal{G}$. We do
this by Noetherian induction on $T$. Passing to the reduction and then to
the finitely many irreducible components reduces the induction step to the
case where $T$ is integral.

We recall the generic-fibre primary-decomposition argument. After a finite
dominant base change $T' \to V$, where $V$ is a nonempty open of $T$, a
reduced irredundant primary decomposition of the generic fibre can be
spread out to finitely many coherent quotients
$$
\mathcal{G}_{T'} \longrightarrow \mathcal{Q}_i.
$$
Let $Z_i$ be the closure of the associated point of the generic fibre of
$\mathcal{Q}_i$. The base change and the quotient list may be chosen so
that the generic fibres of all $Z_i \to T'$ are geometrically irreducible.
Indeed, this is the finite field-extension and spreading-out construction
of More on Morphisms, Lemma
\ref{more-morphisms-lemma-make-components-generic-fibre-geometrically-irreducible},
applied to the finitely many closures, followed by a refinement of the
primary decomposition. After shrinking once more, More on Morphisms,
Lemma \ref{more-morphisms-lemma-geom-irreducible-generic-fibre} makes every
geometric fibre of every $Z_i$ irreducible.

The primary-decomposition theorem near a generic fibre
\cite[IV, 9.8.2--9.8.4]{EGA} now gives, after a further shrinking, the
following facts for every geometric point $\overline{t}'$ of $T'$:
the map
$$
\mathcal{G}_{\overline{t}'} \longrightarrow
\bigoplus_i \mathcal{Q}_{i,\overline{t}'}
$$
is injective; each $\mathcal{Q}_{i,\overline{t}'}$ has no embedded
associated point; its unique associated point is the generic point of
$Z_{i,\overline{t}'}$; these points are pairwise distinct and exhaust
$\operatorname{Ass}(\mathcal{G}_{\overline{t}'})$; and the associated
point is maximal in the support precisely for the indices for which it was
maximal on the generic fibre.

For a maximal associated point the corresponding primary quotient is
unique; equivalently it is the image of localization displayed in (1),
see \cite[IV, 3.2.6]{EGA}. Thus the finitely many $\mathcal{Q}_i$
whose generic associated point is maximal bound the family in (1) over
$V$. For all indices $i$, Lemma
\ref{lemma-bounded-fibrewise-reductions}, applied to $Z_i \to T'$, bounds
the reduced closures of all associated points and hence the family in (2).
Finally apply Lemma \ref{lemma-annihilator-stratification} to the maximal
$\mathcal{Q}_i$. On its strata the quotient by the relative annihilator
has geometric fibres
$$
\mathcal{O}_{X_{\overline{t}'}}/
\mathop{\rm Ann}(\mathcal{Q}_{i,\overline{t}'}),
$$
which are exactly the sheaves in (3).

The cover $T' \to V$ is surjective, so every geometric fibre over $V$
occurs after a field extension among the fibres just constructed. The
closed complement $T \setminus V$, with its reduced induced structure, is
a proper closed subscheme and is handled by the induction hypothesis.
Taking a finite union completes the proof.
\end{proof}





\section{Boundedness for Quot}
\label{section-quot-bounded}

\noindent
Contrary to what happens classically, we already know the Quot
functor is an algebraic space, but we don't know that it is
ever represented by a finite type algebraic space.

\begin{lemma}
\label{lemma-quot-Pn}
Let $n \geq 0$, $r \geq 1$, $P \in \mathbf{Q}[t]$.
The algebraic space
$$
X = \Quotfunctor^P_{\mathcal{O}^{\oplus r}_{\mathbf{P}^n_\mathbf{Z}}/
\mathbf{P}^n_\mathbf{Z}/\mathbf{Z}}
$$
parametrizing quotients of $\mathcal{O}_{\mathbf{P}^n_\mathbf{Z}}^{\oplus r}$
with Hilbert polynomial $P$ is proper over $\Spec(\mathbf{Z})$.
\end{lemma}

\begin{proof}
We already know that $X \to \Spec(\mathbf{Z})$ is separated and
locally of finite presentation (Lemma \ref{lemma-quot-s-lfp}).
We also know that $X \to \Spec(\mathbf{Z})$ satisfies the
existence part of the valuative criterion, see
Lemma \ref{lemma-quot-existence-part}.
By the valuative criterion for properness, it suffices to
prove our Quot space is quasi-compact, see
Morphisms of Spaces,
Lemma \ref{spaces-morphisms-lemma-characterize-proper}.
Thus it suffices to find a quasi-compact scheme $T$ and a surjective
morphism $T \to X$. Let $m$ be the integer found in
Varieties, Lemma \ref{varieties-lemma-bound-quotients-free}.
Let
$$
N = r{m + n \choose n} - P(m)
$$
We will write $\mathbf{P}^n$ for
$\mathbf{P}^n_\mathbf{Z} = \text{Proj}(\mathbf{Z}[T_0, \ldots, T_n])$
and unadorned products will mean products over $\Spec(\mathbf{Z})$.
The idea of the proof is to construct a ``universal'' map
$$
\Psi :
\mathcal{O}_{T \times \mathbf{P}^n}(-m)^{\oplus N}
\longrightarrow
\mathcal{O}_{T \times \mathbf{P}^n}^{\oplus r}
$$
over an affine scheme $T$ and show that every point of $X$
corresponds to a cokernel of this in some point of $T$.

\medskip\noindent
Definition of $T$ and $\Psi$. We take $T = \Spec(A)$ where
$$
A = \mathbf{Z}[a_{i, j, E}]
$$
where $i \in \{1, \ldots, r\}$, $j \in \{1, \ldots, N\}$
and $E = (e_0, \ldots, e_n)$ runs through the multi-indices
of total degree $|E| = \sum_{k = 0, \ldots n} e_k = m$.
Then we define $\Psi$ to be the map whose $(i, j)$ matrix
entry is the map
$$
\sum\nolimits_{E = (e_0, \ldots, e_n)}
a_{i, j, E} T_0^{e_0} \ldots T_n^{e_n} :
\mathcal{O}_{T \times \mathbf{P}^n}(-m)
\longrightarrow
\mathcal{O}_{T \times \mathbf{P}^n}
$$
where the sum is over $E$ as above (but $i$ and $j$ are fixed of course).

\medskip\noindent
Consider the quotient $\mathcal{Q} = \Coker(\Psi)$ on $T \times \mathbf{P}^n$.
By More on Morphisms, Lemma
\ref{more-morphisms-lemma-generic-flatness-stratification}
there exists a $t \geq 0$ and closed subschemes
$$
T = T_0 \supset T_1 \supset \ldots \supset T_t = \emptyset
$$
such that the pullback $\mathcal{Q}_p$ of $\mathcal{Q}$ to
$(T_p \setminus T_{p + 1}) \times \mathbf{P}^n$ is flat over
$T_p \setminus T_{p + 1}$. Observe that we
have an exact sequence
$$
\mathcal{O}_{(T_p \setminus T_{p + 1}) \times \mathbf{P}^n}(-m)^{\oplus N}
\to
\mathcal{O}_{(T_p \setminus T_{p + 1}) \times \mathbf{P}^n}^{\oplus r}
\to
\mathcal{Q}_p
\to
0
$$
by pulling back the exact sequence defining $\mathcal{Q} = \Coker(\Psi)$.
Therefore we obtain a morphism
$$
\coprod (T_p \setminus T_{p + 1})
\longrightarrow
\Quotfunctor_{\mathcal{O}^{\oplus r}/\mathbf{P}/\mathbf{Z}}
\supset
\Quotfunctor^P_{\mathcal{O}^{\oplus r}/\mathbf{P}/\mathbf{Z}} = X
$$
Since the left hand side is a Noetherian scheme and the inclusion
on the right hand side is open, it
suffices to show that any point of $X$ is in the image of this morphism.

\medskip\noindent
Let $k$ be a field and let $x \in X(k)$. Then $x$ corresponds to
a surjection $\mathcal{O}_{\mathbf{P}^n_k}^{\oplus r} \to \mathcal{F}$
of coherent $\mathcal{O}_{\mathbf{P}^n_k}$-modules
such that the Hilbert polynomial of $\mathcal{F}$ is $P$.
Consider the short exact sequence
$$
0 \to \mathcal{K} \to
\mathcal{O}_{\mathbf{P}^n_k}^{\oplus r} \to
\mathcal{F} \to 0
$$
By Varieties, Lemma \ref{varieties-lemma-bound-quotients-free}
and our choice of $m$ we see that $\mathcal{K}$ is $m$-regular.
By Varieties, Lemma \ref{varieties-lemma-m-regular-globally-generated}
we see that $\mathcal{K}(m)$ is globally generated.
By Varieties, Lemma \ref{varieties-lemma-m-regular-up}
and the definition of $m$-regularity we see that
$H^i(\mathbf{P}^n_k, \mathcal{K}(m)) = 0$ for $i > 0$.
Hence we see that
$$
\dim_k H^0(\mathbf{P}^n_k, \mathcal{K}(m)) =
\chi(\mathcal{K}(m)) =
\chi(\mathcal{O}_{\mathbf{P}^n_k}(m)^{\oplus r}) -
\chi(\mathcal{F}(m)) = N
$$
by our choice of $N$. This gives a surjection
$$
\mathcal{O}_{\mathbf{P}^n_k}^{\oplus N}
\longrightarrow
\mathcal{K}(m)
$$
Twisting back down and using the short exact sequence above
we see that $\mathcal{F}$ is the cokernel of a map
$$
\Psi_x :
\mathcal{O}_{\mathbf{P}^n_k}(-m)^{\oplus N}
\to
\mathcal{O}_{\mathbf{P}^n_k}^{\oplus r}
$$
There is a unique ring map $\tau : A \to k$ such that the base change
of $\Psi$ by the corresponding morphism $t = \Spec(\tau) : \Spec(k) \to T$
is $\Psi_x$. This is true because the entries of the $N \times r$
matrix defining $\Psi_x$ are homogeneous polynomials
$\sum \lambda_{i, j, E} T_0^{e_0} \ldots T_n^{e_n}$
of degree $m$ in $T_0, \ldots, T_n$ with coefficients
$\lambda_{i, j, E} \in k$ and we can set
$\tau(a_{i, j, E}) = \lambda_{i, j, E}$.
Then $t \in T_p \setminus T_{p + 1}$ for some $p$ and
the image of $t$ under the morphism above is $x$ as desired.
\end{proof}

\begin{lemma}
\label{lemma-quot-Pn-projective}
\begin{reference}
\cite[Expos\'e 221, Lemma 3.3, Theorem 3.2, and Proposition 3.8,
pp. 260--266]{FGA}.
\end{reference}
Let $n \geq 0$, $r \geq 1$, $P \in \mathbf{Q}[t]$. The algebraic space
$$
X = \Quotfunctor^P_{\mathcal{O}^{\oplus r}_{\mathbf{P}^n_\mathbf{Z}}/
\mathbf{P}^n_\mathbf{Z}/\mathbf{Z}}
$$
is a projective scheme over $\mathbf{Z}$. If $X$ is nonempty, then for
all sufficiently large $m$ the universal quotient $\mathcal{G}$ has the
following properties. The module
$$
\mathcal{E}_m = p_*\mathcal{G}(m),\qquad
p : \mathbf{P}^n_X \longrightarrow X,
$$
is finite locally free of rank $P(m)$ and its formation commutes with
arbitrary base change. The quotient
$$
H^0(\mathbf{P}^n_\mathbf{Z},
\mathcal{O}(m)^{\oplus r}) \otimes_\mathbf{Z} \mathcal{O}_X
\longrightarrow \mathcal{E}_m
$$
defines a closed immersion into the scheme of locally free quotients
of the indicated rank, and the pullback of its Pl\"ucker invertible
module is $\det(\mathcal{E}_m)$.
\end{lemma}

\begin{proof}
By Lemma \ref{lemma-quot-Pn}, $X$ is proper over $\mathbf{Z}$; in
particular it is quasi-compact. If $X$ is empty, the assertion is clear,
so assume it is nonempty. On $\mathbf{P}^n_X$ write the universal
sequence as
$$
0 \longrightarrow \mathcal{K} \longrightarrow
\mathcal{O}^{\oplus r} \longrightarrow \mathcal{G}
\longrightarrow 0.
$$
The modules $\mathcal{K}$ and $\mathcal{G}$ are of finite presentation
and flat over $X$. Choose $m_0 \geq 0$ as in Varieties, Lemma
\ref{varieties-lemma-bound-quotients-free}. For every geometric point
$x$ of $X$, the module $\mathcal{K}_x$ is $m_0$-regular. Hence for
$m \geq m_0$ the higher cohomology of $\mathcal{K}_x(m)$ and
$\mathcal{G}_x(m)$ vanishes, and $\mathcal{K}_x(m)$ is globally
generated; see Varieties, Lemmas
\ref{varieties-lemma-m-regular-up} and
\ref{varieties-lemma-m-regular-globally-generated}.

By Derived Categories of Spaces, Lemmas
\ref{spaces-perfect-lemma-flat-proper-perfect-direct-image-general} and
\ref{spaces-perfect-lemma-open-where-cohomology-in-degree-i-rank-r-geometric},
the modules
$$
\mathcal{W}_m = p_*\mathcal{K}(m)
\quad\text{and}\quad
\mathcal{E}_m = p_*\mathcal{G}(m)
$$
are finite locally free, of ranks
$$
r{n + m \choose n} - P(m)\quad\text{and}\quad P(m),
$$
respectively, their formation commutes with arbitrary base change, and
there is an exact sequence
$$
0 \longrightarrow \mathcal{W}_m \longrightarrow
H^0(\mathbf{P}^n_\mathbf{Z},
\mathcal{O}(m)^{\oplus r}) \otimes_\mathbf{Z} \mathcal{O}_X
\longrightarrow \mathcal{E}_m \longrightarrow 0.
$$
Indeed, the cited open loci contain every geometric point of $X$ and
therefore equal $X$.

Put $M = r{n + m \choose n}$ and $e = P(m)$. The last quotient defines
a morphism
$$
\varphi_m : X \longrightarrow H,
$$
where $H = \mathbf{G}(M - e, M)$ if $0 < e < M$, and where $H$ is
$\Spec(\mathbf{Z})$ in the two extreme cases. The evaluation map
$p^*\mathcal{W}_m \to \mathcal{K}(m)$ is surjective: this can be checked
on geometric fibres, where it follows from regularity. Thus after any
base change the quotient of global sections determines
$\mathcal{K}(m)$, and hence determines the universal quotient
$\mathcal{G}$. It follows that $\varphi_m$ is a monomorphism.

The morphism $\varphi_m$ is locally of finite type. It is also proper:
its graph is a closed immersion because $H$ is separated over
$\mathbf{Z}$, and the projection $X \times_\mathbf{Z} H \to H$ is
proper. Hence $\varphi_m$ is a closed immersion by More on Morphisms of
Spaces, Lemma
\ref{spaces-more-morphisms-lemma-characterize-closed-immersion}.
If $0 < e < M$, Constructions, Lemma
\ref{constructions-lemma-grassmannian-pluecker} gives a closed immersion
of $H$ into projective space and identifies the pullback of
$\mathcal{O}(1)$ with the determinant of its universal quotient. The
same conclusions are immediate in the extreme cases. Thus $X$ is a
projective scheme and the final assertion follows.
\end{proof}

\begin{lemma}
\label{lemma-quot-Pn-over-base}
Let $B$ be an algebraic space. Let $X = B \times \mathbf{P}^n_\mathbf{Z}$.
Let $\mathcal{L}$ be the pullback of $\mathcal{O}_{\mathbf{P}^n}(1)$ to $X$.
Let $\mathcal{F}$ be an $\mathcal{O}_X$-module of finite
presentation. The algebraic space $\Quotfunctor^P_{\mathcal{F}/X/B}$
parametrizing quotients of $\mathcal{F}$
having Hilbert polynomial $P$ with respect to $\mathcal{L}$
is proper over $B$.
\end{lemma}

\begin{proof}
The question is \'etale local over $B$, see
Morphisms of Spaces, Lemma \ref{spaces-morphisms-lemma-proper-local}.
Thus we may assume $B$ is an affine scheme.
In this case $\mathcal{L}$ is an ample invertible module on $X$
(by Constructions, Lemma \ref{constructions-lemma-ample-on-proj}
and the definition of ample invertible modules in
Properties, Definition \ref{properties-definition-ample}).
Thus we can find $r' \geq 0$ and $r \geq 0$ and a surjection
$$
\mathcal{O}_X^{\oplus r} \longrightarrow
\mathcal{F} \otimes_{\mathcal{O}_X} \mathcal{L}^{\otimes r'}
$$
by Properties, Proposition \ref{properties-proposition-characterize-ample}.
By Lemma \ref{lemma-quot-tensor-invertible}
we may replace $\mathcal{F}$ by
$\mathcal{F} \otimes_{\mathcal{O}_X} \mathcal{L}^{\otimes r'}$
and $P(t)$ by $P(t + r')$.
By Lemma \ref{lemma-quot-quotient}
we obtain a closed immersion
$$
\Quotfunctor^P_{\mathcal{F}/X/B}
\longrightarrow
\Quotfunctor^P_{\mathcal{O}_X^{\oplus r}/X/B}
$$
Since we've shown that $\Quotfunctor^P_{\mathcal{O}_X^{\oplus r}/X/B} \to B$
is proper in Lemma \ref{lemma-quot-Pn} we conclude.
\end{proof}

\begin{lemma}
\label{lemma-quot-proper-over-base}
Let $f : X \to B$ be a proper morphism of finite presentation
of algebraic spaces. Let $\mathcal{F}$ be a finitely presented
$\mathcal{O}_X$-module. Let $\mathcal{L}$ be an invertible
$\mathcal{O}_X$-module ample on $X/B$, see
Divisors on Spaces, Definition
\ref{spaces-divisors-definition-relatively-ample}.
The algebraic space $\Quotfunctor^P_{\mathcal{F}/X/B}$
parametrizing quotients of $\mathcal{F}$
having Hilbert polynomial $P$ with respect to $\mathcal{L}$
is proper over $B$.
\end{lemma}

\begin{proof}
The question is \'etale local over $B$, see
Morphisms of Spaces, Lemma \ref{spaces-morphisms-lemma-proper-local}.
Thus we may assume $B$ is an affine scheme.
Then we can find a closed immersion $i : X \to \mathbf{P}^n_B$
such that $i^*\mathcal{O}_{\mathbf{P}^n_B}(1) \cong \mathcal{L}^{\otimes d}$
for some $d \geq 1$. See
Morphisms, Lemma \ref{morphisms-lemma-quasi-projective-finite-type-over-S}.
Changing $\mathcal{L}$ into $\mathcal{L}^{\otimes d}$ and
the numerical polynomial $P(t)$ into $P(dt)$ leaves
$\Quotfunctor^P_{\mathcal{F}/X/B}$ unaffected; some details omitted.
Hence we may assume $\mathcal{L} = i^*\mathcal{O}_{\mathbf{P}^n_B}(1)$.
Then the isomorphism
$\Quotfunctor_{\mathcal{F}/X/B} \to
\Quotfunctor_{i_*\mathcal{F}/\mathbf{P}^n_B/B}$ of
Lemma \ref{lemma-quot-closed} induces an isomorphism
$\Quotfunctor^P_{\mathcal{F}/X/B} \cong
\Quotfunctor^P_{i_*\mathcal{F}/\mathbf{P}^n_B/B}$.
Since $\Quotfunctor^P_{i_*\mathcal{F}/\mathbf{P}^n_B/B}$
is proper over $B$ by Lemma \ref{lemma-quot-Pn-over-base}
we conclude.
\end{proof}

\begin{proposition}
\label{proposition-quot-projective-over-base}
\begin{reference}
\cite[Expos\'e 221, Theorem 3.2 and Proposition 3.8,
pp. 260--266]{FGA}.
\end{reference}
Let $S$ be a Noetherian scheme, let $f : X \to S$ be a projective
morphism of schemes, let $\mathcal{L}$ be an $f$-very ample invertible
module, and let $\mathcal{F}$ be a coherent $\mathcal{O}_X$-module.
For $P \in \mathbf{Q}[t]$ the algebraic space
$$
Q = \Quotfunctor^P_{\mathcal{F}/X/S}
$$
is a projective scheme over $S$.

Let $\mathcal{G}$ be the universal quotient on $X_Q$, and write
$p : X_Q \to Q$. If $Q$ is nonempty, then for all sufficiently large
$m$ the module
$$
\mathcal{E}_m = p_*(\mathcal{G} \otimes \mathcal{L}^{\otimes m})
$$
is finite locally free of rank $P(m)$ and its formation commutes with
arbitrary base change. Zariski locally on $S$, a quotient of a finite
free module onto $\mathcal{E}_m$ defines a closed immersion of $Q$ into
the corresponding Grassmannian. In particular,
$\det(\mathcal{E}_m)$ is very ample on $Q/S$.
\end{proposition}

\begin{proof}
The morphism $Q \to S$ is proper by
Lemma \ref{lemma-quot-proper-over-base}. Since $S$ is Noetherian, it has
a finite affine open covering $S = \bigcup U_i$. After refining this
covering, Morphisms, Lemma
\ref{morphisms-lemma-very-ample-finite-type-over-affine} and properness
give closed immersions
$$
\iota_i : X_{U_i} \longrightarrow \mathbf{P}^{n_i}_{U_i}
$$
such that $\mathcal{L}|_{X_{U_i}} = \iota_i^*\mathcal{O}(1)$. After
choosing $a_i \geq 0$ sufficiently large, there is a surjection
$$
\mathcal{O}_{\mathbf{P}^{n_i}_{U_i}}^{\oplus r_i}
\longrightarrow (\iota_i)_*(\mathcal{F}|_{X_{U_i}})(a_i).
$$
Here we use Properties, Proposition
\ref{properties-proposition-characterize-ample}.

Put $P_i(t) = P(t + a_i)$. Lemmas
\ref{lemma-quot-tensor-invertible}, \ref{lemma-quot-closed}, and
\ref{lemma-quot-quotient} identify $Q_{U_i}$ with a closed subspace of
$$
\Quotfunctor^{P_i}_{
\mathcal{O}^{\oplus r_i}_{\mathbf{P}^{n_i}_{U_i}}/
\mathbf{P}^{n_i}_{U_i}/U_i}.
$$
The latter is the base change to $U_i$ of the projective scheme over
$\mathbf{Z}$ constructed in Lemma \ref{lemma-quot-Pn-projective}.
Consequently $Q_{U_i}$ is a projective scheme over $U_i$. Since being a
scheme is local on the base, $Q$ is a scheme.

If $Q$ is empty, this also proves that it is projective and there is
nothing further to show. Assume $Q$ is nonempty.

The final part of Lemma \ref{lemma-quot-Pn-projective}, restricted to
the displayed closed subspace, shows that for all sufficiently large
$d$ the module
$$
p_*\mathcal{G}(a_i + d)
$$
is finite locally free of rank $P(a_i + d)$, commutes with arbitrary
base change, and is the universal quotient defining a closed immersion
into a Grassmannian. As the covering is finite, one integer works for
all $i$ after writing $m = a_i + d$. These local modules are the
restrictions of $\mathcal{E}_m$, so they glue and prove the assertions
about $\mathcal{E}_m$ and base change.

By Constructions, Lemma
\ref{constructions-lemma-grassmannian-pluecker}, the restriction of
$\det(\mathcal{E}_m)$ to every $Q_{U_i}$ is very ample. Morphisms,
Lemma \ref{morphisms-lemma-relatively-very-ample} shows that
$\det(\mathcal{E}_m)$ is very ample on $Q/S$. It is therefore ample by
Morphisms, Lemma \ref{morphisms-lemma-ample-very-ample}. Thus $Q \to S$
is quasi-projective as well as proper, and hence projective by
Morphisms, Lemma
\ref{morphisms-lemma-projective-is-quasi-projective-proper}.
\end{proof}

\begin{lemma}
\label{corollary-quot-projective-components}
\begin{reference}
\cite[Expos\'e 221, Theorem 3.1, pp. 259--260]{FGA}.
\end{reference}
Under the hypotheses of Proposition
\ref{proposition-quot-projective-over-base}, there is a decomposition
$$
\Quotfunctor_{\mathcal{F}/X/S} =
\coprod_{P \in \mathbf{Q}[t]}
\Quotfunctor^P_{\mathcal{F}/X/S}
$$
as a disjoint union of projective $S$-schemes. In particular, the Quot
space is a scheme locally of finite type over $S$.
\end{lemma}

\begin{proof}
The Hilbert polynomial of a flat family is locally constant on the
base. Thus the open and closed subspaces of
Remark \ref{remark-quot-numerical} give the displayed disjoint union.
Each term is projective by Proposition
\ref{proposition-quot-projective-over-base}.
\end{proof}

\begin{proposition}
\label{proposition-quot-quasi-projective-over-base}
\begin{reference}
\cite[Expos\'e 221, Section 4(a), pp. 266--267]{FGA}.
\end{reference}
Let $S$ be a Noetherian scheme, let $X \to S$ be quasi-projective,
and let $\mathcal{F}$ be a coherent $\mathcal{O}_X$-module. Choose a
factorization
$$
X \xrightarrow{j} \overline{X} \longrightarrow S
$$
where $j$ is an open immersion and $\overline{X} \to S$ is projective,
and choose a very ample invertible module $\overline{\mathcal{L}}$ on
$\overline{X}/S$. For every numerical polynomial $P$, the functor
$$
\Quotfunctor^P_{\mathcal{F}/X/S}
$$
of quotients flat over the variable base and with proper support is a
quasi-projective scheme over $S$. Consequently
$$
\Quotfunctor_{\mathcal{F}/X/S} =
\coprod_{P \in \mathbf{Q}[t]}
\Quotfunctor^P_{\mathcal{F}/X/S}
$$
is a scheme locally of finite type over $S$ whose fixed-polynomial
components are quasi-projective. Here the Hilbert polynomial is computed
with respect to $j^*\overline{\mathcal{L}}$.
\end{proposition}

\begin{proof}
The factorization exists by Morphisms, Lemma
\ref{morphisms-lemma-quasi-projective-open-projective}. Since $S$ is
Noetherian, both $X$ and $\overline{X}$ are Noetherian and $j$ is
quasi-compact. Properties, Lemma
\ref{properties-lemma-lift-finite-presentation} gives a finitely presented,
hence coherent, $\mathcal{O}_{\overline{X}}$-module
$\overline{\mathcal{F}}$ with
$j^*\overline{\mathcal{F}} \cong \mathcal{F}$.

Lemma \ref{lemma-quot-better-open} gives an open immersion
$$
\Quotfunctor^P_{\mathcal{F}/X/S} \longrightarrow
\Quotfunctor^P_{\overline{\mathcal{F}}/\overline{X}/S}.
$$
The target is a projective scheme by Proposition
\ref{proposition-quot-projective-over-base}. It is Noetherian, so the open
subscheme on the left is quasi-compact. The open immersion is therefore
quasi-projective by Morphisms, Lemma
\ref{morphisms-lemma-quasi-affine-finite-type-quasi-projective}, and the
composite with the projective morphism to $S$ is quasi-projective by
Morphisms, Lemma \ref{morphisms-lemma-composition-quasi-projective}.

Finally, the Hilbert polynomial is locally constant in a flat family.
Remark \ref{remark-quot-numerical} supplies the open and closed
decomposition, and the preceding argument applies to every component.
\end{proof}

\begin{proposition}
\label{proposition-bounded-coherent-Hilbert-polynomials}
\begin{reference}
\cite[Expos\'e 221, Theorem 2.1 and its erratum, pp. 253--254 and 302]{FGA}.
\end{reference}
Let $S$ be a Noetherian scheme, let $X \to S$ be projective, and fix a
relatively very ample invertible module $\mathcal{O}_X(1)$. A family $E$
of fibrewise coherent sheaves over $X/S$ is bounded if and only if the
following conditions hold.
\begin{enumerate}
\item There is a coherent $\mathcal{O}_X$-module $\mathcal{L}$ such that
every member $(s,K,\mathcal{F})$ of $E$, after a field extension if
necessary, is represented by a quotient
$$
\mathcal{L}_K \longrightarrow \mathcal{F} \longrightarrow 0.
$$
One may require $\mathcal{L}$ to be a finite direct sum of negative twists
of $\mathcal{O}_X$.
\item The Hilbert polynomials of the members of $E$, computed with respect
to $\mathcal{O}_X(1)$, belong to a fixed finite set.
\end{enumerate}
\end{proposition}

\begin{proof}
Suppose first that $(T,\mathcal{G})$ bounds $E$. The scheme $T$ is
quasi-compact. On each member of a finite affine open covering of $T$,
relative Serre vanishing and relative global generation give a surjection
$$
\mathcal{O}_{X_T}(-n)^{\oplus N} \longrightarrow \mathcal{G}
\longrightarrow 0
$$
for suitable $n$ and $N$, which may depend on the member of the covering.
The direct sum of the finitely many corresponding modules on $X$ gives a
single $\mathcal{L}$ satisfying (1).

By Lemma \ref{lemma-bounded-coherent-elementary}, we may replace $T$ by a
finite stratification on which $\mathcal{G}$ is flat. On each stratum its
Hilbert polynomial is locally constant. Each stratum is Noetherian and has
only finitely many connected components. Hence only finitely many Hilbert
polynomials occur, proving (2).

Conversely, let $P_1,\ldots,P_m$ be the finite list in (2). For each $i$,
Lemma \ref{lemma-quot-proper-over-base} shows that
$$
\Quotfunctor^{P_i}_{\mathcal{L}/X/S}
$$
is proper, and in particular of finite type, over $S$. The finite disjoint
union of these Quot spaces is quasi-compact. Choose a quasi-compact scheme
mapping surjectively and \'etale to it. This scheme is of finite type over
$S$, and the pullback of the universal quotient to it bounds all quotients
allowed by (1) whose Hilbert polynomial is in the list. It therefore bounds
$E$.
\end{proof}


\begin{lemma}
\label{lemma-coherent-dimension-truncation}
\begin{reference}
\cite[Expos\'e 221, the paragraph preceding Theorem 2.2,
pp. 254--255]{FGA}.
\end{reference}
Let $k$ be a field, let $X$ be a scheme of finite type over $k$, and let
$\mathcal{F}$ be a coherent $\mathcal{O}_X$-module. For every integer $r$
there is a unique largest coherent submodule
$$
\mathcal{N}_r(\mathcal{F}) \subset \mathcal{F}
$$
whose support has dimension $< r$. It is the subsheaf whose sections are the
sections of $\mathcal{F}$ having support of dimension $< r$. Put
$$
\mathcal{F}_{(r)} = \mathcal{F}/\mathcal{N}_r(\mathcal{F}).
$$
Then
\begin{enumerate}
\item the $\mathcal{N}_r(\mathcal{F})$ form a finite increasing filtration,
with $\mathcal{N}_r(\mathcal{F})=0$ for $r \leq 0$ and
$\mathcal{N}_r(\mathcal{F})=\mathcal{F}$ for
$r>\dim(\operatorname{Supp}(\mathcal{F}))$,
\item
$$
\operatorname{Ass}(\mathcal{F}_{(r)})=
\{x\in\operatorname{Ass}(\mathcal{F})\mid
\dim(\overline{\{x\}})\geq r\},
$$
\item
$$
\operatorname{Ass}(\mathcal{N}_{r+1}(\mathcal{F})/
\mathcal{N}_r(\mathcal{F}))=
\{x\in\operatorname{Ass}(\mathcal{F})\mid
\dim(\overline{\{x\}})=r\},
$$
and
\item formation of $\mathcal{N}_r(\mathcal{F})$ and
$\mathcal{F}_{(r)}$ commutes with extension of the ground field.
\end{enumerate}
\end{lemma}

\begin{proof}
The sum of two coherent submodules whose supports have dimension $<r$ has
the same property. Since $\mathcal{F}$ is a Noetherian object, the sum of all
such submodules is already a finite sum. This gives the largest submodule
$\mathcal{N}_r(\mathcal{F})$. The same construction after restriction to an
open subscheme gives its restriction. A local section belongs to this
submodule if and only if the coherent submodule it generates has support of
dimension $<r$. This proves the description by sections and (1).

Set $\mathcal{N}=\mathcal{N}_r(\mathcal{F})$ and
$\mathcal{Q}=\mathcal{F}/\mathcal{N}$. The maximality of $\mathcal{N}$ says
that $\mathcal{Q}$ has no nonzero coherent submodule with support of
dimension $<r$. Using the usual properties of associated primes in the exact
sequence
$$
0\longrightarrow\mathcal{N}\longrightarrow\mathcal{F}
\longrightarrow\mathcal{Q}\longrightarrow0
$$
we obtain
$$
\operatorname{Ass}(\mathcal{Q})=
\{x\in\operatorname{Ass}(\mathcal{F})\mid
\dim(\overline{\{x\}})\geq r\}.
$$
Indeed, an associated point of $\mathcal{Q}$ of smaller dimension would
generate a forbidden submodule. Conversely, if $x$ is associated to
$\mathcal{F}$ and has dimension at least $r$, then
$x\notin\operatorname{Supp}(\mathcal{N})$, and localization at $x$ shows that
$x$ is associated to $\mathcal{Q}$. This proves (2). The associated points of
$\mathcal{N}_{r+1}(\mathcal{F})$ are precisely the associated points of
$\mathcal{F}$ having dimension at most $r$. Applying (2) to this submodule
proves (3).

Let $K/k$ be a field extension. The support of
$\mathcal{N}_r(\mathcal{F})_K$ has dimension $<r$. By Algebra, Lemma
\ref{algebra-lemma-bourbaki}, the associated points of
$\mathcal{F}_{(r),K}$ lie over associated points of
$\mathcal{F}_{(r)}$; the corresponding irreducible components have the same
dimension after extension of the ground field. Hence
$\mathcal{F}_{(r),K}$ has no nonzero coherent submodule with support of
dimension $<r$. Thus $\mathcal{N}_r(\mathcal{F})_K$ is the largest such
submodule of $\mathcal{F}_K$, proving (4).
\end{proof}

\begin{remark}
\label{remark-FGA-Hilbert-coefficient-inequality}
The inequality in the printed hypothesis of
\cite[Expos\'e 221, Theorem 2.2, pp. 254--255]{FGA} is reversed. It asks that
the coefficients in degrees $\leq s-1$ be bounded. For a counterexample, let
$X=\mathbf{P}^2_k$ over an algebraically closed field of characteristic zero.
For $d\geq3$, choose a smooth plane curve $C_d$ of degree $d$ and a
zero-dimensional subscheme $Z_d$ of length $d(d-3)/2$. Then
$$
\mathcal{F}_d=\mathcal{O}_{C_d}\oplus\mathcal{O}_{Z_d}
$$
is a quotient of $\mathcal{O}_X^{\oplus2}$ and
$$
P_{\mathcal{F}_d}(n)=
dn+\frac{3d-d^2}{2}+\frac{d(d-3)}2=dn.
$$
Thus for $s=1$ all coefficients in degrees $\leq0$ are fixed, whereas
$(\mathcal{F}_d)_{(1)}=\mathcal{O}_{C_d}$ is not a bounded family: its leading
coefficient is $d$. The proof of the cited theorem and its Corollary 2.3 both
require coefficients in degrees $\geq s-1$. We use this corrected inequality
below.
\end{remark}


\begin{lemma}
\label{lemma-bounded-finite-pushforward}
Let $S$ be a Noetherian scheme. Let $X$ and $Y$ be schemes proper and of
finite type over $S$, and let $f:X\to Y$ be finite. A family $E$ of fibrewise
coherent sheaves on $X/S$ is bounded if and only if the family
$$
\{f_*\mathcal{F}\mid\mathcal{F}\text{ represents a member of }E\}
$$
is bounded on $Y/S$.
\end{lemma}

\begin{proof}
If $(T,\mathcal{G})$ bounds $E$, then
$(T,(f_T)_*\mathcal{G})$ bounds the pushforwards. Here we use that a finite
morphism is affine, that its pushforward preserves coherent modules, and that
this pushforward commutes with base change.

Conversely, choose $(T,\mathcal{H})$ bounding the pushforwards and put
$$
\mathcal{A}=(f_T)_*\mathcal{O}_{X_T}.
$$
After a finite stratification of $T$, we may assume that the coherent modules
which occur below are flat over $T$. The functor of homomorphisms
$$
\mathcal{A}\otimes\mathcal{H}\longrightarrow\mathcal{H}
$$
is then represented by an affine scheme of finite presentation over $T$, as
in the proof of Proposition
\ref{proposition-bounded-coherent-morphisms-extensions}. The unit and
associativity identities cut out a closed subscheme $M$ of this Hom scheme.
Thus $M$ parametrizes all $\mathcal{A}$-module structures on the fibres of
$\mathcal{H}$. The universal $\mathcal{A}$-module on $Y_M$ corresponds to a
coherent module on $X_M$ by Morphisms, Lemma
\ref{morphisms-lemma-affine-equivalence-modules}. Every module in $E$ occurs
as a fibre after extending the ground field and transporting its
$f_*\mathcal{O}_X$-module structure to the corresponding fibre of
$\mathcal{H}$. Since $M$ is of finite type over $S$, the universal module on
$X_M$ bounds $E$.
\end{proof}

\begin{lemma}
\label{lemma-bounded-reduced-equidimensional-degree}
\begin{reference}
\cite[Expos\'e 221, Lemma 2.4, p. 255]{FGA}.
\end{reference}
Let $S$ be a Noetherian scheme, let $X\to S$ be projective, and fix a
relatively very ample invertible module $\mathcal{O}_X(1)$. Fix integers
$r,d\geq0$. The structure sheaves of the reduced closed subschemes
$$
Z\subset X_K
$$
which are pure of dimension $r$ and have degree at most $d$ form a bounded
family over $X/S$.
\end{lemma}

\begin{proof}
Choose a closed immersion $X\to\mathbf{P}^N_S$ defined by
$\mathcal{O}_X(1)$. For each $e\leq d$, the Chow-coordinate construction gives
a projective parameter scheme for effective $r$-cycles of degree $e$ in
$\mathbf{P}^N$, together with its incidence subscheme. Restricting to the
closed locus of cycles supported on $X$ and taking the product with $S$ gives
a finite-type parameter scheme over $S$. A reduced pure-dimensional
$Z\subset X_K$ determines the cycle which is the sum of its irreducible
components with coefficient $1$, and the reduction of the corresponding
incidence fibre is exactly $Z$. Lemma
\ref{lemma-bounded-fibrewise-reductions}, applied to the finitely many
incidence families for $0\leq e\leq d$, bounds their reduced fibres and hence
all the stated structure sheaves.
\end{proof}

\begin{lemma}
\label{lemma-bounded-top-dimensional-small-ambient}
\begin{reference}
\cite[Expos\'e 221, Lemma 2.5, pp. 255--257]{FGA}.
\end{reference}
Let $S$ be a Noetherian scheme, let $X\to S$ be projective with fibres of
dimension at most $r$, and fix a relatively very ample invertible module
$\mathcal{O}_X(1)$. Let $\mathcal{L}$ be a coherent $\mathcal{O}_X$-module,
and let $E$ be a family of quotient modules of the modules $\mathcal{L}_K$.
If $r\geq1$, write
$$
P_{\mathcal{F}}(n)=
a_{\mathcal{F}}\frac{n^r}{r!}+
b_{\mathcal{F}}\frac{n^{r-1}}{(r-1)!}+
\text{terms of degree }<r-1.
$$
Then the $a_{\mathcal{F}}$ are bounded and the $b_{\mathcal{F}}$ are bounded
below. If the $b_{\mathcal{F}}$ are bounded above, then the family of
$\mathcal{F}_{(r)}$ is bounded. If $r=0$, then $E$ is bounded.
\end{lemma}

\begin{proof}
Embed $X$ into $\mathbf{P}^N_S$ using $\mathcal{O}_X(1)$. The linear
projections $\mathbf{P}^N\dashrightarrow\mathbf{P}^r$ are parametrized by a
scheme of finite type. The condition that a projection be defined on $X$ and
induce a finite morphism on $X$ is open, and the standard generic-projection
argument shows that this open meets every geometric fibre over $S$. Hence,
after replacing $S$ by a finite-type parameter scheme which meets every
geometric fibre, we may assume there is a finite morphism
$$
f:X\longrightarrow\mathbf{P}^r_S,
\qquad
\mathcal{O}_X(1)=f^*\mathcal{O}_{\mathbf{P}^r_S}(1).
$$
This replacement is harmless for boundedness in Definition
\ref{definition-bounded-coherent}, since every member is allowed a ground
field extension.

Finite pushforward preserves the Hilbert polynomial by the projection formula
and the vanishing of higher direct images. Since a finite morphism preserves
dimensions of supports, the description by sections in Lemma
\ref{lemma-coherent-dimension-truncation} gives
$$
f_*(\mathcal{F}_{(r)})=(f_*\mathcal{F})_{(r)}.
$$
Lemma \ref{lemma-bounded-finite-pushforward} therefore reduces the question to
$X=\mathbf{P}^r_S$. After taking a finite affine covering of $S$, replacing it
by the corresponding disjoint union, and twisting by a fixed integer, we may
also assume that every member is a quotient of
$\mathcal{O}_{\mathbf{P}^r_K}^{\oplus q}$ for one fixed $q$. Twisting changes
$b_{\mathcal{F}}$ by a fixed multiple of $a_{\mathcal{F}}$, so it does not
change any of the boundedness assertions.

The leading coefficient $a_{\mathcal{F}}$ is the generic rank of
$\mathcal{F}$ and hence
$$
0\leq a_{\mathcal{F}}\leq q.
$$
The quotient $\mathcal{G}=\mathcal{F}_{(r)}$ is either zero or a torsion-free
module of rank $a=a_{\mathcal{F}}$ on $\mathbf{P}^r_K$. Assume it is nonzero.
On the complement of a closed subset of codimension at least $2$ it is
locally free. Since projective space is regular and locally factorial, the
rank-one reflexive module
$$
\det(\mathcal{G})=(\wedge^a\mathcal{G})^{**}
$$
is invertible. Divisors, Lemma
\ref{divisors-lemma-Pic-projective-space-UFD} gives a unique integer $c$ such
that $\det(\mathcal{G})\cong\mathcal{O}(c)$. The quotient map from
$\mathcal{O}^{\oplus q}$ gives a nonzero homomorphism
$$
s:\wedge^a\mathcal{O}^{\oplus q}\longrightarrow\mathcal{O}(c).
$$
Consequently $c\geq0$.

Modulo coherent modules supported in codimension at least $2$, the class of a
torsion-free module on the regular scheme $\mathbf{P}^r_K$ is determined by
its rank and determinant. Thus the Hilbert polynomials of $\mathcal{G}$ and
$\mathcal{O}^{\oplus(a-1)}\oplus\mathcal{O}(c)$ differ in degree at most
$r-2$. Expanding
$$
P_{\mathcal{O}(c)}(n)={n+c+r\choose r}
$$
shows that
$$
b_{\mathcal{G}}=a\frac{r+1}{2}+c.
$$
The coefficient of degree $r-1$ in
$P_{\mathcal{N}_r(\mathcal{F})}$ is nonnegative. Hence
$b_{\mathcal{F}}\geq b_{\mathcal{G}}\geq0$, which proves the asserted lower
bound.

If the $b_{\mathcal{F}}$ are bounded above, then so are the
$b_{\mathcal{G}}$, and therefore only finitely many pairs $(a,c)$ occur. For
fixed $(a,c)$ the homomorphisms $s$ above form a finite-dimensional affine
parameter space. Moreover, $\mathcal{G}$ is the image of the homomorphism
$$
\mathcal{O}^{\oplus q}\longrightarrow
\SheafHom(\wedge^{a-1}\mathcal{O}^{\oplus q},\mathcal{O}(c)),
\qquad
v\longmapsto(w\mapsto s(v\wedge w)).
$$
Indeed, this is true where $\mathcal{G}$ is locally free, and both sides are
torsion-free. Proposition
\ref{proposition-bounded-coherent-morphisms-extensions} now shows that all
these images form a bounded family. This proves the assertion for $r\geq1$.

If $r=0$, finite projection reduces to $\mathbf{P}^0_S=S$. The dimensions of
quotients of one fixed finite module are bounded, so only finitely many
constant Hilbert polynomials occur. Proposition
\ref{proposition-bounded-coherent-Hilbert-polynomials} finishes the proof.
\end{proof}

\begin{proposition}
\label{proposition-bounded-top-dimensional-pieces}
\begin{reference}
\cite[Expos\'e 221, Lemma 2.6, pp. 257--258]{FGA}.
\end{reference}
Let $S$ be a Noetherian scheme, let $X\to S$ be projective, and fix a
relatively very ample invertible module $\mathcal{O}_X(1)$. Let $E$ be a
family satisfying condition (1) of Proposition
\ref{proposition-bounded-coherent-Hilbert-polynomials}. Suppose that every
member has support of dimension at most $r$. For $r\geq1$, write
$$
P_{\mathcal{F}}(n)=
a_{\mathcal{F}}\frac{n^r}{r!}+
b_{\mathcal{F}}\frac{n^{r-1}}{(r-1)!}+
\text{terms of degree }<r-1.
$$
If the $a_{\mathcal{F}}$ are bounded, then the $b_{\mathcal{F}}$ are bounded
below. If in addition the $b_{\mathcal{F}}$ are bounded above, then the family
of $\mathcal{F}_{(r)}$ is bounded. For $r=0$, boundedness of the constant
coefficients implies that $E$ is bounded.
\end{proposition}

\begin{proof}
We may extend the ground fields and assume they are algebraically closed.
For a member $\mathcal{F}$, let
$$
Z_{\mathcal{F}}=
(\operatorname{Supp}(\mathcal{F}_{(r)}))_{\mathrm{red}}.
$$
By Varieties, Lemma
\ref{varieties-lemma-numerical-polynomial-leading-term},
$$
a_{\mathcal{F}}=
\sum_Z
\mathop{\rm length}_{\mathcal{O}_{X_K,\xi_Z}}
((\mathcal{F}_{(r)})_{\xi_Z})\deg(Z),
$$
where the sum is over the irreducible components of $Z_{\mathcal{F}}$.
Thus boundedness of $a_{\mathcal{F}}$ bounds both the degree of
$Z_{\mathcal{F}}$ and every displayed generic length. Lemma
\ref{lemma-bounded-reduced-equidimensional-degree} shows that the
$Z_{\mathcal{F}}$ form a bounded family.

Choose an integer $A \geq 1$ bounding all the generic lengths, and let
$\mathcal{J}_{\mathcal{F}}$ be the ideal of $Z_{\mathcal{F}}$ in $X_K$.
Then
$$
\mathcal{J}_{\mathcal{F}}^A\mathcal{F}_{(r)}=0.
$$
Indeed, at every associated point of $\mathcal{F}_{(r)}$ this follows from
the elementary fact that the $A$th power of the maximal ideal annihilates a
module of length at most $A$. A nonzero submodule
$\mathcal{J}_{\mathcal{F}}^A\mathcal{F}_{(r)}$ would have an associated point
among those of $\mathcal{F}_{(r)}$, giving a contradiction.

Choose a finite-type family $Z\subset X_T$ which contains all the
$Z_{\mathcal{F}}$ as geometric fibres. The closed subscheme
$$
Y=V(\mathcal{I}_Z^A)\subset X_T
$$
has fibres of dimension at most $r$, and every $\mathcal{F}_{(r)}$ is a module
on a geometric fibre of $Y/T$. It is a quotient of the restriction to that
fibre of the fixed coherent source in condition (1). Apply Lemma
\ref{lemma-bounded-top-dimensional-small-ambient} to $Y/T$. This proves the
lower bound for the $b_{\mathcal{F}}$ and, when they are bounded above, the
boundedness of the $\mathcal{F}_{(r)}$. The same argument and the last
sentence of that lemma handle $r=0$.
\end{proof}

\begin{proposition}
\label{proposition-bounded-dimension-truncations}
\begin{reference}
\cite[Expos\'e 221, Theorem 2.2, pp. 254--258]{FGA}.
\end{reference}
Let $S$ be a Noetherian scheme, let $X\to S$ be projective, and fix a
relatively very ample invertible module $\mathcal{O}_X(1)$. Let $E$ be a
family satisfying condition (1) of Proposition
\ref{proposition-bounded-coherent-Hilbert-polynomials}, and let $s$ be an
integer. Suppose that the coefficients in degrees at least $s-1$ of the
Hilbert polynomials of members of $E$ are bounded. Then the family of
$\mathcal{F}_{(s)}$ is bounded. Moreover, the coefficients in degree $s-2$
are bounded below, where coefficients in negative degrees are understood to
be zero.
\end{proposition}

\begin{proof}
Remark \ref{remark-FGA-Hilbert-coefficient-inequality} explains the correction
of the inequality in the source. If $s\leq0$, all coefficients in nonnegative
degrees are bounded. Numerical polynomials of bounded degree form a lattice
in their coefficient space, so only finitely many Hilbert polynomials occur;
Proposition \ref{proposition-bounded-coherent-Hilbert-polynomials} applies.
We may therefore assume $s\geq1$.

We prove both assertions simultaneously by induction on an upper bound $r$
for the dimensions of the supports. If $r\leq s-2$, then
$\mathcal{F}_{(s)}=0$ and the coefficient in degree $s-2$ is either zero or a
leading coefficient, hence is nonnegative. If $r=s-1$, then again
$\mathcal{F}_{(s)}=0$. When $s=1$ the coefficient in degree $s-2=-1$ is zero
by convention. When $s\geq2$, Proposition
\ref{proposition-bounded-top-dimensional-pieces}, applied with $r=s-1$,
gives the required lower bound from the bounded coefficient in degree $s-1$.

Assume $r\geq s$ and the result is known with upper bound $r-1$. The
coefficients $a_{\mathcal{F}}$ and $b_{\mathcal{F}}$ in degrees $r$ and
$r-1$ are bounded. Proposition
\ref{proposition-bounded-top-dimensional-pieces} shows that the family of
$\mathcal{F}_{(r)}$ is bounded. Let $\mathcal{L}_K\to\mathcal{F}$ be the
fixed-source quotient from condition (1). The kernels of the composites
$$
\mathcal{L}_K\longrightarrow\mathcal{F}
\longrightarrow\mathcal{F}_{(r)}
$$
form a bounded family by Proposition
\ref{proposition-bounded-coherent-morphisms-extensions}. Each
$\mathcal{N}_r(\mathcal{F})$ is a quotient of the corresponding kernel, so
the family of the $\mathcal{N}_r(\mathcal{F})$ has a common coherent quotient
source. Since
$$
P_{\mathcal{N}_r(\mathcal{F})}=P_{\mathcal{F}}-P_{\mathcal{F}_{(r)}}
$$
and a bounded family has only finitely many Hilbert polynomials, the
coefficients in degrees at least $s-1$ on the left are bounded. The induction
hypothesis therefore applies to the $\mathcal{N}_r(\mathcal{F})$.

The dimension filtration gives a short exact sequence
$$
0\longrightarrow
(\mathcal{N}_r(\mathcal{F}))_{(s)}\longrightarrow
\mathcal{F}_{(s)}\longrightarrow
\mathcal{F}_{(r)}\longrightarrow0.
$$
Its left and right terms form bounded families, so Proposition
\ref{proposition-bounded-coherent-morphisms-extensions} bounds the middle
terms. The coefficient in degree $s-2$ of the left term is bounded below by
induction, and all coefficients of the right term are bounded. Additivity of
Hilbert polynomials proves the asserted lower bound for $\mathcal{F}$ and
completes the induction.
\end{proof}

\begin{lemma}
\label{corollary-bounded-associated-dimensions}
\begin{reference}
\cite[Expos\'e 221, Corollary 2.3, p. 255]{FGA}.
\end{reference}
In the situation of Proposition
\ref{proposition-bounded-coherent-Hilbert-polynomials}, suppose there are
integers $s\leq r$ such that every associated point of every member
$\mathcal{F}$ has closure of dimension between $s$ and $r$. Then, in addition
to the common-source condition, it is enough to bound the coefficients of
$P_{\mathcal{F}}$ in degrees from $s-1$ through $r$ in order to conclude that
the family is bounded.
\end{lemma}

\begin{proof}
There are no coefficients in degrees greater than $r$, and Lemma
\ref{lemma-coherent-dimension-truncation} gives
$\mathcal{F}_{(s)}=\mathcal{F}$. Thus Proposition
\ref{proposition-bounded-dimension-truncations} applies. The converse follows
from Proposition \ref{proposition-bounded-coherent-Hilbert-polynomials}.
\end{proof}


\begin{lemma}
\label{lemma-quot-qc-over-base}
Let $f : X \to B$ be a separated morphism of finite presentation
of algebraic spaces. Let $\mathcal{F}$ be a finitely presented
$\mathcal{O}_X$-module. Let $\mathcal{L}$ be an invertible
$\mathcal{O}_X$-module ample on $X/B$, see
Divisors on Spaces, Definition
\ref{spaces-divisors-definition-relatively-ample}.
The algebraic space $\Quotfunctor^P_{\mathcal{F}/X/B}$
parametrizing quotients of $\mathcal{F}$
having Hilbert polynomial $P$ with respect to $\mathcal{L}$
is separated of finite presentation over $B$.
\end{lemma}

\begin{proof}
We have already seen that $\Quotfunctor_{\mathcal{F}/X/B} \to B$
is separated and locally of finite presentation, see
Lemma \ref{lemma-quot-s-lfp}. Thus it suffices to show that
the open subspace $\Quotfunctor^P_{\mathcal{F}/X/B}$
of Remark \ref{remark-quot-numerical} is quasi-compact over $B$.

\medskip\noindent
The question is \'etale local on $B$
(Morphisms of Spaces, Lemma \ref{spaces-morphisms-lemma-quasi-compact-local}).
Thus we may assume $B$ is affine.

\medskip\noindent
Assume $B = \Spec(\Lambda)$. Write $\Lambda = \colim \Lambda_i$ as the
colimit of its finite type $\mathbf{Z}$-subalgebras. Then
we can find an $i$ and a system $X_i, \mathcal{F}_i, \mathcal{L}_i$
as in the lemma over $B_i = \Spec(\Lambda_i)$ whose base change to
$B$ gives $X, \mathcal{F}, \mathcal{L}$.
This follows from
Limits of Spaces, Lemmas
\ref{spaces-limits-lemma-descend-finite-presentation} (to find $X_i$),
\ref{spaces-limits-lemma-descend-modules-finite-presentation} (to find
$\mathcal{F}_i$), \ref{spaces-limits-lemma-descend-invertible-modules}
(to find $\mathcal{L}_i$), and \ref{spaces-limits-lemma-descend-separated}
(to make $X_i$ separated). Because
$$
\Quotfunctor_{\mathcal{F}/X/B} = B \times_{B_i}
\Quotfunctor_{\mathcal{F}_i/X_i/B_i}
$$
and similarly for $\Quotfunctor^P_{\mathcal{F}/X/B}$ we reduce
to the case discussed in the next paragraph.

\medskip\noindent
Assume $B$ is affine and Noetherian. We may replace $\mathcal{L}$
by a positive power, see Lemma \ref{lemma-quot-power-invertible}.
Thus we may assume there exists an immersion $i : X \to \mathbf{P}^n_B$
such that $i^*\mathcal{O}_{\mathbf{P}^n}(1) = \mathcal{L}$. By
Morphisms, Lemma \ref{morphisms-lemma-quasi-compact-immersion}
there exists a closed subscheme $X' \subset \mathbf{P}^n_B$
such that $i$ factors through an open immersion $j : X \to X'$.
By Properties, Lemma \ref{properties-lemma-lift-finite-presentation}
there exists a finitely presented $\mathcal{O}_{X'}$-module
$\mathcal{G}$ such that $j^*\mathcal{G} = \mathcal{F}$.
Thus we obtain an open immersion
$$
\Quotfunctor_{\mathcal{F}/X/B}
\longrightarrow
\Quotfunctor_{\mathcal{G}/X'/B}
$$
by Lemma \ref{lemma-quot-better-open}. Clearly this open immersion
sends $\Quotfunctor^P_{\mathcal{F}/X/B}$ into
$\Quotfunctor^P_{\mathcal{G}/X'/B}$. Now
$\Quotfunctor^P_{\mathcal{G}/X'/B}$ is proper over $B$ by
Lemma \ref{lemma-quot-proper-over-base}.
Therefore it is Noetherian and since any open of a Noetherian
algebraic space is quasi-compact we win.
\end{proof}




\section{Properties of the Hilbert functor}
\label{section-hilb}

\noindent
Let $f : X \to B$ be a morphism of algebraic spaces which is
separated and of finite presentation. Then
$\Hilbfunctor_{X/B}$ is an algebraic space locally of finite
presentation over $B$. See Quot, Proposition \ref{quot-proposition-hilb}.

\begin{lemma}
\label{lemma-hilb-diagonal-closed}
The diagonal of $\Hilbfunctor_{X/B} \to B$ is a closed immersion
of finite presentation.
\end{lemma}

\begin{proof}
In Quot, Lemma \ref{quot-lemma-hilb-is-quot} we have seen that
$\Hilbfunctor_{X/B} = \Quotfunctor_{\mathcal{O}_X/X/B}$.
Hence this follows from Lemma \ref{lemma-quot-diagonal-closed}.
\end{proof}

\begin{lemma}
\label{lemma-hilb-s-lfp}
The morphism $\Hilbfunctor_{X/B} \to B$ is separated
and locally of finite presentation.
\end{lemma}

\begin{proof}
To check $\Hilbfunctor_{X/B} \to B$ is separated we have to
show that its diagonal is a closed immersion. This
is true by Lemma \ref{lemma-hilb-diagonal-closed}.
The second statement is part of
Quot, Proposition \ref{quot-proposition-hilb}.
\end{proof}

\begin{lemma}
\label{lemma-hilb-existence-part}
Assume $X \to B$ is proper as well as of finite presentation.
Then $\Hilbfunctor_{X/B} \to B$ satisfies the existence part
of the valuative criterion (Morphisms of Spaces, Definition
\ref{spaces-morphisms-definition-valuative-criterion}).
\end{lemma}

\begin{proof}
In Quot, Lemma \ref{quot-lemma-hilb-is-quot} we have seen that
$\Hilbfunctor_{X/B} = \Quotfunctor_{\mathcal{O}_X/X/B}$.
Hence this follows from Lemma \ref{lemma-quot-existence-part}.
\end{proof}

\begin{lemma}
\label{lemma-hilb-open}
Let $B$ be an algebraic space. Let $\pi : X \to Y$ be an open immersion
of algebraic spaces which are separated and of finite presentation over $B$.
Then $\pi$ induces an open immersion
$\Hilbfunctor_{X/B} \to \Hilbfunctor_{Y/B}$.
\end{lemma}

\begin{proof}
Omitted. Hint: If $Z \subset X_T$ is a closed subscheme which is
proper over $T$, then $Z$ is also closed in $Y_T$. Thus we obtain
the transformation $\Hilbfunctor_{X/B} \to \Hilbfunctor_{Y/B}$.
If $Z \subset Y_T$ is an element of $\Hilbfunctor_{Y/B}(T)$
and for $t \in T$ we have $|Z_t| \subset |X_t|$,
then the same is true for $t' \in T$ in a neighbourhood of $t$.
\end{proof}

\begin{lemma}
\label{lemma-hilb-closed}
Let $B$ be an algebraic space. Let $\pi : X \to Y$ be a closed immersion
of algebraic spaces which are separated and of finite presentation
over $B$. Then $\pi$ induces a closed immersion
$\Hilbfunctor_{X/B} \to \Hilbfunctor_{Y/B}$.
\end{lemma}

\begin{proof}
Since $\pi$ is a closed immersion, it is immediate that given a
closed subscheme $Z \subset X_T$,
we can view $Z$ as a closed subscheme of $X_T$. Thus we obtain
the transformation $\Hilbfunctor_{X/B} \to \Hilbfunctor_{Y/B}$.
This transformation is immediately seen to be a monomorphism.
To prove that it is a closed immersion, you can use
Lemma \ref{lemma-quot-quotient} for the map
$\mathcal{O}_Y \to \mathcal{O}_X$ and the identifications
$\Hilbfunctor_{X/B} = \Quotfunctor_{\mathcal{O}_X/X/B}$,
$\Hilbfunctor_{Y/B} = \Quotfunctor_{\mathcal{O}_Y/Y/B}$
of Quot, Lemma \ref{quot-lemma-hilb-is-quot}.
\end{proof}

\begin{remark}[Numerical invariants]
\label{remark-hilb-numerical}
Let $f : X \to B$ be as in the introduction to this section.
Let $I$ be a set and for $i \in I$ let $E_i \in D(\mathcal{O}_X)$ be perfect.
Let $P : I \to \mathbf{Z}$ be a function. Recall that
$\Hilbfunctor_{X/B} = \Quotfunctor_{\mathcal{O}_X/X/B}$, see
Quot, Lemma \ref{quot-lemma-hilb-is-quot}.
Thus we can define
$$
\Hilbfunctor^P_{X/B} = \Quotfunctor^P_{\mathcal{O}_X/X/B}
$$
where $\Quotfunctor^P_{\mathcal{O}_X/X/B}$ is as in
Remark \ref{remark-quot-numerical}. The morphism
$$
\Hilbfunctor^P_{X/B} \longrightarrow \Hilbfunctor_{X/B}
$$
is a flat closed immersion which is an open and closed immersion
for example if $I$ is finite, or $B$ is locally Noetherian, or
$I = \mathbf{Z}$ and $E_i = \mathcal{L}^{\otimes i}$ for some
invertible $\mathcal{O}_X$-module $\mathcal{L}$. In the last case
we sometimes use the notation $\Hilbfunctor^{P, \mathcal{L}}_{X/B}$.
\end{remark}

\begin{lemma}
\label{lemma-hilb-proper-over-base}
Let $f : X \to B$ be a proper morphism of finite presentation
of algebraic spaces. Let $\mathcal{L}$ be an invertible
$\mathcal{O}_X$-module ample on $X/B$, see
Divisors on Spaces, Definition
\ref{spaces-divisors-definition-relatively-ample}.
The algebraic space $\Hilbfunctor^P_{X/B}$
parametrizing closed subschemes
having Hilbert polynomial $P$ with respect to $\mathcal{L}$
is proper over $B$.
\end{lemma}

\begin{proof}
Recall that $\Hilbfunctor_{X/B} = \Quotfunctor_{\mathcal{O}_X/X/B}$, see
Quot, Lemma \ref{quot-lemma-hilb-is-quot}.
Thus this lemma is an immediate consequence of
Lemma \ref{lemma-quot-proper-over-base}.
\end{proof}

\begin{lemma}
\label{lemma-hilb-qc-over-base}
Let $f : X \to B$ be a separated morphism of finite presentation
of algebraic spaces. Let $\mathcal{L}$ be an invertible
$\mathcal{O}_X$-module ample on $X/B$, see
Divisors on Spaces, Definition
\ref{spaces-divisors-definition-relatively-ample}.
The algebraic space $\Hilbfunctor^P_{X/B}$
parametrizing closed subschemes
having Hilbert polynomial $P$ with respect to $\mathcal{L}$
is separated of finite presentation over $B$.
\end{lemma}

\begin{proof}
Recall that $\Hilbfunctor_{X/B} = \Quotfunctor_{\mathcal{O}_X/X/B}$, see
Quot, Lemma \ref{quot-lemma-hilb-is-quot}.
Thus this lemma is an immediate consequence of
Lemma \ref{lemma-quot-qc-over-base}.
\end{proof}

\begin{proposition}
\label{proposition-quot-open-fibre-conditions}
\begin{reference}
\cite[Expos\'e 221, Section 4(b), p. 267]{FGA} and
\cite[Chapter IV, Theorems 12.2.1 and 12.2.4]{EGA4}.
\end{reference}
Let $S$ be a Noetherian scheme, let $X \to S$ be quasi-projective, let
$\mathcal{F}$ be a coherent $\mathcal{O}_X$-module, and fix a Hilbert
polynomial $P$ as in Proposition
\ref{proposition-quot-quasi-projective-over-base}. Write
$$
Q = \Quotfunctor^P_{\mathcal{F}/X/S}
$$
and let $\mathcal{G}$ be the universal quotient on $X_Q$. If
$D \subset \mathbf{Z}$ is finite, then the locus in $Q$ where the
dimensions of all prime cycles associated to the fibre of $\mathcal{G}$
belong to $D$ is open.

If $\mathcal{F} = \mathcal{O}_X$, write $H = \Hilbfunctor^P_{X/S}$ and
let $Z \subset X_H$ be the universal closed subscheme. The following define
open subfunctors of $H$:
\begin{enumerate}
\item the locus where $Z \to H$ is smooth on every point of the fibre;
\item the locus where the fibre of $Z \to H$ is geometrically normal;
\item if $X \to S$ is flat, the locus where
$Z_h \to X_h$ is a local complete intersection for every $h \in H$.
\end{enumerate}
Finite intersections of these loci represent simultaneous conditions.
\end{proposition}

\begin{proof}
By Proposition \ref{proposition-quot-quasi-projective-over-base}, $Q$ is a
locally Noetherian scheme. The support of $\mathcal{G}$ is proper over $Q$,
and $\mathcal{G}$ is finitely presented and flat over $Q$. The first claim
is therefore exactly the openness theorem of
\cite[Chapter IV, Theorem 12.2.1(i)]{EGA4}.

The universal family $Z \to H$ is proper, flat, and of finite presentation.
The smooth and geometrically normal fibre loci are open by
\cite[Chapter IV, Theorem 12.2.4(iii), (iv)]{EGA4}. These properties commute
with field extension, so the resulting opens have the asserted functorial
meaning. For the final locus, apply More on Morphisms of Spaces, Lemma
\ref{spaces-more-morphisms-lemma-where-lci} to
$$
\xymatrix{
Z \ar[rr] \ar[rd] & & X_H \ar[ld] \\
& H.
}
$$
The first vertical morphism is flat, proper, and of finite presentation,
and the second is flat and of finite presentation. Intersections of open
subschemes are open.
\end{proof}

\begin{proposition}
\label{proposition-hilbert-lci-tangent-smooth}
\begin{reference}
\cite[Expos\'e 221, Corollaries 5.3 and 5.4]{FGA}.
\end{reference}
Let $S$ be a Noetherian scheme and let $X \to S$ be flat and
quasi-projective. Fix a Hilbert polynomial $P$ and put
$$
H = \Hilbfunctor^P_{X/S}.
$$
Let $h \in H$ be a point with residue field $k$, let $s \in S$ be its
image, and let $Y \subset X_k$ be the corresponding closed subscheme.
Assume that $Y \to X_k$ is a local complete intersection. Then
$$
T_{H_s,h} = H^0(Y, \mathcal{N}_{Y/X_k})
$$
canonically. If
$H^1(Y, \mathcal{N}_{Y/X_k}) = 0$, then $H \to S$ is smooth at $h$.
\end{proposition}

\begin{proof}
The scheme $H$ is locally of finite type over $S$ by Proposition
\ref{proposition-quot-quasi-projective-over-base}. It is the fixed-polynomial
part of the Quot functor for $\mathcal{O}_X$. If $\mathcal{I}$ is the ideal
of $Y$ in $X_k$, Quot, Lemma
\ref{quot-lemma-cotangent-space-quot} gives
$$
T_{H_s,h} =
\Hom_{\mathcal{O}_{X_k}}(\mathcal{I}, \mathcal{O}_Y).
$$
Every such homomorphism annihilates $\mathcal{I}^2$. Since the conormal
sheaf $\mathcal{I}/\mathcal{I}^2$ is finite locally free on $Y$, the right
hand side is
$H^0(Y, \mathcal{N}_{Y/X_k})$ by the definition of the normal sheaf.

Assume the asserted $H^1$ vanishes. We use the pointwise infinitesimal
lifting criterion in More on Morphisms, Lemma
\ref{more-morphisms-lemma-lifting-along-artinian-at-point}. Let
$A' \to A$ be a small extension of local Artinian rings with residue field
$k$ and kernel $J$, and consider a solid diagram in that criterion centered
at $h$. The map $\Spec(A) \to H$ corresponds to a closed subscheme
$Y_A \subset X_A$ flat over $A$.

The local-complete-intersection subfunctor is open by Proposition
\ref{proposition-quot-open-fibre-conditions}. Since $\Spec(A)$ has one
point, the map $\Spec(A) \to H$ factors through this open and
$Y_A \to X_A$ is a local complete intersection. Locally on $X_A$, choose a
quasi-regular sequence defining $Y_A$ and lift its members to $X_{A'}$.
The morphism $X_{A'} \to \Spec(A')$ is flat and of finite presentation.
More on Algebra, Lemma \ref{more-algebra-lemma-cut-by-koszul-flat} shows
that the lifted equations define a flat lift over $A'$. Thus the quotient
$\mathcal{O}_{X_A} \to \mathcal{O}_{Y_A}$ has lifts locally on $X_A$.

Apply Quot, Lemma \ref{quot-lemma-quotient-lifts-torsor}. Because
$A' \to A$ is a small extension, $J$ is a $k$-vector space, and its
obstruction sheaf is canonically the pushforward from $Y$ of
$$
\mathcal{N}_{Y/X_k} \otimes_k J.
$$
Consequently its first cohomology group is
$$
H^1(Y, \mathcal{N}_{Y/X_k}) \otimes_k J = 0.
$$
The lemma therefore gives a global flat lift $Y_{A'} \subset X_{A'}$.
Finite presentation and properness are unchanged by the nilpotent
thickening, so this is an $A'$-point of $H$ lifting the given $A$-point.
The pointwise infinitesimal lifting criterion now proves that $H \to S$ is
smooth at $h$.
\end{proof}

\begin{lemma}
\label{corollary-finite-etale-Hilbert-open}
Let $S$ be a Noetherian scheme and let $X \to S$ be quasi-projective. For
$r \geq 0$, the functor of closed subschemes of $X_T$ which are finite
\'etale of rank $r$ over $T$ is represented by an open subscheme of
$\Hilbfunctor^r_{X/S}$ and hence by a quasi-projective $S$-scheme.
\end{lemma}

\begin{proof}
The constant Hilbert polynomial $r$ says that every fibre of the universal
family is zero dimensional of length $r$. On the smooth open from
Proposition \ref{proposition-quot-open-fibre-conditions}, the universal
family is smooth of relative dimension zero, hence \'etale by definition.
It is proper with finite
fibres, hence finite by More on Morphisms, Lemma
\ref{more-morphisms-lemma-characterize-finite}, and its finite locally free
rank is $r$. The converse is immediate. Quasi-projectivity follows from
Proposition \ref{proposition-quot-quasi-projective-over-base} and Quot,
Lemma \ref{quot-lemma-hilb-is-quot}.
\end{proof}









\section{Zero cycles and the Grothendieck--Deligne norm map}
\label{section-zero-cycles}

\noindent
The symmetric product is the expected parameter space for effective zero
cycles when the morphism to the base is flat. Without flatness, divided
powers give a better behaved parameter space: they commute with arbitrary
base change. We record the comparison before constructing the norm map.

\begin{theorem}
\label{theorem-divided-power-zero-cycles}
\begin{reference}
\cite[Theorem 3.4.1, Corollary 4.2.5, and Proposition
4.2.6]{Rydh-zero-cycles-I}.
\end{reference}
Let $X \to S$ be a separated morphism of algebraic spaces and let $d \geq 0$.
There is a separated algebraic space
$$
\Gamma^d(X/S) \longrightarrow S
$$
parametrizing effective zero cycles of degree $d$, with the following
properties.
\begin{enumerate}
\item Formation of $\Gamma^d(X/S)$ commutes with arbitrary base change on
$S$.
\item If $S = \Spec(A)$ and $X = \Spec(B)$, then
$$
\Gamma^d(X/S) = \Spec(\Gamma_A^d(B)),
$$
where $\Gamma_A^d(B)$ represents homogeneous multiplicative polynomial
laws of degree $d$ from $B$.
\item Addition of cycles gives a morphism
$$
\Psi : (X/S)^d \longrightarrow \Gamma^d(X/S).
$$
The locus $\Gamma^d(X/S)_{\mathrm{nd}}$ where the corresponding geometric
cycle has $d$ distinct points is open, and $\Psi$ is finite \'etale of degree
$d!$ over this locus.
\item The uniform categorical quotient
$$
\operatorname{Sym}^d_S(X) = (X/S)^d/\mathfrak S_d
$$
exists as a separated algebraic space. There is a canonical morphism
$$
\operatorname{Sym}^d_S(X) \longrightarrow \Gamma^d(X/S)
$$
which is a universal homeomorphism with trivial residue field extensions.
It is an isomorphism if $X \to S$ is flat, and it is an isomorphism over
$\Gamma^d(X/S)_{\mathrm{nd}}$ without a flatness assumption.
\end{enumerate}
For an $A$-algebra $C$, a homogeneous multiplicative polynomial law
$F : B \to C$ of degree $d$ is a collection of maps
$$
F_{A'} : B \otimes_A A' \longrightarrow C \otimes_A A'
$$
natural in the $A$-algebra $A'$, multiplicative and unital, and satisfying
$F_{A'}(a x) = a^dF_{A'}(x)$.
\end{theorem}

\begin{proof}
The affine description is the universal property of the algebra of divided
powers. It also proves arbitrary base change in the affine case. The global
space is obtained from the affine spaces of divided powers using \'etale
presentations. Addition of polynomial laws gives $\Psi$ and shows that the
nondegenerate locus is open. The quotient morphism is integral and has the
same geometric fibres as $\Psi$, which proves the assertion about its
topological and residue-field behaviour. In the affine flat case the
comparison is induced by the isomorphism between divided powers and
symmetric tensors; the general flat case follows from the \'etale-local
construction.
\end{proof}

\noindent
Let $f : X \to S$ be separated and let $\mathcal{F}$ be a finitely presented
quasi-coherent $\mathcal{O}_X$-module. We write
$\Quotfunctor^d_{\mathcal{F}/X/S}$ for the subfunctor of the Quot functor
whose objects over $T \to S$ are quotients
$$
\mathcal{F}_T \longrightarrow \mathcal{G}
$$
such that $\mathcal{G}$ is finitely presented and flat over $T$, its support
is finite over $T$, and $p_*\mathcal{G}$ is locally free of rank $d$, where
$p : X_T \to T$. Similarly, $\Hilbfunctor^d_{X/S}$ denotes the functor of
closed subspaces $Z \subset X_T$ finite locally free of degree $d$ over $T$.

\begin{proposition}
\label{proposition-grothendieck-deligne-norm}
\begin{reference}
\cite[Expos\'e 221, Section 6]{FGA} and
\cite[Sections 3 and 6]{Rydh-zero-cycles-II}.
\end{reference}
There are canonical transformations
$$
\Quotfunctor^d_{\mathcal{F}/X/S} \longrightarrow \Gamma^d(X/S)
\quad\text{and}\quad
\Hilbfunctor^d_{X/S} \longrightarrow \Gamma^d(X/S).
$$
They commute with arbitrary base change on $S$. The second transformation,
obtained from the first by taking $\mathcal{F} = \mathcal{O}_X$, is called
the Grothendieck--Deligne norm map.

If $k$ is an algebraically closed field and $\mathcal{G}$ is a $k$-point of
the first functor, then its image is
$$
\sum_{x \in \operatorname{Supp}(\mathcal{G})}
\operatorname{length}_{\mathcal{O}_{X_k,x}}(\mathcal{G}_x)[x].
$$
\end{proposition}

\begin{proof}
The construction is local on the base. Suppose first that $T = \Spec(A)$
and that the finite support of $\mathcal{G}$ is contained in an affine
$\Spec(B) \subset X_T$. Multiplication gives an $A$-algebra homomorphism
$$
B \longrightarrow \operatorname{End}_A(p_*\mathcal{G}).
$$
For every $A$-algebra $A'$, take the determinant of the induced endomorphism
of the rank $d$ locally free $A'$-module
$p_*\mathcal{G} \otimes_A A'$. These determinants form a homogeneous
multiplicative polynomial law of degree $d$. By Theorem
\ref{theorem-divided-power-zero-cycles}, it determines a morphism
$T \to \Gamma^d(X/S)$. Determinants commute with arbitrary base change.
The construction is independent of the affine neighbourhood and glues
\'etale locally, giving the asserted transformations.

For the formula over $k$, decompose $p_*\mathcal{G}$ into the direct sum of
its local factors. The determinant of multiplication by a function is the
product of its values at the support points, each raised to the length of
the corresponding local factor. This is exactly the displayed cycle.
\end{proof}

\begin{remark}
\label{remark-norm-symmetric-morphism}
\begin{reference}
\cite[Expos\'e 221, Section 6]{FGA}.
\end{reference}
Assume that $X \to S$ is flat and let $Y \to S$ be an algebraic space. A
symmetric morphism $u : (X/S)^d \to Y$ factors uniquely through the
categorical quotient $\operatorname{Sym}^d_S(X)$. Since
$\operatorname{Sym}^d_S(X) = \Gamma^d(X/S)$ by Theorem
\ref{theorem-divided-power-zero-cycles}, composition with Proposition
\ref{proposition-grothendieck-deligne-norm} gives the operation customarily
denoted $\pi^u$. If $X = Y$ is a commutative monoid and $u$ is
multiplication, this specializes to the usual norm operation. For a nonflat
$X/S$, the canonical and base-change-compatible target of the determinant
construction is $\Gamma^d(X/S)$, which need not equal the symmetric product.
\end{remark}

\begin{proposition}
\label{proposition-hilbert-symmetric-product}
\begin{reference}
\cite[Definition 4.3.1]{Rydh-hilbert-chow-points}.
\end{reference}
For every separated morphism $X \to S$ there is a canonical transformation
$$
\Hilbfunctor^d_{X/S} \longrightarrow \operatorname{Sym}^d_S(X)
$$
whose composition with
$\operatorname{Sym}^d_S(X) \to \Gamma^d(X/S)$ is the
Grothendieck--Deligne norm map.
\end{proposition}

\begin{proof}
Let $Z \subset X_T$ be finite locally free of degree $d$ over $T$. Since
$Z \to T$ is flat, Theorem \ref{theorem-divided-power-zero-cycles} gives
$$
\operatorname{Sym}^d_T(Z) \xrightarrow{\sim} \Gamma^d(Z/T).
$$
The norm of $\mathcal{O}_Z$ is a $T$-point of the right hand side. Transport
it to the symmetric product and then use the morphism
$\operatorname{Sym}^d_T(Z) \to
\operatorname{Sym}^d_S(X) \times_S T$ induced by $Z \to X_T$.
This construction commutes with pullback in $T$ and has the required
compatibility with the norm map.
\end{proof}

\begin{proposition}
\label{proposition-hilbert-chow-nondegenerate}
\begin{reference}
\cite[Section 6]{Rydh-zero-cycles-II}.
\end{reference}
Let $X \to S$ be separated. The Grothendieck--Deligne norm map induces an
isomorphism from the subfunctor of $\Hilbfunctor^d_{X/S}$ parametrizing
finite \'etale closed subspaces onto
$\Gamma^d(X/S)_{\mathrm{nd}}$. Equivalently, the morphism of Proposition
\ref{proposition-hilbert-symmetric-product} is an isomorphism over the
nondegenerate locus of the symmetric product.
\end{proposition}

\begin{proof}
The assertion is \'etale local on the base. A nondegenerate cycle becomes a
sum of $d$ pairwise distinct sections after an \'etale base change. Their
disjoint union is the unique finite \'etale closed subspace producing the
cycle. This construction descends, and it is inverse to the norm map.
\end{proof}

\begin{proposition}
\label{proposition-hilbert-chow-smooth-curves}
\begin{reference}
\cite[Expos\'e 221, Section 6(b)]{FGA} and
\cite[Section 6]{Rydh-zero-cycles-II}.
\end{reference}
Let $X \to S$ be a separated smooth morphism of schemes of relative
dimension $1$.
Then the canonical morphisms
$$
\Hilbfunctor^d_{X/S}
\longrightarrow \operatorname{Sym}^d_S(X)
\longrightarrow \Gamma^d(X/S)
$$
are isomorphisms.
\end{proposition}

\begin{proof}
The second morphism is an isomorphism because $X \to S$ is flat. A finite
locally free closed subspace of a smooth relative curve is a relative
effective Cartier divisor; see Picard Schemes of Curves, Lemma
\ref{pic-lemma-divisors-on-curves}. Addition of divisors gives an
$\mathfrak S_d$-invariant morphism
$$
(X/S)^d \longrightarrow \Hilbfunctor^d_{X/S}.
$$
It is finite locally free of degree $d!$ by repeated application of Picard
Schemes of Curves, Lemma \ref{pic-lemma-universal-object}. Hence it is an
fpqc covering. It factors through a morphism
$$
\operatorname{Sym}^d_S(X) \longrightarrow \Hilbfunctor^d_{X/S}.
$$
Let $a$ denote this morphism and let $b$ denote the
Hilbert-to-symmetric-product morphism. The equality
$b \circ a = \operatorname{id}$ follows from the categorical quotient
property after composition with $(X/S)^d \to \operatorname{Sym}^d_S(X)$.
The equality $a \circ b = \operatorname{id}$ may be checked after the
displayed fpqc covering of $\Hilbfunctor^d_{X/S}$, where it follows
directly from the construction by adding the $d$ sections.
\end{proof}

\begin{example}
\label{example-hilbert-chow-not-injective}
Let $X$ be smooth of dimension at least $2$ over an algebraically closed
field $k$, and let $x \in X(k)$. Put
$R = \mathcal{O}_{X,x}$ and let $\mathfrak m$ be its maximal ideal. Ideals
$$
\mathfrak m^2 \subset I \subset \mathfrak m
\quad\text{with}\quad
\dim_k(\mathfrak m/I) = 1
$$
give distinct length-two closed subschemes supported at $x$. They are
parametrized by the one-dimensional quotients of
$\mathfrak m/\mathfrak m^2$, while all have norm cycle $2[x]$. Thus the
Hilbert--Chow morphism need not be injective already for $d = 2$. Adding
$d - 2$ fixed reduced points away from $x$ gives the same conclusion for
every $d > 1$.
\end{example}





\section{Relative effective Cartier divisors and complete linear systems}
\label{section-relative-effective-Cartier-divisors}

\noindent
Let $f : X \to S$ be a flat, proper morphism of finite presentation of
schemes. For a scheme $T$ over $S$, set
$$
\mathit{Div}_{X/S}(T) =
\{D \subset X_T \mid D\text{ is a relative effective Cartier divisor
on }X_T/T\}.
$$
Formation of a relative effective Cartier divisor commutes with arbitrary
base change by Divisors, Lemma
\ref{divisors-lemma-relative-Cartier}. Thus this rule defines a subfunctor of
$\Hilbfunctor_{X/S}$.

\begin{lemma}
\label{lemma-divisors-open-in-hilbert}
\begin{reference}
\cite[Expos\'e 232, Section 4, pp. 150--152]{FGA}.
\end{reference}
The transformation
$$
\mathit{Div}_{X/S} \longrightarrow \Hilbfunctor_{X/S}
$$
is representable by open immersions. In particular,
$\mathit{Div}_{X/S}$ is a separated algebraic space locally of finite
presentation over $S$.
\end{lemma}

\begin{proof}
A relative effective Cartier divisor on $X_T/T$ is proper over $T$ because
it is closed in the proper $T$-scheme $X_T$. It is flat and of finite
presentation over $T$, so it gives a point of the Hilbert functor.

Put $H = \Hilbfunctor_{X/S}$ and let $Z \subset X_H$ be the universal
closed subspace. We construct the required open after replacing $H$ by a
scheme $U$ \'etale over it. Then $Z_U \subset X_U$ is a closed subscheme,
$Z_U \to U$ is flat, proper, and of finite presentation, and
$X_U \to U$ is flat and of finite presentation.

Let $W \subset Z_U$ be the set of points $z$ such that the inclusion of the
fibre
$$
(Z_U)_u \subset (X_U)_u
$$
is an effective Cartier divisor in a neighbourhood of $z$. This is open.
Indeed, at such a point Divisors, Lemma
\ref{divisors-lemma-fibre-Cartier} shows that $Z_U$ is a relative effective
Cartier divisor in a neighbourhood of $z$. The converse follows by base
change. Since $Z_U \to U$ is proper, the image of the closed subset
$Z_U \setminus W$ is closed in $U$ by Morphisms, Lemma
\ref{morphisms-lemma-image-proper-scheme-closed}. Its complement $U^0$ has
the following functorial description: a morphism $T \to U$ factors through
$U^0$ if and only if $Z_T \subset X_T$ is a relative effective Cartier
divisor. The opens $U^0$ are compatible with \'etale pullback and hence
descend to the desired open subspace of $H$.

The last assertion follows from Quot, Proposition
\ref{quot-proposition-hilb} and Lemma \ref{lemma-hilb-s-lfp}.
\end{proof}

\begin{proposition}
\label{proposition-divisors-projective-flat}
\begin{reference}
\cite[Expos\'e 232, Proposition 4.1, p. 152]{FGA}.
\end{reference}
Let $S$ be a Noetherian scheme and let $X \to S$ be projective and flat.
Choose an $f$-very ample invertible module $\mathcal{O}_X(1)$. Then
$\mathit{Div}_{X/S}$ is a scheme and there is a decomposition
$$
\mathit{Div}_{X/S} =
\coprod_{P \in \mathbf{Q}[t]} \mathit{Div}^P_{X/S}
$$
into open and closed subschemes. Every $\mathit{Div}^P_{X/S}$ is
quasi-projective over $S$ and is the open Cartier-divisor locus in
$\Hilbfunctor^{P, \mathcal{O}_X(1)}_{X/S}$.
\end{proposition}

\begin{proof}
The Hilbert functor has the corresponding open and closed decomposition by
Remark \ref{remark-hilb-numerical}. Its fixed-polynomial terms are
projective schemes by Proposition
\ref{proposition-quot-projective-over-base}, applied to
$\mathcal{F} = \mathcal{O}_X$. Lemma
\ref{lemma-divisors-open-in-hilbert} identifies
$\mathit{Div}^P_{X/S}$ with an open subscheme of the corresponding
projective Hilbert scheme. It is therefore quasi-projective over $S$.
\end{proof}

\noindent
If $D \subset X_T$ is a relative effective Cartier divisor with invertible
ideal $\mathcal{I}_D$, write
$$
\mathcal{O}_{X_T}(D) = \mathcal{I}_D^{-1}.
$$
This construction commutes with base change and defines the Abel
transformation
\begin{equation}
\label{equation-divisor-to-picard}
a : \mathit{Div}_{X/S} \longrightarrow \Picardfunctor_{X/S},
\qquad
D \longmapsto [\mathcal{O}_{X_T}(D)].
\end{equation}

\begin{lemma}
\label{lemma-universal-module-of-sections}
\begin{reference}
\cite[Expos\'e 232, Section 4, pp. 152--153]{FGA}.
\end{reference}
Let $f : X \to S$ be a proper morphism of finite presentation and let
$\mathcal{F}$ be an $\mathcal{O}_X$-module of finite presentation, flat over
$S$. Put $K = Rf_*\mathcal{F}$ and
$$
\mathcal{Q}_{\mathcal{F}} =
H^0\left(R\SheafHom_{\mathcal{O}_S}(K, \mathcal{O}_S)\right).
$$
Then $\mathcal{Q}_{\mathcal{F}}$ is of finite presentation, its formation
commutes with arbitrary base change, and for every $g : T \to S$ there is a
functorial isomorphism
\begin{equation}
\label{equation-universal-module-of-sections}
f_{T, *}\mathcal{F}_T \cong
\SheafHom_{\mathcal{O}_T}(g^*\mathcal{Q}_{\mathcal{F}}, \mathcal{O}_T).
\end{equation}
\end{lemma}

\begin{proof}
By Derived Categories of Schemes, Lemma
\ref{perfect-lemma-flat-proper-perfect-direct-image-general}, $K$ is perfect,
its formation commutes with arbitrary base change, and every derived pullback
of $K$ has vanishing cohomology in negative degrees. Quot, Lemma
\ref{quot-lemma-cohomology-perfect-complex} shows that locally on $S$ the
complex $K$ has tor amplitude in $[0, b]$ for some $b \geq 0$.

On an affine open of $S$, represent $K$ by a complex
$$
K^0 \longrightarrow K^1 \longrightarrow \ldots \longrightarrow K^b
$$
of finite locally free modules. On this open we have
$$
\mathcal{Q}_{\mathcal{F}} =
\Coker\left((K^1)^\vee \longrightarrow (K^0)^\vee\right).
$$
This description proves finite presentation and commutation with arbitrary
base change. After pulling back to $T$, it also gives
\begin{align*}
\SheafHom_{\mathcal{O}_T}(g^*\mathcal{Q}_{\mathcal{F}}, \mathcal{O}_T)
& = \Ker(g^*K^0 \longrightarrow g^*K^1) \\
& = H^0(Lg^*K) \\
& = f_{T, *}\mathcal{F}_T,
\end{align*}
where the last equality uses the arbitrary base change assertion for $K$.
The construction is canonical, so these local isomorphisms glue.
\end{proof}




\begin{theorem}[Complete linear systems]
\label{theorem-complete-linear-system}
\begin{reference}
\cite[Expos\'e 232, Theorem 4.3, p. 153]{FGA}.
\end{reference}
Let $f : X \to S$ be proper, flat, and of finite presentation with
geometrically integral fibres, and let $\mathcal{L}$ be an invertible
$\mathcal{O}_X$-module. Let $\mathcal{Q} =
\mathcal{Q}_{\mathcal{L}}$ be the module of Lemma
\ref{lemma-universal-module-of-sections}. The functor which associates to
$T \to S$ the relative effective Cartier divisors $D \subset X_T$ for which
there exist an invertible $\mathcal{O}_T$-module $\mathcal{N}$ and an
isomorphism
$$
\mathcal{O}_{X_T}(D) \cong
\mathcal{L}_T \otimes_{\mathcal{O}_{X_T}} f_T^*\mathcal{N}
$$
is represented by the projective bundle $\mathbf{P}(\mathcal{Q})$.
Equivalently, this is the fibre of the Abel transformation
(\ref{equation-divisor-to-picard}) over the class of $\mathcal{L}$.
\end{theorem}

\begin{proof}
For every $T \to S$ we have
$$
\mathcal{O}_T \xrightarrow{\sim} f_{T, *}\mathcal{O}_{X_T}
$$
by Derived Categories of Schemes, Lemma
\ref{perfect-lemma-proper-flat-h0}. Consequently the final assertion about
the fibre follows from Quot, Lemma
\ref{quot-lemma-flat-geometrically-connected-fibres}: two invertible modules
on $X_T$ define the same element of the Picard functor precisely when they
differ by the pullback of an invertible module on $T$.

By Constructions, Definition
\ref{constructions-definition-projective-bundle}, a $T$-point of
$\mathbf{P}(\mathcal{Q})$ is an invertible quotient
$$
g^*\mathcal{Q} \longrightarrow \mathcal{N}.
$$
Dualizing and using
(\ref{equation-universal-module-of-sections}) gives a map
$$
\mathcal{N}^{-1} \longrightarrow f_{T, *}\mathcal{L}_T,
$$
or, by adjunction, a section
$$
s : f_T^*\mathcal{N}^{-1} \longrightarrow \mathcal{L}_T.
$$
The quotient is surjective if and only if the restriction of $s$ to every
geometric fibre is nonzero. Since the geometric fibres of $X \to S$ are
integral, the zero scheme of every such fibrewise section is an effective
Cartier divisor. The zero scheme of $s$ is locally principal on $X_T$, so
Divisors, Lemma \ref{divisors-lemma-fibre-Cartier} shows that it is a relative
effective Cartier divisor $D$ on $X_T/T$. Moreover,
$$
\mathcal{O}_{X_T}(D) \cong
\mathcal{L}_T \otimes f_T^*\mathcal{N}.
$$

Conversely, the canonical section of a divisor with such an isomorphism
gives a fibrewise nonzero section of
$\mathcal{L}_T \otimes f_T^*\mathcal{N}$ and hence an invertible quotient of
$g^*\mathcal{Q}$. These constructions are inverse. Their independence from
the displayed isomorphism follows from
$\mathcal{O}_T^* = f_{T, *}\mathcal{O}_{X_T}^*$, and they commute with base
change. Thus they give the asserted isomorphism of functors.
\end{proof}

\begin{lemma}
\label{corollary-divisor-to-picard-projective}
\begin{reference}
\cite[Expos\'e 232, Corollary 4.4, p. 153]{FGA}.
\end{reference}
Let $S$ be a locally Noetherian scheme and let $X \to S$ be projective and
flat with geometrically integral fibres. Then the Abel transformation
$$
\mathit{Div}_{X/S} \longrightarrow \Picardfunctor_{X/S}
$$
is representable by projective morphisms.
\end{lemma}

\begin{proof}
Let $T \to \Picardfunctor_{X/S}$ be a morphism from a scheme and set
$$
Y = T \times_{\Picardfunctor_{X/S}} \mathit{Div}_{X/S}.
$$
After an fppf covering of $T$, the given Picard class is represented by an
invertible module $\mathcal{L}$ on $X_T$. Theorem
\ref{theorem-complete-linear-system} identifies the corresponding pullback
of $Y$ with $\mathbf{P}(\mathcal{Q}_{\mathcal{L}})$. It follows that
$Y \to T$ is proper by Descent on Spaces, Lemma
\ref{spaces-descent-lemma-descending-property-proper}.

Choose an $f$-very ample invertible module on $X$. Fppf locally on $T$, the
Hilbert polynomial of a divisor in $Y$ is
$$
P_{\mathcal{O}_{X_t}} - P_{\mathcal{L}_t^{-1}}.
$$
It is locally constant by Derived Categories of Schemes, Lemma
\ref{perfect-lemma-chi-locally-constant-geometric}, and it is unchanged when
$\mathcal{L}$ is tensored by the pullback of an invertible module on the
base. Hence it descends to a locally constant function on $T$. Write $T_P$
for its open and closed level sets and $Y_P = Y \times_T T_P$.

The morphism $Y_P \to T_P$ factors through the open subscheme
$\mathit{Div}^P_{X/S}$ of the projective Hilbert scheme
$\Hilbfunctor^{P, \mathcal{O}_X(1)}_{X/S}$. Since $Y_P \to T_P$ is proper
and $\Picardfunctor_{X/S} \to S$ is separated by Lemma
\ref{lemma-pic-functor-uniqueness-part}, $Y_P$ is a closed subspace of
$\mathit{Div}^P_{X/S} \times_S T_P$ and hence is a scheme. The projective
Hilbert scheme is separated over $T_P$, so the resulting immersion of
$Y_P$ into it is a closed immersion; see Morphisms, Lemma
\ref{morphisms-lemma-image-proper-scheme-closed}. Thus $Y_P$ is a projective
$T_P$-scheme. The $T_P$ are pairwise disjoint open and closed subschemes, so
these projective embeddings glue. Therefore $Y \to T$ is projective.
\end{proof}

\begin{remark}[Cohomology and smooth complete linear systems]
\label{remark-complete-linear-system-base-change}
\begin{reference}
\cite[Expos\'e 232, Remark 4.5, pp. 153--154]{FGA}.
\end{reference}
In Theorem \ref{theorem-complete-linear-system}, the coherent module
$\mathcal{Q}$ need not be locally free: the dimension of
$H^0(X_s, \mathcal{L}_s)$ can jump. At a point $s \in S$, the following are
equivalent:
\begin{enumerate}
\item $\mathcal{Q}$ is finite locally free in a neighbourhood of $s$;
\item the base change map
$$
(f_*\mathcal{L}) \otimes_{\mathcal{O}_S} \kappa(s)
\longrightarrow H^0(X_s, \mathcal{L}_s)
$$
is surjective.
\end{enumerate}
This follows directly by representing $Rf_*\mathcal{L}$ near $s$ by the
finite locally free complex used in the proof of Lemma
\ref{lemma-universal-module-of-sections}. If these conditions hold, then
$\mathbf{P}(\mathcal{Q})$ is a projective space bundle and is smooth over
$S$ near its fibre over $s$. Conversely, the elementary affine-chart
criterion for $\mathbf{P}(\mathcal{Q})$ shows that flatness at every point
over $s$ forces $\mathcal{Q}$ to be finite locally free at $s$. Thus the
displayed surjectivity is the exact criterion for smoothness along the whole
complete linear system. The vanishing
$$
H^1(X_s, \mathcal{L}_s) = 0
$$
is a sufficient condition for this surjectivity, by the same finite-complex
calculation and Nakayama's lemma.

Applied after base change to the Abel transformation, this gives the same
criterion for smoothness along one of its complete-linear-system fibres. The
word translated here as ``smooth'' is the historical term ``simple'' in the
printed source.
\end{remark}




\section{The positive Picard locus and the classical quotient construction}
\label{section-positive-picard-locus}

\noindent
Let $S$ be a locally Noetherian scheme and let $f : X \to S$ be projective
and flat with geometrically integral fibres. Fix an $f$-very ample invertible
module $\mathcal{O}_X(1)$ and denote its class in the Picard functor by $\xi$.
For $T \to S$, say that a class $\alpha \in \Picardfunctor_{X/S}(T)$ is
\emph{positive} if, fppf locally on $T$, it is represented by an invertible
module $\mathcal{L}$ such that for every geometric point $\bar t$ of $T$
we have
\begin{align*}
H^i(X_{\bar t}, \mathcal{L}_{\bar t}(n)) & = 0
&& \text{for all }i > 0\text{ and }n \geq 0, \\
H^0(X_{\bar t}, \mathcal{L}_{\bar t}(n)) & \neq 0
&& \text{for all }n \geq 0.
\end{align*}
This condition is unchanged if $\mathcal{L}$ is tensored by the pullback of
an invertible module on $T$. Thus it defines a subfunctor
$\mathit{Pic}^+_{X/S}$ of $\Picardfunctor_{X/S}$.

\begin{lemma}
\label{lemma-positive-picard-open}
\begin{reference}
\cite[Expos\'e 232, Section 5, pp. 154--155]{FGA}.
\end{reference}
The functor $\mathit{Pic}^+_{X/S}$ is represented by an open subspace of
$\Picardfunctor_{X/S}$. Its formation commutes with arbitrary base change,
it satisfies
$$
\mathit{Pic}^+_{X/S} + \xi \subset \mathit{Pic}^+_{X/S},
$$
and
\begin{equation}
\label{equation-picard-covered-positive-translates}
\Picardfunctor_{X/S} =
\bigcup_{n \geq 0}(\mathit{Pic}^+_{X/S} - n\xi).
\end{equation}
\end{lemma}

\begin{proof}
The condition is formulated on geometric fibres, so it commutes with base
change and is fppf local on the base. To prove openness, pull back to a
locally Noetherian scheme $T$ over the Picard functor and then pass to an
fppf cover on which the class is represented by an invertible module
$\mathcal{L}$ on $X_T$.

For every $n$, the complex
$$
K_n = Rf_{T, *}\mathcal{L}(n)
$$
is perfect and commutes with arbitrary base change by Derived Categories of
Schemes, Lemma
\ref{perfect-lemma-flat-proper-perfect-direct-image-general}. The dimensions
of its fibre cohomology are upper semi-continuous by Lemma
\ref{perfect-lemma-jump-loci-geometric}. At a positive geometric point
$\bar t$, relative Serre vanishing, Cohomology of Schemes, Lemma
\ref{coherent-lemma-kill-by-twisting}, gives an integer $N$ and a
neighbourhood on which the higher fibre cohomology vanishes for every
$n \geq N$. For the finitely many $0 \leq n < N$, upper semi-continuity
gives the same conclusion after shrinking again.

On this neighbourhood the dimension of $H^0$ equals the Euler
characteristic and is therefore locally constant by Lemma
\ref{perfect-lemma-chi-locally-constant-geometric}. Shrink once more so that
it is positive for $0 \leq n \leq N$. If $n > N$, multiply a nonzero
section of $\mathcal{L}_{\bar t}(N)$ by a nonzero section of
$\mathcal{O}_{X_{\bar t}}(n - N)$. The product is nonzero because
$X_{\bar t}$ is integral. Thus positivity is open. The opens obtained on an
fppf cover descend, proving the first assertion.

If $\mathcal{L}$ is positive, then $\mathcal{L}(1)$ is positive, which proves
the inclusion under translation. Finally, for any invertible module
$\mathcal{M}$ on a geometric fibre, Serre vanishing and global generation
show that $\mathcal{M}(n)$ is positive for all sufficiently large $n$.
Consequently every geometric point of the Picard functor belongs to one of
the opens on the right hand side of
(\ref{equation-picard-covered-positive-translates}), proving the equality.
\end{proof}

\begin{lemma}
\label{lemma-positive-divisor-cover}
\begin{reference}
\cite[Expos\'e 232, Section 5, pp. 154--155]{FGA}.
\end{reference}
Set
$$
\mathit{Div}^+_{X/S} =
\mathit{Div}_{X/S} \times_{\Picardfunctor_{X/S}}
\mathit{Pic}^+_{X/S}.
$$
Then $\mathit{Div}^+_{X/S}$ is an open subscheme of
$\mathit{Div}_{X/S}$ and the Abel transformation restricts to a smooth,
projective, and surjective morphism
\begin{equation}
\label{equation-positive-divisors-to-positive-picard}
q : \mathit{Div}^+_{X/S} \longrightarrow \mathit{Pic}^+_{X/S}.
\end{equation}
There are open and closed decompositions
$$
\mathit{Div}^+_{X/S} = \coprod_P \mathit{Div}^{+,P}_{X/S},
\qquad
\mathit{Pic}^+_{X/S} = \coprod_P \mathit{Pic}^{+,P}_{X/S},
$$
where $P$ is the Hilbert polynomial of the divisor and
$\mathit{Pic}^{+,P}_{X/S}$ is the image of
$\mathit{Div}^{+,P}_{X/S}$. Each $\mathit{Div}^{+,P}_{X/S}$ is
quasi-projective over $S$, and
$$
R^P = \mathit{Div}^{+,P}_{X/S}
\times_{\mathit{Pic}^{+,P}_{X/S}}
\mathit{Div}^{+,P}_{X/S}
$$
is a scheme defining an equivalence relation whose two projections are
smooth and projective.
\end{lemma}

\begin{proof}
The first assertion follows from Lemma \ref{lemma-positive-picard-open} and
Proposition \ref{proposition-divisors-projective-flat}. Fppf locally on a
scheme mapping to $\mathit{Pic}^+_{X/S}$, the Picard class is represented by
an invertible module $\mathcal{L}$. Positivity and cohomology and base change
show that $f_*\mathcal{L}$ is finite locally free of positive rank and
commutes with base change. Lemma
\ref{lemma-universal-module-of-sections} identifies
$\mathcal{Q}_{\mathcal{L}}$ with $(f_*\mathcal{L})^\vee$. Theorem
\ref{theorem-complete-linear-system} therefore identifies the pullback of
$q$ with the projective space bundle
$\mathbf{P}((f_*\mathcal{L})^\vee)$. This proves that $q$ is smooth,
projective, and surjective.

The decomposition of the divisor scheme and its quasi-projectivity were
proved in Proposition \ref{proposition-divisors-projective-flat}. All
divisors above the same Picard class have the same Hilbert polynomial: on a
geometric fibre the exact sequence of a divisor gives
$$
P_{\mathcal{O}_D} =
P_{\mathcal{O}_X} - P_{\mathcal{L}^{-1}}.
$$
Since $q$ is open, its fixed-polynomial images are open; since they are
pairwise disjoint and cover, they are also closed. The projections from
$R^P$ are base changes of $q$ and hence are smooth and projective. Finally,
the separatedness of the Picard functor, Lemma
\ref{lemma-pic-functor-uniqueness-part}, identifies $R^P$ with a closed
subspace of the scheme
$\mathit{Div}^{+,P}_{X/S} \times_S \mathit{Div}^{+,P}_{X/S}$.
Thus $R^P$ is a scheme.
\end{proof}

\begin{remark}[The extra scheme-quotient step]
\label{remark-classical-positive-divisor-quotient}
\begin{reference}
\cite[Expos\'e 232, Section 5 and Remark 5.1, pp. 154--155]{FGA}.
\end{reference}
Lemma \ref{lemma-positive-divisor-cover} gives an fppf presentation of every
$\mathit{Pic}^{+,P}_{X/S}$ by a quasi-projective scheme with a projective
flat equivalence relation. Bootstrap, Theorem
\ref{bootstrap-theorem-final-bootstrap}, proves from this presentation that
the quotient is an algebraic space. The classical proof invokes the stronger
effective-quotient theorem for a projective flat equivalence relation on a
quasi-projective scheme. That theorem makes
$$
\mathit{Pic}^{+,P}_{X/S} = \mathit{Div}^{+,P}_{X/S}/R^P
$$
a quasi-projective scheme, compatibly with base change. The disjoint union
over $P$, followed by the translate covering
(\ref{equation-picard-covered-positive-translates}), is the additional step
which gives the scheme assertion in Remark
\ref{remark-classical-projective-picard-theorem}. It must not be inferred
from algebraic-space bootstrap alone.

The source describes this as Matsusaka's construction. Its original
effective quotient can also be constructed from the Hilbert scheme; older
constructions used Chow coordinates.
\end{remark}

\begin{remark}[The corrected nonintegral boundary]
\label{remark-corrected-picard-existence-boundary}
\begin{reference}
\cite[Expos\'e 232, Remark 5.2, pp. 155--156;
correction to Expos\'e 232, p. 303]{FGA}.
\end{reference}
The printed remark conjectures that properness, flatness, and the
base-change-stable equality
$$
\mathcal{O}_T \xrightarrow{\sim} f_{T, *}\mathcal{O}_{X_T}
$$
should suffice for existence of the Picard \emph{scheme}. The linked
correction says explicitly that this conjecture is false; it adds that
Mumford's methods prove a slightly weaker theorem. The same correction
inserts the omitted hypothesis ``with algebraically closed residue field''
in the printed example involving a complete local base and says that
Mumford's counterexample makes that restriction indispensable.

There is no conflict with Quot, Proposition
\ref{quot-proposition-pic-functor}: under the displayed universal
global-functions hypothesis the current theorem represents the Picard
functor by an algebraic space, not necessarily by a scheme. The divisor
argument also locates the obstruction. Without geometrically integral
fibres, fibrewise nonzero sections cutting out relative effective Cartier
divisors form only an open subfunctor of the projective bundle in Theorem
\ref{theorem-complete-linear-system}. The resulting equivalence relation can
be flat without being proper, so the classical scheme-quotient step above no
longer applies.
\end{remark}

\begin{remark}[Picard spaces before abelian schemes]
\label{remark-picard-before-abelian-schemes}
\begin{reference}
\cite[Expos\'e 232, Remarks 5.3 and 5.4, pp. 156--157]{FGA}.
\end{reference}
The construction above does not use a prior construction of Jacobians or the
theory of abelian schemes. This order is useful: Picard spaces exist in
situations where the answer is not an abelian scheme, as already happens for
generalized Jacobians of singular curves, and Picard theory in turn supplies
the natural setting for duality of abelian schemes.

For a projective flat family of curves with geometrically integral fibres,
the classical theorem applies even when the family is not smooth. The
separate smooth-curve construction is discussed in Picard Schemes of Curves,
Proposition \ref{pic-proposition-pic-curve}. Statements in the 1962 remark
about then-unknown reducible-fibre cases and compactifications are historical
status reports, not current nonexistence assertions; the modern
algebraic-space theorem is Quot, Proposition
\ref{quot-proposition-pic-functor}.
\end{remark}




\section{Reductions for Picard spaces}
\label{section-reductions-picard-spaces}

\noindent
The results in this section explain how representability of a Picard functor
can be transported through a finite Stein factor, a union of closed
subschemes, and a nilpotent thickening. We also record two nonflat descent
theorems over a field. The word ``scheme'' in the last two results is
essential: algebraic-space representability alone does not contain these
statements.

\begin{proposition}[Picard functors and a finite Stein factor]
\label{proposition-picard-stein-restriction}
\begin{reference}
\cite[Expos\'e 232, Proposition 6.1 and Corollary 6.2,
pp. 157--158]{FGA}.
\end{reference}
Let $S$ be a locally Noetherian scheme and let $X \to S' \to S$ be
morphisms of schemes such that
\begin{enumerate}
\item $X \to S'$ is projective and flat with geometrically integral fibres,
and
\item $S' \to S$ is finite locally free.
\end{enumerate}
Then there is a canonical isomorphism of fppf sheaves
\begin{equation}
\label{equation-picard-stein-restriction}
\Picardfunctor_{X/S} \longrightarrow
\mathop{\rm Res}\nolimits_{S'/S}(\Picardfunctor_{X/S'}).
\end{equation}
Moreover, both sides are represented by schemes.
\end{proposition}

\begin{proof}
For a scheme $T$ over $S$ we have a canonical equality
$$
X \times_S T = X \times_{S'}(S' \times_S T).
$$
It identifies the isomorphism-class presheaves of invertible modules which
occur on the two sides of
(\ref{equation-picard-stein-restriction}). The identification is compatible
with fppf pullback, and hence gives the displayed isomorphism after
sheafification.

By Remark \ref{remark-classical-projective-picard-theorem}, the Picard
functor over $S'$ is a scheme covered by open subschemes which are
quasi-projective over $S'$. The classical construction also has the
finite-subset property used in the cited Proposition 6.1: every finite set
in a fibre is contained in an affine open. Restriction of scalars of any one
of these opens is a scheme by Criteria for Representability, Proposition
\ref{criteria-proposition-restriction-of-scalars-properties}.
These restrictions of scalars cover the right hand side. Indeed, over a
point of $S$ the finite fibre of $S' \to S$ has finite image in the Picard
scheme. Such an image is contained in one of the quasi-projective opens,
and the same containment holds over an open neighbourhood of the point
because $S' \to S$ is finite. Thus the algebraic space supplied by Criteria
for Representability, Proposition
\ref{criteria-proposition-restriction-of-scalars-algebraic-space} is covered
by open subspaces which are schemes, and is therefore a scheme.
\end{proof}

\begin{lemma}
\label{corollary-picard-geometrically-locally-integral}
Let $S$ be a locally Noetherian scheme and let $X \to S$ be projective and
flat with geometrically locally integral fibres. Then
$\Picardfunctor_{X/S}$ is represented by a scheme.
\end{lemma}

\begin{proof}
Let $X \to S' \to S$ be the Stein factorization. The morphism $S' \to S$
is finite \'etale by More on Morphisms of Spaces, Lemma
\ref{spaces-more-morphisms-lemma-stein-factorization-etale}. Its geometric
fibres index the connected components of the geometric fibres of $X/S$.
A connected locally integral scheme is integral, so $X \to S'$ has
geometrically integral fibres. It is flat, and the result follows from
Proposition \ref{proposition-picard-stein-restriction}.
\end{proof}

\begin{example}
\label{example-picard-disjoint-union}
If $X = \coprod_{i = 1}^r X_i$ is a finite disjoint union over $S$, then
Proposition \ref{proposition-picard-stein-restriction} specializes to
$$
\Picardfunctor_{X/S} = \prod_{i = 1}^r \Picardfunctor_{X_i/S}.
$$
\end{example}

\begin{proposition}[Gluing along a flat closed intersection]
\label{proposition-picard-closed-union-affine}
\begin{reference}
\cite[Expos\'e 232, Proposition 6.3, pp. 158--159]{FGA}. The final letter
in the printed description of the intersection is $X$; the displayed
flatness hypothesis makes the intended base $S$.
\end{reference}
Let $S$ be locally Noetherian, let $X \to S$ be proper, and let
$X_1, X_2 \subset X$ be closed subschemes defined by coherent ideals
$\mathcal{I}_1, \mathcal{I}_2$. Assume that
\begin{enumerate}
\item $\mathcal{I}_1 \cap \mathcal{I}_2 = 0$,
\item $X_1$ and $X_2$ are flat over $S$,
\item $Z = X_1 \cap X_2$ is flat over $S$, and
\item for every $s \in S$ the maps
$$
\kappa(s) \longrightarrow
H^0((X_i)_s, \mathcal{O}_{(X_i)_s}), \qquad i = 1, 2,
$$
are isomorphisms.
\end{enumerate}
Then restriction induces a morphism
\begin{equation}
\label{equation-picard-closed-union}
\Picardfunctor_{X/S} \longrightarrow
\Picardfunctor_{X_1/S} \times_S \Picardfunctor_{X_2/S}
\end{equation}
which is representable by affine morphisms. In particular, if the two
Picard functors on the right are schemes, then so is the Picard functor on
the left.
\end{proposition}

\begin{proof}
The first three assumptions give an exact sequence
$$
0 \longrightarrow \mathcal{O}_X \longrightarrow
\mathcal{O}_{X_1} \oplus \mathcal{O}_{X_2} \longrightarrow
\mathcal{O}_Z \longrightarrow 0.
$$
In particular $X$ and $Z$ are flat and of finite presentation over $S$.
The assertion that a base change of
(\ref{equation-picard-closed-union}) is affine is fppf local on the base.
We may therefore work locally where the two Picard classes are represented
by invertible modules $\mathcal{L}_1$ and $\mathcal{L}_2$. After separating
the open and closed locus on which $Z$ is empty, we may also make an fppf
base change such that $Z \to S$ has a section $\sigma$. This section gives
sections of $X_1$ and $X_2$ and permits the modules to be rigidified; compare
Quot, Lemmas \ref{quot-lemma-flat-geometrically-connected-fibres-with-section}
and \ref{quot-lemma-pic-with-section-stack}.

An invertible module on the scheme-theoretic union $X_1 \cup X_2 = X$ is
then the same as a triple
$$
(\mathcal{L}_1, \mathcal{L}_2,
u : \mathcal{L}_1|_Z \longrightarrow \mathcal{L}_2|_Z),
$$
where $u$ is an isomorphism compatible with the rigidifications. The
isomorphism space between the two restrictions is affine over $Z$. Its
functor of sections over the proper flat morphism $Z \to S$ is affine by
Proposition \ref{proposition-sections-proper-flat}; compatibility at
$\sigma$ cuts out the fibre over the identity and is again affine. Hence
the fibre of (\ref{equation-picard-closed-union}) is affine. These affine
schemes and their descent data descend along the fppf coverings used above,
which proves representability and affineness.
\end{proof}

\begin{lemma}
\label{corollary-picard-reduced-components}
\begin{reference}
\cite[Expos\'e 232, Corollary 6.4, p. 159]{FGA}.
\end{reference}
Let $X$ be a proper geometrically reduced scheme over a field $k$. After a
finite separable extension $k'/k$, write $X_1, \ldots, X_r$ for the
irreducible components of $X_{k'}$. If the Picard functors of the $X_i$ are
represented by schemes, then the Picard functor of $X$ is represented by a
scheme, and
after that extension the canonical morphism
$$
\Picardfunctor_{X_{k'}/k'} \longrightarrow
\prod_{i = 1}^r \Picardfunctor_{X_i/k'}
$$
is affine.
\end{lemma}

\begin{proof}
The finite separable extension can be chosen so that all irreducible
components are geometrically irreducible by Varieties, Lemma
\ref{varieties-lemma-finite-extension-geometrically-irreducible-components}.
Their reduced scheme-theoretic intersections are flat over the field.
Applying Proposition \ref{proposition-picard-closed-union-affine}
successively proves the assertion after this extension. The construction is
equivariant for the finite descent datum; the affine morphism and its affine
source descend. This gives the assertion over $k$.
\end{proof}

\begin{proposition}[Nilpotent thickenings]
\label{proposition-picard-nilpotent-thickening-affine}
\begin{reference}
\cite[Expos\'e 232, Proposition 6.5, p. 159]{FGA}.
\end{reference}
Let $X$ be a proper scheme over a field $k$ and let $X_0 \subset X$ be a
closed subscheme with the same underlying topological space. Then
\begin{equation}
\label{equation-picard-nilpotent-thickening}
\Picardfunctor_{X/k} \longrightarrow \Picardfunctor_{X_0/k}
\end{equation}
is representable by affine morphisms. Consequently, representability of
the Picard functor of $X_0$ by a scheme implies the same for $X$.
\end{proposition}

\begin{proof}
Filtering the nilpotent ideal by its powers reduces the assertion to a first
order thickening with ideal $\mathcal{I}$. More on Morphisms, Lemma
\ref{more-morphisms-lemma-picard-group-first-order-thickening} gives the
deformation sequence
$$
H^1(X_0, \mathcal{I}) \longrightarrow \Pic(X) \longrightarrow
\Pic(X_0) \longrightarrow H^2(X_0, \mathcal{I}).
$$
The same sequence applies after every base change over $k$. Since $X_0$ is
proper, the cohomology groups of $\mathcal{I}$ are finite dimensional and
cohomology commutes with extension of the ground field. Thus, after pulling
back a class of $\Picardfunctor_{X_0/k}$ to a test scheme, its obstruction to
lifting is a section of the vector bundle associated to
$H^2(X_0, \mathcal{I})$. Its zero scheme is affine over the test scheme.
Over that zero scheme the sheaf of lifts is a torsor under the vector group
associated to $H^1(X_0, \mathcal{I})$, and is therefore affine. This
description is compatible with fppf descent and proves the claim. Iterating
over the filtration by powers of the ideal finishes the proof.
\end{proof}

\begin{lemma}[The projective Picard scheme over a field]
\label{corollary-picard-projective-field}
\begin{reference}
\cite[Expos\'e 232, Corollary 6.6, p. 159]{FGA}.
\end{reference}
For every projective scheme $X$ over a field $k$, the fppf Picard functor
$\Picardfunctor_{X/k}$ is represented by a scheme.
\end{lemma}

\begin{proof}
After a finite extension of the ground field, the reduced irreducible
components are covered by the classical projective Picard theorem, Remark
\ref{remark-classical-projective-picard-theorem}. Corollary
\ref{corollary-picard-reduced-components} glues the components and
Proposition \ref{proposition-picard-nilpotent-thickening-affine} restores
the nilpotent structure. The inseparable and descent steps are the Oort
reduction in the cited source.
\end{proof}

\begin{theorem}[Affine pullback under a proper surjection]
\label{theorem-picard-pullback-proper-surjective-affine}
\begin{reference}
\cite[Expos\'e 232, Remark 6.7, p. 160 and correction, p. 303;
Comments on Expos\'e 236, final paragraph, p. 307]{FGA}.
\end{reference}
Let $g : Y \to X$ be a surjective morphism of proper schemes over a field
$k$. Pullback gives an affine morphism
\begin{equation}
\label{equation-picard-proper-surjective-pullback}
g^* : \Picardfunctor_{X/k} \longrightarrow \Picardfunctor_{Y/k}.
\end{equation}
\end{theorem}

\begin{proof}
This is the affirmative nonflat-descent theorem referenced in the correction
to Expos\'e 232. The Comments on Expos\'e 236 first give the relative
finite-type statement and then state explicitly that over the spectrum of a
field the same proof makes the pullback morphism affine. Proposition
\ref{proposition-picard-closed-union-affine} and Proposition
\ref{proposition-picard-nilpotent-thickening-affine} are two elementary
instances of the same theorem.
\end{proof}

\begin{theorem}[Murre]
\label{theorem-picard-proper-field}
\begin{reference}
\cite[Expos\'e 232, Remark 6.6, pp. 159--160 and correction,
p. 303]{FGA}.
\end{reference}
For every proper scheme $X$ over a field $k$, the fppf Picard functor
$\Picardfunctor_{X/k}$ is represented by a scheme.
\end{theorem}

\begin{proof}
The projective case is Corollary \ref{corollary-picard-projective-field}.
After the component and nilpotent reductions, Chow's lemma supplies a proper
surjection $Y \to X$ with $Y$ projective. The Picard functor of $Y$ is a
scheme and Theorem \ref{theorem-picard-pullback-proper-surjective-affine}
makes the pullback from the Picard functor of $X$ affine. This is Murre's
affirmative solution of the question printed in Remark 6.6, as recorded by
the linked correction.
\end{proof}

\begin{lemma}[Laurent invariance for a normal scheme]
\label{lemma-pic-normal-laurent-invariance}
\begin{reference}
\cite[Expos\'e 236, proof of Theorem 2.1(ii), pp. 231--232 and correction,
p. 303]{FGA}.
\end{reference}
Let $X$ be a normal scheme of finite type over a field $k$. The projection
$p : X\times_k\mathbf G_{m,k}\to X$ induces an isomorphism
$$
p^* : \Pic(X)\longrightarrow \Pic(X\times_k\mathbf G_{m,k}).
$$
\end{lemma}

\begin{proof}
The section defined by $t=1$ shows that $p^*$ is injective. Since the
irreducible components of a normal Noetherian scheme are open and closed, we
may assume that $X$ is integral. Put $Y=X\times_k\mathbf G_{m,k}$ and let
$K$ be the function field of $X$.

Represent an invertible module on $Y$ by a Cartier divisor $D$. Its
restriction to
$$
Y_K=\Spec(K[t,t^{-1}])
$$
is principal because $K[t,t^{-1}]$ is a unique factorization domain. After
subtracting the divisor of a rational function, we may therefore assume that
no irreducible component of $D$ dominates $X$.

Every prime divisor of $Y$ which does not dominate $X$ is of the form
$Z\times_k\mathbf G_{m,k}$ for a unique prime divisor $Z$ of $X$. Hence
$D=p^*E$ as a Weil divisor for some Weil divisor $E$ on $X$. The rank one
reflexive module $\mathcal{O}_X(E)$ pulls back to the invertible module
$\mathcal{O}_Y(D)$. Invertibility is fpqc local, and $p$ is faithfully flat;
thus $\mathcal{O}_X(E)$ is invertible. Equivalently, $E$ is Cartier. It
follows that the original invertible module is pulled back from $X$, proving
surjectivity.
\end{proof}

\begin{theorem}[Proper smooth families]
\label{theorem-picard-smooth-proper-pieces}
\begin{reference}
\cite[Expos\'e 236, Theorem 2.1(i) and Corollary 2.4,
pp. 231--232]{FGA}.
\end{reference}
Let $S$ be a locally Noetherian scheme and let $f:X\to S$ be proper and
smooth. Suppose that the fppf Picard functor is represented by a scheme
$$
P=\Picardfunctor_{X/S}.
$$
Then $P\to S$ is separated. Every closed subscheme $Z\subset P$ which is of
finite type over $S$ is proper over $S$.
\end{theorem}

\begin{proof}
The representing scheme $P$ is locally of finite type over $S$. By Limits,
Lemma \ref{limits-lemma-Noetherian-dvr-valuative-separation}, separatedness
may be checked using discrete valuation rings. Consider two lifts to $P$ over
a discrete valuation ring which agree over its fraction field. Equality of
the lifts is fpqc local on the valuation ring, so we may pass to its strict
henselization.

The Stein factorization of the base change of $X$ is finite \'etale by More on
Morphisms of Spaces, Lemma
\ref{spaces-more-morphisms-lemma-stein-factorization-etale}. It therefore
splits the base change of $X$ into a finite disjoint union of proper smooth
families with geometrically connected fibres. These fibres are geometrically
integral. Example \ref{example-picard-disjoint-union} identifies the Picard
functor with the product of their Picard functors, and Lemma
\ref{lemma-pic-functor-uniqueness-part} says that every factor is separated.
The two lifts are therefore equal after strict henselization, hence equal over
the original valuation ring. Thus $P\to S$ is separated.

Let $Z\subset P$ be closed and of finite type over $S$. Given a valuative
diagram for $Z$ over a discrete valuation ring $R$, pass to a faithfully flat
extension $R\subset R'$ of valuation rings over which the generic Picard
class is represented by an invertible module. Lemma
\ref{lemma-pic-existence-part} extends that module over the proper smooth
family and hence gives a point of $P(R')$. On
$\Spec(R'\otimes_RR')$ its two pullbacks agree over the schematically dense
generic fibre. They agree everywhere because $P\to S$ is separated. Thus the
point descends to $P(R)$; it factors through $Z$ because its generic point
does and $Z$ is closed. Uniqueness follows from separatedness. Limits, Lemma
\ref{limits-lemma-Noetherian-dvr-valuative-proper}, now shows that $Z\to S$
is proper.

The classical proof in the cited source obtains the same extension by closing
a Cartier divisor from the generic fibre. In its final sentence the printed
phrase ``divisor on $S$'' has to be read as ``divisor on $X$'': it is the
regular local rings of the smooth scheme $X$ which are factorial.
\end{proof}

\begin{lemma}[The codimension two extension theorem]
\label{lemma-pic-flat-lci-codimension-two-extension}
\begin{reference}
\cite[Expos\'e 236, Remark 2.2, pp. 231--232]{FGA} and
\cite[Expos\'e XI, Corollaire 3.14]{SGA2}.
\end{reference}
Let $R$ be a discrete valuation ring with fraction field $K$. Let $X\to
\Spec(R)$ be proper and flat. Assume every geometric fibre is a local complete
intersection and is smooth at all points of codimension at most $2$. Then
every invertible module on $X_K$ extends to an invertible module on $X$.
\end{lemma}

\begin{proof}
The fibre hypotheses imply that the geometric fibres are normal: a local
complete intersection is Cohen--Macaulay, hence $(S_2)$, and smoothness in
codimension at most $2$ gives $(R_1)$. The total space $X$ is normal as well.
Indeed, at a point of the closed fibre its local ring $A$ is a complete
intersection by Divided Power Algebra, Lemma \ref{dpa-lemma-avramov}, and is
therefore $(S_2)$. At a prime of height at most $1$, either one is on the
generic fibre, where regularity follows from the fibre hypothesis, or one is
on the closed fibre and regularity follows from Algebra, Lemma
\ref{algebra-lemma-flat-over-regular-with-regular-fibre}. Thus $A$ is normal
by Serre's criterion.

Work componentwise. Represent an invertible module on $X_K$ by a Cartier
divisor $D_K$, and let $D$ be its closure in $X$, viewed as a Weil divisor.
It remains to prove that $D$ is Cartier at every point of the closed fibre.
Fix such a point and put $A=\mathcal{O}_{X,x}$. As above, $A$ is a normal
local complete intersection. Let $\mathfrak p\subset A$ have height at most
$3$. If the uniformizer of $R$ is not in $\mathfrak p$, then $D$ is Cartier
at $\mathfrak p$ because it restricts to $D_K$. If the uniformizer is in
$\mathfrak p$, then
$$
\dim(A_\mathfrak p/\pi A_\mathfrak p)
=\dim(A_\mathfrak p)-1\leq 2.
$$
The fibre hypothesis makes this quotient regular, and Algebra, Lemma
\ref{algebra-lemma-flat-over-regular-with-regular-fibre} makes
$A_\mathfrak p$ regular. Hence $D$ is Cartier at $\mathfrak p$. Divisors,
Theorem \ref{divisors-theorem-ci-Cartier-codimension-three} now shows that
$D$ is Cartier at $x$. Thus $\mathcal{O}_X(D)$ is the required extension.
\end{proof}

\begin{theorem}[Proper local-complete-intersection families]
\label{theorem-picard-flat-lci-codimension-two}
\begin{reference}
\cite[Expos\'e 236, Remark 2.2, pp. 231--232]{FGA}.
\end{reference}
Let $S$ be a locally Noetherian scheme and let $f:X\to S$ be proper and flat.
Assume every geometric fibre is a local complete intersection and is smooth
at all points of codimension at most $2$. Suppose the fppf Picard functor is
represented by a scheme
$$
P=\Picardfunctor_{X/S}.
$$
Then $P\to S$ is separated. Every closed subscheme $Z\subset P$ which is of
finite type over $S$ is proper over $S$.
\end{theorem}

\begin{proof}
The scheme $P$ is locally of finite type over $S$. The separation argument of
Theorem \ref{theorem-picard-smooth-proper-pieces} applies without change.
Indeed, after strict henselization the finite \'etale Stein factorization
splits $X$ into proper flat families with geometrically integral fibres, and
Lemma \ref{lemma-pic-functor-uniqueness-part} applies to every factor.

Let $Z\subset P$ be closed and of finite type over $S$, and consider its
valuative criterion over a discrete valuation ring $R$. After a faithfully
flat extension which is again a discrete valuation ring, the generic Picard
class is represented by an invertible module. Lemma
\ref{lemma-pic-flat-lci-codimension-two-extension} extends it over the family.
Separatedness makes the two pullbacks agree on the overlap, so the resulting
point descends to $P(R)$ and factors through $Z$. Uniqueness follows from
separatedness. Limits, Lemma
\ref{limits-lemma-Noetherian-dvr-valuative-proper} proves that $Z\to S$ is
proper.
\end{proof}

\begin{example}[The quadric-cone boundary]
\label{example-picard-quadric-cone-boundary}
\begin{reference}
\cite[Expos\'e 236, Remark 2.2, pp. 231--232]{FGA}.
\end{reference}
Let $k$ be an algebraically closed field of characteristic different from
$2$, let $R=k[[t]]$, and let $X\subset\mathbf P^3_R$ be defined by
$$
x_0x_1-x_2^2+t^2x_3^2=0.
$$
This is a projective flat family whose generic fibre is a smooth split
quadric surface and whose special fibre is a normal quadric cone. Its Picard
functor is represented by a scheme, but it has a finite type closed piece
which is not proper over $R$.
\end{example}

\begin{proof}
Representability follows from the projective geometrically integral-fibre
criterion recorded in Remark
\ref{remark-mumford-picard-scheme-boundary}.

Over $K=k((t))$, put $y=x_2+tx_3$ and $z=x_2-tx_3$. The generic equation is
$x_0x_1-yz=0$, so $X_K\cong\mathbf P^1_K\times\mathbf P^1_K$. Let
$$
\mathcal{L}_K=\mathcal{O}_{\mathbf P^1\times\mathbf P^1}(1,0)
\quad\text{and}\quad
\mathcal{H}=\mathcal{O}_X(1).
$$
Then $\deg(\mathcal{L}_K\cdot\mathcal{H}_K)=1$.

Write $Q=X_k$. The resolution of the vertex of $Q$ is the Hirzebruch surface
$\mathbf F_2$. Pullback on Picard groups is injective. If $C$ is the
exceptional section and $F$ a ruling, then
$C^2=-2$ and the pullback of $\mathcal{O}_Q(1)$ is
$\mathcal{O}_{\mathbf F_2}(C+2F)$. A line bundle
$\mathcal{O}_{\mathbf F_2}(aC+bF)$ in the image has degree zero on $C$,
namely $b=2a$. Conversely, every such class is the pullback of
$\mathcal{O}_Q(a)$. Consequently
$$
\Pic(Q)=\mathbf Z[\mathcal{O}_Q(1)].
$$
In particular, the intersection degree with $\mathcal{H}_k$ of every
invertible module on $Q$ is even.

If $\mathcal{L}_K$ extended after any dominating extension of discrete
valuation rings, constancy of the Hilbert polynomial would make its
intersection degree with $\mathcal{H}$ equal to $1$ on the special fibre.
This contradicts the preceding parity calculation. Hence the corresponding
$K$-point of the Picard scheme has no specialization, even after extending
the valuation ring. It is isolated in the generic Picard scheme; its closure
therefore gives a finite type closed copy of $\Spec(K)$ in the relative
Picard scheme, and this copy is not proper over $R$.

The special fibre is a hypersurface with one singular point of codimension
$2$. Thus it is smooth in codimension at most $1$ and normal, but not smooth
in codimension at most $2$. This proves the sharpness asserted in the cited
remark.
\end{proof}

\begin{remark}[Boundary of the proper-pieces theorem]
\label{remark-picard-proper-pieces-boundary}
\begin{reference}
\cite[Expos\'e 236, Corollary 2.4 and the following comment,
p. 232]{FGA}.
\end{reference}
In Theorem \ref{theorem-picard-smooth-proper-pieces}, suppose that
$$
P=\coprod_i P^{(i)}
$$
and every $P^{(i)}$ is of finite type over $S$. Each $P^{(i)}$ is open and
closed in $P$, so the theorem makes it proper over $S$. Theorem
\ref{theorem-picard-flat-lci-codimension-two} gives the same conclusion under
its weaker local-complete-intersection fibre hypotheses.

Example \ref{example-picard-quadric-cone-boundary} shows that geometrically
normal fibres alone do not suffice for this conclusion. The nonproper point
in that example is represented on the generic split quadric by
$\mathcal{O}(1,0)$. Its class has infinite order in
$$
\mathop{\rm NS}(\mathbf P^1\times\mathbf P^1)=\mathbf Z^2,
$$
so it does not belong to $P^\tau$. Consequently the example does not answer
the narrower question left open in the cited source: whether $P^\tau$ must
be proper for a proper normal family when it is of finite type over the
base.
\end{remark}

\begin{theorem}[The identity component over a field]
\label{theorem-picard-identity-component-proper-normal}
\begin{reference}
\cite[Expos\'e 236, Theorem 2.1(ii), pp. 231--232 and correction,
p. 303]{FGA}.
\end{reference}
Let $X$ be a proper geometrically normal scheme over a field $k$. Let
$$
P=\Picardfunctor_{X/k}
$$
be its Picard scheme, whose existence is Theorem
\ref{theorem-picard-proper-field}. Then the identity component $P^0$ is
proper over $k$.
\end{theorem}

\begin{proof}
Properness may be checked after extending the ground field, and the assertion
is componentwise in $X$. We may therefore assume that $k$ is algebraically
closed and that $X$ is nonempty and connected. Choose a point $x\in X(k)$;
rigidification along $x$ identifies maps to $P$ with rigidified invertible
modules.

Let $H=(P^0)_{red}$. This is a smooth connected subgroup scheme by Groupoids,
Lemmas \ref{groupoids-lemma-reduced-subgroup-scheme-perfect} and
\ref{groupoids-lemma-reduced-group-scheme-prefect-field-characteristic-p-smooth}
in positive characteristic, and by
Lemma \ref{groupoids-lemma-group-scheme-characteristic-zero-smooth} in
characteristic zero. Suppose that $H$ is not proper. Chevalley's structure
theorem for the connected commutative algebraic group $H$ then gives a
positive dimensional connected affine subgroup. Over an algebraically closed
field such a group contains a subgroup isomorphic to $\mathbf G_a$ or
$\mathbf G_m$. In the additive case compose its inclusion with, for example,
$u\mapsto u-1$ on $\mathbf G_m$; in the multiplicative case use the subgroup
itself. In either case there is a nonconstant morphism
$$
\mathbf G_m\longrightarrow H\longrightarrow P.
$$

The corresponding rigidified invertible module on
$X\times_k\mathbf G_m$ is pulled back from $X$ by
Lemma \ref{lemma-pic-normal-laurent-invariance}. Hence the displayed morphism
is constant, a contradiction. Thus $H$ is proper. Finally $P^0$ is of finite
type over $k$ and $H\subset P^0$ is a thickening. More on Morphisms, Lemma
\ref{more-morphisms-lemma-thicken-property-morphisms} shows that $P^0$ is
separated and universally closed because $H$ is; hence $P^0$ is proper.
Properness descends from the algebraic closure, proving the theorem over $k$.
\end{proof}

\begin{theorem}[Normal families]
\label{theorem-picard-normal-relative-components}
\begin{reference}
\cite[Expos\'e 236, Corollary 2.3, p. 232]{FGA}.
\end{reference}
Let $S$ be a locally Noetherian scheme and let $f:X\to S$ be proper and
flat with geometrically normal fibres. Suppose the fppf Picard functor is
represented by a scheme
$$
P=\Picardfunctor_{X/S}.
$$
Then $P\to S$ is separated. The identity-component locus
$$
P^0=\bigcup_{s\in S}P_s^0
$$
is closed in $P$ and, with its reduced induced closed subscheme structure,
is proper over $S$. The torsion-component locus $P^\tau$ is open and closed,
and the residue-characteristic-primary component locus $P^\rho$ is closed.
If all residue fields of $S$ have the same characteristic exponent, then the
prime-to-characteristic component locus $P^\sigma$ is closed as well.
\end{theorem}

\begin{proof}
The scheme $P$ is locally of finite type over $S$. We first prove that it is
separated. By the discrete-valuation criterion, fpqc descent, and strict
henselization, it is enough to prove uniqueness over a strictly henselian
discrete valuation ring. The Stein factorization of $X$ is finite \'etale by
More on Morphisms of Spaces, Lemma
\ref{spaces-more-morphisms-lemma-stein-factorization-etale}, since the
geometric fibres are normal and hence reduced. It splits $X$ into finitely
many proper flat families with geometrically connected normal fibres. Those
fibres are geometrically integral. Example
\ref{example-picard-disjoint-union} identifies $P$ with the product of the
Picard functors of the factors, each of which is separated by Lemma
\ref{lemma-pic-functor-uniqueness-part}. This proves that $P\to S$ is
separated.

For every $s\in S$, Theorem
\ref{theorem-picard-identity-component-proper-normal}, applied over
$\kappa(s)$, shows that $P_s^0$ is proper. More on Morphisms, Lemma
\ref{more-morphisms-lemma-proper-connected-component-neighbourhood} gives,
locally on $S$, closedness of $P^0$ and properness of its reduced induced
structure. Both conclusions are local on the base, so they hold over $S$.

Finally, More on Morphisms, Lemma
\ref{more-morphisms-lemma-torsion-component-locus-open} makes $P^\tau$ open.
Lemma \ref{more-morphisms-lemma-order-in-component-group-constructible},
using the closedness of $P^0$, makes $P^\tau$ and $P^\rho$ closed and also
makes $P^\sigma$ closed under the stated equal-characteristic-exponent
hypothesis.
\end{proof}

\begin{theorem}[Prime-to-$n$ multiplication on a Picard space]
\label{theorem-picard-multiplication-etale}
\begin{reference}
\cite[Expos\'e 236, Theorem 2.5, p. 233]{FGA}. The printed statement
places the point $x$ in $X$. Since the power map is an endomorphism of the
Picard space, the point lies in the Picard space instead.
\end{reference}
Let $B$ be a scheme and let $X \to B$ be a proper flat morphism of schemes
of finite presentation. Suppose that the fppf Picard functor is represented
by an algebraic space
$$
P=\Picardfunctor_{X/B}.
$$
Let $n \geq 1$. Multiplication by $n$
$$
[n] : P \longrightarrow P
$$
is \'etale over the open subscheme of $B$ on which $n$ is invertible. In
particular, it is \'etale at every point of $P$ whose residue characteristic
is prime to $n$.
\end{theorem}

\begin{proof}
First, $P \to B$ is locally of finite presentation. Indeed, the usual limit
argument for invertible modules and fppf descent shows that the fppf Picard
functor commutes with directed limits of affine $B$-schemes: a finite
presentation fppf covering, the invertible modules on its pullback of $X$,
the finite descent data, and all their relations descend to a finite stage.
Thus the assertion follows from Limits, Proposition
\ref{limits-proposition-characterize-locally-finite-presentation}.
Consequently $[n]$ is locally of finite presentation by Morphisms of Spaces,
Lemma \ref{spaces-morphisms-lemma-finite-presentation-permanence}.

We may now replace $B$ by the open subscheme on which $n$ is invertible. Let
$T_0 \subset T$ be a first order thickening of affine schemes over $B$. On
the fppf site of $T$, put $U_0=U\times_TT_0$ for every $U/T$ and consider
the square of presheaves whose value on $U$ is
$$
\xymatrix{
\Pic(X_U) \ar[r]^{[n]} \ar[d] & \Pic(X_U) \ar[d] \\
\Pic(X_{U_0}) \ar[r]^{[n]} & \Pic(X_{U_0}).
}
$$
It is cartesian for every $U$ by More on Morphisms, Lemma
\ref{more-morphisms-lemma-pic-roots-first-order-thickening}. Sheafification
commutes with finite limits by Sites, Lemma
\ref{sites-lemma-sheafification-exact}. After fppf sheafification and
evaluation on $T$, we therefore obtain a cartesian square
$$
\xymatrix{
P(T) \ar[r]^{[n]} \ar[d] & P(T) \ar[d] \\
P(T_0) \ar[r]^{[n]} & P(T_0).
}
$$
Thus $[n]$ is formally \'etale. Since it is locally of finite presentation,
it is \'etale by More on Morphisms of Spaces, Lemma
\ref{spaces-more-morphisms-lemma-etale-formally-etale}.
\end{proof}

\begin{example}[A relative Picard space which is not a scheme]
\label{example-picard-conic-not-scheme}
\begin{reference}
\cite[Expos\'e 236, Section 0(a), p. 221]{FGA}.
\end{reference}
Set
$$
S = \Spec(\mathbf{R}[[t]])
$$
and let $X \subset \mathbf{P}^2_S$ be the closed subscheme defined by
\begin{equation}
\label{equation-picard-conic-not-scheme}
x^2 + y^2 = tz^2.
\end{equation}
The morphism $X \to S$ is a projective flat family of curves with
geometrically reduced and geometrically connected fibres. Its fppf Picard
functor is an algebraic space, but it is not a scheme.
\end{example}

\begin{proof}
Let $S' = \Spec(\mathbf{C}[[t]])$. The morphism $S' \to S$ is finite
\'etale. The explicit Picard calculation in the cited reference represents
$\Picardfunctor_{X_{S'}/S'}$ by a disjoint union of copies of a scheme
obtained from the trait $S'$ by replacing its closed point by infinitely
many origins. The nontrivial element of
$\operatorname{Gal}(\mathbf{C}/\mathbf{R})$ exchanges some of these origins.
The resulting descent datum is not effective in schemes: some of its
two-point orbits have no common affine open neighbourhood.

On the other hand, fppf descent data for algebraic spaces are effective by
Bootstrap, Lemma \ref{bootstrap-lemma-descend-algebraic-space}. Hence the
same descent datum produces an algebraic space over $S$. Since formation of
the fppf Picard functor commutes with base change, the descended algebraic
space represents $\Picardfunctor_{X/S}$. If it were a scheme, its pullback
and descent datum would give the scheme whose nonexistence was just
established. Thus it is not a scheme.
\end{proof}

\begin{remark}[The scheme-level boundary]
\label{remark-mumford-picard-scheme-boundary}
\begin{reference}
\cite[Expos\'e 236, Section 0(a), pp. 221--222]{FGA}.
\end{reference}
The historical adjective ``separable'' for the family in
Example \ref{example-picard-conic-not-scheme} means flat with geometrically
reduced fibres. Thus neither projectivity nor this fibrewise condition makes
the relative Picard algebraic space a scheme.

The cited source also records Mumford's positive scheme-level result: if a
projective flat family has geometrically reduced fibres and every
irreducible component of every fibre is geometrically irreducible over its
residue field, then its Picard functor is represented by a scheme. This uses
a strengthened effective-quotient theorem and is not a consequence of
algebraic-space representability alone. The same discussion points out that
the union of torsion components can still be representable by a scheme when
the full Picard functor is not; the relative torsion-component loci are
studied in More on Morphisms, Section
\ref{more-morphisms-section-connected-components}.
\end{remark}

\begin{example}[N\'eron--Severi jumps obstruct universal openness]
\label{example-picard-neron-severi-jump-not-open}
\begin{reference}
\cite[Expos\'e 236, Remark 1.6, p. 228]{FGA}.
\end{reference}
Let $N \geq 3$, let $M$ be a connected fine modular curve over
$\mathbf{C}$ parametrizing elliptic curves with full level-$N$ structure,
and let $E \to M$ be the universal elliptic curve. Set
$$
X=E\times_M E.
$$
Then $X \to M$ is smooth and projective with geometrically connected fibres,
but the relative Picard scheme $\Picardfunctor_{X/M} \to M$ is not
universally open.
\end{example}

\begin{proof}
For a geometric point $s$ of $M$, the N\'eron--Severi group of
$E_s\times E_s$ identifies, through the product principal polarization, with
the symmetric homomorphisms
$$
E_s\times E_s \longrightarrow (E_s\times E_s)^\vee.
$$
At a non-CM point, where $\operatorname{End}(E_s)=\mathbf{Z}$, this group has
rank $3$:
it is generated rationally by the two factors and the diagonal. At a CM
point it has rank $4$. Concretely, the graph of a complex multiplication
endomorphism supplies an additional divisor class.

That additional class cannot extend to a neighbourhood of the CM point. If
it did, the associated symmetric homomorphism of the abelian surface would
extend and would restrict on the geometric generic fibre to an endomorphism
not contained in the symmetric matrices over $\mathbf{Z}$, whereas the
generic elliptic curve has endomorphism ring $\mathbf{Z}$. Consequently the
connected component of the special Picard fibre indexed by this class does
not meet a neighbouring fibre. Its closure is a vertical irreducible
component of the relative Picard scheme, and an open neighbourhood of its
generic point has image supported at the CM point. Thus the Picard morphism
is not open there, hence is not universally open.

Passing to the level cover is the scheme-theoretic version of the source's
shorthand ``modular family over the $j$-line''; it removes the automorphisms
which prevent a universal elliptic curve from existing over the coarse
$j$-line itself.
\end{proof}

\begin{remark}[Picard spaces of formal neighbourhoods]
\label{remark-picard-formal-neighbourhoods-lefschetz}
\begin{reference}
\cite[Expos\'e 232, Remark 6.8, pp. 160--161 and correction,
p. 303]{FGA}.
\end{reference}
Let $X$ be a smooth projective irreducible scheme over a field $k$, let
$Y \subset X$ be a hyperplane section, and let $Y_n$ be its $n$th
infinitesimal neighbourhood. The classical Lefschetz statement recorded in
the source says that for $n$ sufficiently large the morphism
\begin{equation}
\label{equation-picard-formal-neighbourhood}
\Picardfunctor_{X/k} \longrightarrow \Picardfunctor_{Y_n/k}
\end{equation}
is an isomorphism if $\dim(X) \geq 4$. If $\dim(X) \geq 3$, it induces an
isomorphism on the inverse images of the torsion subgroups of the
N\'eron--Severi groups.

This is precisely a setting in which the nilpotent Picard spaces in
Proposition \ref{proposition-picard-nilpotent-thickening-affine} carry
information that the reduced hyperplane alone does not record. The current
formal-coherent infrastructure is Algebraization, Proposition
\ref{algebraization-proposition-lefschetz} and Proposition
\ref{algebraization-proposition-lefschetz-equivalence}; the deformation of
invertible modules between successive neighbourhoods is controlled by More
on Morphisms, Lemma
\ref{more-morphisms-lemma-picard-group-first-order-thickening}.

The printed source says that $X$ is ``regular''. The linked correction
replaces this by ``smooth over $k$'', which is the hypothesis used here.
The subsequent equal-characteristic description of two-dimensional local
class groups by algebraic groups and the proposed higher-dimensional
pro-algebraic extension are historical applications and expectations; they
are not needed for the statements above.
\end{remark}





\section{Properties of the Picard stack}
\label{section-picard-stack}

\noindent
Let $f : X \to B$ be a morphism of algebraic spaces which is flat,
proper, and of finite presentation. Then the stack
$\Picardstack_{X/B}$ parametrizing invertible sheaves on $X/B$
is algebraic, see Quot, Proposition \ref{quot-proposition-pic}.

\begin{lemma}
\label{lemma-pic-diagonal-affine-fp}
The diagonal of $\Picardstack_{X/B}$ over $B$ is affine
and of finite presentation.
\end{lemma}

\begin{proof}
In Quot, Lemma \ref{quot-lemma-picard-stack-open-in-coh} we have seen that
$\Picardstack_{X/B}$ is an open substack of
$\Cohstack_{X/B}$. Hence this follows from
Lemma \ref{lemma-coherent-diagonal-affine-fp}.
\end{proof}

\begin{lemma}
\label{lemma-pic-qs-lfp}
The morphism $\Picardstack_{X/B} \to B$ is quasi-separated and
locally of finite presentation.
\end{lemma}

\begin{proof}
In Quot, Lemma \ref{quot-lemma-picard-stack-open-in-coh} we have seen that
$\Picardstack_{X/B}$ is an open substack of
$\Cohstack_{X/B}$. Hence this follows from
Lemma \ref{lemma-coherent-qs-lfp}.
\end{proof}

\begin{lemma}
\label{lemma-pic-existence-part}
Assume $X \to B$ is smooth in addition to being proper.
Then $\Picardstack_{X/B} \to B$ satisfies the existence part
of the valuative criterion (Morphisms of Stacks, Definition
\ref{stacks-morphisms-definition-existence}).
\end{lemma}

\begin{proof}
Taking base change, this immediately reduces to the following
problem: given a valuation ring $R$ with fraction field $K$ and
an algebraic space $X$ proper and smooth over $R$ and an invertible
$\mathcal{O}_{X_K}$-module $\mathcal{L}_K$, show there exists
an invertible $\mathcal{O}_X$-module $\mathcal{L}$
whose generic fibre is $\mathcal{L}_K$.
Observe that $X_K$ is Noetherian, separated, and regular
(use Morphisms of Spaces, Lemma
\ref{spaces-morphisms-lemma-finite-presentation-noetherian}
and
Spaces over Fields, Lemma \ref{spaces-over-fields-lemma-smooth-regular}).
Thus we can write
$\mathcal{L}_K$ as the difference in the Picard group of
$\mathcal{O}_{X_K}(D_K)$ and $\mathcal{O}_{X_K}(D'_K)$
for two effective Cartier divisors $D_K, D'_K$ in $X_K$, see
Divisors on Spaces, Lemma
\ref{spaces-divisors-lemma-Noetherian-regular-separated-pic-effective-Cartier}.
Finally, we know that $D_K$ and $D'_K$ are restrictions of
effective Cartier divisors $D, D' \subset X$, see
Divisors on Spaces, Lemma
\ref{spaces-divisors-lemma-smooth-over-valuation-ring-effective-Cartier}.
\end{proof}

\begin{lemma}
\label{lemma-pic-inertia}
Assume $f_{T, *}\mathcal{O}_{X_T} \cong \mathcal{O}_T$ for all
schemes $T$ over $B$. Then the inertia stack of $\Picardstack_{X/B}$
is equal to $\mathbf{G}_m \times \Picardstack_{X/B}$.
\end{lemma}

\begin{proof}
This is explained in Examples of Stacks, Example
\ref{examples-stacks-example-inertia-stack-of-picard}.
\end{proof}

\begin{lemma}
\label{lemma-pic-curves-smooth}
Assume $f : X \to B$ has relative dimension $\leq 1$ in addition to
the other assumptions in this section. Then $\Picardstack_{X/B} \to B$
is smooth.
\end{lemma}

\begin{proof}
We already know that $\Picardstack_{X/B} \to B$ is
locally of finite presentation, see Lemma \ref{lemma-pic-qs-lfp}.
Thus it suffices to show that $\Picardstack_{X/B} \to B$ is
formally smooth, see More on Morphisms of Stacks, Lemma
\ref{stacks-more-morphisms-lemma-smooth-formally-smooth}.
Taking base change, this immediately reduces to the following
problem: given a first order thickening $T \subset T'$
of affine schemes, given $X' \to T'$ proper, flat, of finite
presentation and of relative dimension $\leq 1$, and
for $X = T \times_{T'} X'$ given an invertible $\mathcal{O}_X$-module
$\mathcal{L}$, prove that there exists an invertible
$\mathcal{O}_{X'}$-module $\mathcal{L}'$ whose
restriction to $X$ is $\mathcal{L}$.
Since $T \subset T'$ is a first order thickening, the
same is true for $X \subset X'$, see
More on Morphisms of Spaces, Lemma
\ref{spaces-more-morphisms-lemma-base-change-thickening}.
By More on Morphisms of Spaces, Lemma
\ref{spaces-more-morphisms-lemma-picard-group-first-order-thickening}
we see that it suffices to show $H^2(X, \mathcal{I}) = 0$
where $\mathcal{I}$ is the quasi-coherent ideal cutting out $X$ in $X'$.
Denote $f : X \to T$ the structure morphism.
By Cohomology of Spaces, Lemma
\ref{spaces-cohomology-lemma-higher-direct-images-zero-above-dimension-fibre}
we see that $R^pf_*\mathcal{I} = 0$ for $p > 1$.
Hence we get the desired vanishing by
Cohomology of Spaces, Lemma
\ref{spaces-cohomology-lemma-quasi-coherence-higher-direct-images-application}
(here we finally use that $T$ is affine).
\end{proof}




\section{Properties of the Picard functor}
\label{section-picard-functor}

\noindent
Let $f : X \to B$ be a morphism of algebraic spaces which is flat,
proper, and of finite presentation such that moreover for every $T/B$
the canonical map
$$
\mathcal{O}_T \longrightarrow f_{T, *}\mathcal{O}_{X_T}
$$
is an isomorphism. Then the Picard functor $\Picardfunctor_{X/B}$ is an
algebraic space, see Quot, Proposition \ref{quot-proposition-pic-functor}.
There is a closed relationship with the Picard stack.

\begin{lemma}
\label{lemma-pic-gerbe-over-pic-functor}
The morphism $\Picardstack_{X/B} \to \Picardfunctor_{X/B}$
turns the Picard stack into a gerbe over the Picard functor.
\end{lemma}

\begin{proof}
The definition of $\Picardstack_{X/B} \to \Picardfunctor_{X/B}$ being
a gerbe is given in Morphisms of Stacks, Definition
\ref{stacks-morphisms-definition-gerbe}, which in turn refers to
Stacks, Definition \ref{stacks-definition-gerbe-over-stack-in-groupoids}.
To prove it, we will check conditions (2)(a) and (2)(b) of
Stacks, Lemma \ref{stacks-lemma-when-gerbe}. This follows immediately from
Quot, Lemma \ref{quot-lemma-pic-over-pic}; here is a detailed explanation.

\medskip\noindent
Condition (2)(a).
Suppose that $\xi \in \Picardfunctor_{X/B}(U)$ for some scheme $U$ over $B$.
Since $\Picardfunctor_{X/B}$ is the fppf sheafification of the rule
$T \mapsto \Pic(X_T)$ on schemes over $B$
(Quot, Situation \ref{quot-situation-pic}), we see that there exists an
fppf covering $\{U_i \to U\}$ such that $\xi|_{U_i}$ corresponds
to some invertible module $\mathcal{L}_i$ on $X_{U_i}$.
Then $(U_i \to B, \mathcal{L}_i)$ is an object of
$\Picardstack_{X/B}$ over $U_i$ mapping to $\xi|_{U_i}$.

\medskip\noindent
Condition (2)(b). Suppose that $U$ is a scheme over $B$ and
$\mathcal{L}, \mathcal{N}$ are invertible modules on $X_U$
which map to the same element of $\Picardfunctor_{X/B}(U)$.
Then there exists an fppf covering $\{U_i \to U\}$
such that $\mathcal{L}|_{X_{U_i}}$ is isomorphic to $\mathcal{N}|_{X_{U_i}}$.
Thus we find isomorphisms between
$(U \to B, \mathcal{L})|_{U_i} \to (U \to B, \mathcal{N})|_{U_i}$
as desired.
\end{proof}

\begin{lemma}
\label{lemma-pic-functor-diagonal-qc-immersion}
The diagonal of $\Picardfunctor_{X/B}$ over $B$ is a quasi-compact immersion.
\end{lemma}

\begin{proof}
The diagonal is an immersion by Quot, Lemma \ref{quot-lemma-diagonal-pic}.
To finish we show that the diagonal is quasi-compact.
The diagonal of $\Picardstack_{X/B}$ is quasi-compact
by Lemma \ref{lemma-pic-diagonal-affine-fp} and
$\Picardstack_{X/B}$ is a gerbe over $\Picardfunctor_{X/B}$ by
Lemma \ref{lemma-pic-gerbe-over-pic-functor}.
We conclude by Morphisms of Stacks, Lemma
\ref{stacks-morphisms-lemma-gerbe-diagonal-quasi-compact}.
\end{proof}

\begin{lemma}
\label{lemma-pic-functor-qs-lfp}
The morphism $\Picardfunctor_{X/B} \to B$ is quasi-separated and
locally of finite presentation.
\end{lemma}

\begin{proof}
To check $\Picardfunctor_{X/B} \to B$ is quasi-separated we have to
show that its diagonal is quasi-compact. This is immediate from
Lemma \ref{lemma-pic-functor-diagonal-qc-immersion}.
Since the morphism $\Picardstack_{X/B} \to \Picardfunctor_{X/B}$
is surjective, flat, and locally of finite presentation
(by Lemma \ref{lemma-pic-gerbe-over-pic-functor} and
Morphisms of Stacks, Lemma \ref{stacks-morphisms-lemma-gerbe-fppf})
it suffices to prove that $\Picardstack_{X/B} \to B$
is locally of finite presentation, see
Morphisms of Stacks, Lemma
\ref{stacks-morphisms-lemma-flat-finite-presentation-permanence}.
This follows
from Lemma \ref{lemma-pic-qs-lfp}.
\end{proof}

\begin{lemma}
\label{lemma-pic-functor-multiplication-etale}
Assume $X$ is a scheme in addition to the other assumptions in this
section. Let $n$ be an integer which is invertible on $B$. Then
multiplication by $n$
$$
[n] : \Picardfunctor_{X/B} \longrightarrow \Picardfunctor_{X/B}
$$
is an \'etale morphism.
\end{lemma}

\begin{proof}
The morphism $[n]$ is locally of finite presentation by
Lemma \ref{lemma-pic-functor-qs-lfp} and
Morphisms of Spaces, Lemma
\ref{spaces-morphisms-lemma-finite-presentation-permanence}.
We prove that it is formally \'etale. Let $T_0 \subset T$ be a first
order thickening of affine schemes over $B$. Suppose we have
$$
\xi_0 \in \Picardfunctor_{X/B}(T_0), \qquad
\zeta \in \Picardfunctor_{X/B}(T)
$$
such that $n\xi_0 = \zeta|_{T_0}$.

By Lemma \ref{lemma-pic-gerbe-over-pic-functor} and
Morphisms of Stacks, Lemma
\ref{stacks-morphisms-lemma-gerbe-smooth}, elements of the Picard
functor are represented by invertible modules \'etale locally on the
base; here we also use More on Morphisms, Lemma
\ref{more-morphisms-lemma-etale-nbhd-dominates-smooth}.
Moreover, an \'etale covering of $T_0$ lifts to an \'etale covering of
$T$ by \'Etale Morphisms, Theorem
\ref{etale-theorem-remarkable-equivalence}. Consequently, after
replacing $T$ by an \'etale covering, we may choose
$$
\mathcal{L} \in \Pic(X_T), \qquad
\mathcal{M}_0 \in \Pic(X_{T_0})
$$
representing $\zeta$ and $\xi_0$. By Quot, Lemma
\ref{quot-lemma-flat-geometrically-connected-fibres}, and after a
further Zariski refinement, we may arrange that
$$
n[\mathcal{M}_0] = [\mathcal{L}|_{X_{T_0}}]
\quad\text{in}\quad \Pic(X_{T_0}).
$$

The closed immersion $X_{T_0} \subset X_T$ is a first order
thickening. More on Morphisms, Lemma
\ref{more-morphisms-lemma-pic-roots-first-order-thickening}
gives a unique $\mathcal{M} \in \Pic(X_T)$ restricting to
$\mathcal{M}_0$ and satisfying $n[\mathcal{M}] = [\mathcal{L}]$.

We claim that any two lifts in the original lifting problem are equal.
Their difference is an element
$\delta \in \Picardfunctor_{X/B}(T)$ such that
$\delta|_{T_0} = 0$ and $n\delta = 0$. After \'etale localization it
is represented by some $\mathcal{D} \in \Pic(X_T)$. Using
Quot, Lemma \ref{quot-lemma-flat-geometrically-connected-fibres},
and further Zariski localization, we may arrange that
$\mathcal{D}|_{X_{T_0}}$ and $\mathcal{D}^{\otimes n}$ are trivial.
More on Morphisms, Lemma
\ref{more-morphisms-lemma-torsion-pic-thickening} then shows that
$\mathcal{D}$ is trivial. Thus $\delta$ is zero. It follows that the
local lifts constructed above agree on overlaps. They descend because
$\Picardfunctor_{X/B}$ is an fppf sheaf. Hence $[n]$ is formally
\'etale.

We conclude that $[n]$ is \'etale by More on Morphisms of Spaces,
Lemma \ref{spaces-more-morphisms-lemma-etale-formally-etale}.
\end{proof}

\begin{lemma}
\label{lemma-pic-functor-uniqueness-part}
\begin{reference}
\cite[Expos\'e 232, Theorem 3.1, pp. 148--149]{FGA}.
\end{reference}
Assume the geometric fibres of $X \to B$ are integral
in addition to the other assumptions in this section.
Then $\Picardfunctor_{X/B} \to B$ is separated.
\end{lemma}

\begin{proof}
Since $\Picardfunctor_{X/B} \to B$ is quasi-separated, it suffices
to check the uniqueness part of the valuative criterion, see
Morphisms of Spaces, Lemma
\ref{spaces-morphisms-lemma-valuative-criterion-separatedness}.
This immediately reduces to the following problem: given
\begin{enumerate}
\item a valuation ring $R$ with fraction field $K$,
\item an algebraic space $X$ proper and flat over $R$
with integral geometric fibre,
\item an element $a \in \Picardfunctor_{X/R}(R)$ with
$a|_{\Spec(K)} = 0$,
\end{enumerate}
then we have to prove $a = 0$. Applying
Morphisms of Stacks, Lemma
\ref{stacks-morphisms-lemma-lift-valuation-ring-through-flat-morphism}
to the surjective flat morphism
$\Picardstack_{X/R} \to \Picardfunctor_{X/R}$
(surjective and flat by Lemma \ref{lemma-pic-gerbe-over-pic-functor} and
Morphisms of Stacks, Lemma \ref{stacks-morphisms-lemma-gerbe-fppf})
after replacing $R$ by an extension we may assume
$a$ is given by an invertible $\mathcal{O}_X$-module
$\mathcal{L}$. Since $a|_{\Spec(K)} = 0$ we find
$\mathcal{L}_K \cong \mathcal{O}_{X_K}$ by
Quot, Lemma \ref{quot-lemma-flat-geometrically-connected-fibres}.

\medskip\noindent
Denote $f : X \to \Spec(R)$ the structure morphism.
Let $\eta, 0 \in \Spec(R)$ be the generic and closed point.
Consider the perfect complexes
$K = Rf_*\mathcal{L}$ and $M = Rf_*(\mathcal{L}^{\otimes -1})$
on $\Spec(R)$, see Derived Categories of Spaces, Lemma
\ref{spaces-perfect-lemma-flat-proper-perfect-direct-image-general}.
Consider the functions
$\beta_{K, i}, \beta_{M, i} : \Spec(R) \to \mathbf{Z}$
of Derived Categories of Spaces, Lemma
\ref{spaces-perfect-lemma-jump-loci} associated to $K$ and $M$.
Since the formation of $K$ and $M$ commutes with
base change (see lemma cited above) we find
$\beta_{K, 0}(\eta) = \beta_{M, 0}(\eta) = 1$ by
Spaces over Fields, Lemma
\ref{spaces-over-fields-lemma-proper-geometrically-reduced-global-sections}
and our assumption on the fibres of $f$.
By upper semi-continuity we find
$\beta_{K, 0}(0) \geq 1$ and $\beta_{M, 0}(0) \geq 1$.
By 
Spaces over Fields, Lemma
\ref{spaces-over-fields-lemma-characterize-trivial-pic-integral}
we conclude that the restriction of $\mathcal{L}$
to the special fibre $X_0$ is trivial. In turn this gives
$\beta_{K, 0}(0) = \beta_{M, 0}(0) = 1$ as above.
Then by More on Algebra, Lemma
\ref{more-algebra-lemma-lift-pseudo-coherent-from-residue-field}
we can represent $K$ by a complex of the form
$$
\ldots \to 0 \to R \to R^{\oplus \beta_{K, 1}(0)} \to
R^{\oplus \beta_{K, 2}(0)} \to \ldots
$$
Now $R \to R^{\oplus \beta_{K, 1}(0)}$ is zero
because $\beta_{K, 0}(\eta) = 1$. In other words
$K = R \oplus \tau_{\geq 1}(K)$ in $D(R)$ where $\tau_{\geq 1}(K)$
has tor amplitude in $[1, b]$ for some $b \in \mathbf{Z}$.
Hence there is a global section $s \in H^0(X, \mathcal{L})$
whose restriction $s_0$
to $X_0$ is nonvanishing (again because formation of $K$
commutes with base change). Then $s : \mathcal{O}_X \to \mathcal{L}$
is a map of invertible sheaves whose restriction to $X_0$
is an isomorphism and hence is an isomorphism as desired.
\end{proof}

\begin{remark}[The classical projective Picard theorem]
\label{remark-classical-projective-picard-theorem}
\begin{reference}
\cite[Expos\'e 232, Theorem 3.1 and Remark 3.3, pp. 148--150;
correction to Expos\'e 232, p. 303]{FGA}.
\end{reference}
Let $X$ and $B$ be locally Noetherian schemes and suppose that
$X \to B$ is projective and flat with integral geometric fibres.
The classical projective Picard theorem strengthens
Quot, Proposition \ref{quot-proposition-pic-functor},
Lemma \ref{lemma-pic-functor-qs-lfp}, and
Lemma \ref{lemma-pic-functor-uniqueness-part}: in this situation
$\Picardfunctor_{X/B}$ is represented by a separated scheme.
If $\mathcal{O}_X(1)$ is relatively very ample and $\xi$ is its class,
the classical construction gives an open subspace $U$ which is a disjoint
union of open subschemes quasi-projective over $B$ and such that
$$
U + \xi \subset U
\quad\text{and}\quad
\Picardfunctor_{X/B} = \bigcup_{n \geq 0}(U - n\xi).
$$
Translation identifies every $U - n\xi$ with $U$. The scheme and
quasi-projectivity assertions use the divisor and Hilbert scheme
construction; they do not follow from algebraic-space representability and
local finite presentation alone.

The printed source also asks whether the Picard scheme is a disjoint union
of open subschemes of finite type, and hence quasi-projective in this
setting, over $B$. The correction appended to the collected expos\'es says
that Mumford answered this question affirmatively. Thus the question in the
printed Remark 3.3 is not open.
\end{remark}

\begin{remark}[A separatedness boundary]
\label{remark-picard-separatedness-boundary}
\begin{reference}
\cite[Expos\'e 232, Remark 3.2, p. 149]{FGA}.
\end{reference}
Fibrewise constancy of global functions is not a substitute for the
integrality hypothesis in Lemma
\ref{lemma-pic-functor-uniqueness-part}. More precisely, the condition
$$
\kappa(b) \longrightarrow H^0(X_b, \mathcal{O}_{X_b})
$$
being an isomorphism for every $b \in B$ does not by itself force the
relative Picard functor to be separated. The cited source records examples
over discrete valuation rings in which an integral generic curve specializes
either to two intersecting irreducible components, as for a degeneration of
a conic to two lines, or to a nonreduced multiple fibre with irreducible
support, as for a double elliptic curve. Both kinds of example occur in every
characteristic.
\end{remark}

\begin{remark}[Locally projective families and components]
\label{remark-picard-locally-projective-components}
\begin{reference}
\cite[Expos\'e 232, Remark 3.3, pp. 149--150]{FGA}.
\end{reference}
Replacing projectivity by projectivity locally on the base preserves the
local representability problem, but it does not force connected components
of the Picard scheme to be of finite type. Here is the construction recorded
in the cited source. Let $X_0$ be a smooth projective variety over an
algebraically closed field $k$. Choose an automorphism $u$ and a class
$\xi$ in its N\'eron--Severi group such that the classes $u^n(\xi)$ are
pairwise distinct. One example is $X_0 = E \times_k E$ for an elliptic curve
$E$, with
$$
u(x, y) = (x, y + x).
$$
Let $B$ be the union of two smooth irreducible curves meeting in two points
and let $P \to B$ be a connected principal $\mathbf{Z}$-cover. The family
associated to the action of $\mathbf{Z}$ on $X_0$ is locally projective over
$B$; for the displayed example it is an abelian scheme. Its Picard scheme has
a connected component isomorphic to
$$
P \times_k \underline{\Picardfunctor}^{0}_{X_0/k}.
$$
This component is not of finite type over $B$. In particular, local finite
presentation of the entire Picard space does not imply that each of its
connected components is of finite type.
\end{remark}

\begin{lemma}
\label{lemma-pic-functor-curves-smooth}
Assume $f : X \to B$ has relative dimension $\leq 1$ in addition to
the other assumptions in this section. Then $\Picardfunctor_{X/B} \to B$
is smooth.
\end{lemma}

\begin{proof}
By Lemma \ref{lemma-pic-curves-smooth} we know that
$\Picardstack_{X/B} \to B$ is smooth. The morphism
$\Picardstack_{X/B} \to \Picardfunctor_{X/B}$ is surjective
and smooth by combining Lemma \ref{lemma-pic-gerbe-over-pic-functor} with
Morphisms of Stacks, Lemma \ref{stacks-morphisms-lemma-gerbe-smooth}.
Thus if $U$ is a scheme and $U \to \Picardstack_{X/B}$ is surjective
and smooth, then $U \to \Picardfunctor_{X/B}$ is surjective and smooth
and $U \to B$ is surjective and smooth (because these properties
are preserved by composition). Thus $\Picardfunctor_{X/B} \to B$
is smooth for example by
Descent on Spaces, Lemma
\ref{spaces-descent-lemma-syntomic-smooth-etale-permanence}.
\end{proof}






\section{Properties of relative morphisms}
\label{section-relative-morphisms}

\noindent
Let $B$ be an algebraic space. Let $X$ and $Y$ be algebraic spaces
over $B$ such that $Y \to B$ is flat, proper, and of finite presentation
and $X \to B$ is separated and of finite presentation.
Then the functor $\mathit{Mor}_B(Y, X)$ of relative morphisms
is an algebraic space locally of finite presentation over $B$.
See Quot, Proposition \ref{quot-proposition-Mor}.

\begin{lemma}
\label{lemma-Mor-diagonal-closed}
The diagonal of $\mathit{Mor}_B(Y, X) \to B$ is a closed immersion
of finite presentation.
\end{lemma}

\begin{proof}
There is an open immersion
$\mathit{Mor}_B(Y, X) \to \Hilbfunctor_{Y \times_B X/B}$, see
Quot, Lemma \ref{quot-lemma-Mor-into-Hilb-open}.
Thus the lemma follows from
Lemma \ref{lemma-hilb-diagonal-closed}.
\end{proof}

\begin{lemma}
\label{lemma-Mor-s-lfp}
The morphism $\mathit{Mor}_B(Y, X) \to B$ is separated
and locally of finite presentation.
\end{lemma}

\begin{proof}
To check $\mathit{Mor}_B(Y, X) \to B$ is separated we have to
show that its diagonal is a closed immersion. This
is true by Lemma \ref{lemma-Mor-diagonal-closed}.
The second statement is part of
Quot, Proposition \ref{quot-proposition-Mor}.
\end{proof}

\begin{lemma}
\label{lemma-Isom-in-Mor}
With $B, X, Y$ as in the introduction of this section, in addition
assume $X \to B$ is proper. Then the
subfunctor $\mathit{Isom}_B(Y, X) \subset \mathit{Mor}_B(Y, X)$
of isomorphisms is an open subspace.
\end{lemma}

\begin{proof}
Follows immediately from More on Morphisms of Spaces, Lemma
\ref{spaces-more-morphisms-lemma-where-isomorphism}.
\end{proof}

\begin{remark}[Numerical invariants]
\label{remark-Mor-numerical}
Let $B, X, Y$ be as in the introduction to this section.
Let $I$ be a set and for $i \in I$ let
$E_i \in D(\mathcal{O}_{Y \times_B X})$ be perfect.
Let $P : I \to \mathbf{Z}$ be a function. Recall that
$$
\mathit{Mor}_B(Y, X) \subset
\Hilbfunctor_{Y \times_B X/B}
$$
is an open subspace, see Quot, Lemma \ref{quot-lemma-Mor-into-Hilb-open}.
Thus we can define
$$
\mathit{Mor}^P_B(Y, X) =
\mathit{Mor}_B(Y, X) \cap \Hilbfunctor^P_{Y \times_B X/B}
$$
where $\Hilbfunctor^P_{Y \times_B X/B}$ is as in
Remark \ref{remark-hilb-numerical}. The morphism
$$
\mathit{Mor}^P_B(Y, X) \longrightarrow \mathit{Mor}_B(Y, X)
$$
is a flat closed immersion which is an open and closed immersion
for example if $I$ is finite, or $B$ is locally Noetherian, or
$I = \mathbf{Z}$, $E_i = \mathcal{L}^{\otimes i}$
for some invertible $\mathcal{O}_{Y \times_B X}$-module $\mathcal{L}$.
In the last case we sometimes use the notation
$\mathit{Mor}^{P, \mathcal{L}}_B(Y, X)$.
\end{remark}

\begin{lemma}
\label{lemma-Mor-qc-over-base}
With $B, X, Y$ as in the introduction of this section, let
$\mathcal{L}$ be ample on $X/B$ and let $\mathcal{N}$ be ample on $Y/B$.
See Divisors on Spaces, Definition
\ref{spaces-divisors-definition-relatively-ample}.
Let $P$ be a numerical polynomial. Then
$$
\mathit{Mor}^{P, \mathcal{M}}_B(Y, X) \longrightarrow B
$$
is separated and of finite presentation where
$\mathcal{M} = \text{pr}_1^*\mathcal{N}
\otimes_{\mathcal{O}_{Y \times_B X}} \text{pr}_2^*\mathcal{L}$.
\end{lemma}

\begin{proof}
By Lemma \ref{lemma-Mor-s-lfp} the morphism $\mathit{Mor}_B(Y, X) \to B$
is separated and locally of finite presentation. Thus it suffices to
show that the open and closed subspace $\mathit{Mor}^{P, \mathcal{M}}_B(Y, X)$
of Remark \ref{remark-Mor-numerical} is quasi-compact over $B$.

\medskip\noindent
The question is \'etale local on $B$
(Morphisms of Spaces, Lemma \ref{spaces-morphisms-lemma-quasi-compact-local}).
Thus we may assume $B$ is affine.

\medskip\noindent
Assume $B = \Spec(\Lambda)$. Note that $X$ and
$Y$ are schemes and that $\mathcal{L}$ and $\mathcal{N}$ are ample
invertible sheaves on $X$ and $Y$ (this follows immediately from the
definitions). Write $\Lambda = \colim \Lambda_i$ as the
colimit of its finite type $\mathbf{Z}$-subalgebras. Then
we can find an $i$ and a system $X_i, Y_i, \mathcal{L}_i, \mathcal{N}_i$
as in the lemma over $B_i = \Spec(\Lambda_i)$ whose base change to
$B$ gives $X, Y, \mathcal{L}, \mathcal{N}$. This follows from
Limits, Lemmas
\ref{limits-lemma-descend-finite-presentation} (to find $X_i$, $Y_i$),
\ref{limits-lemma-descend-invertible-modules} (to find $\mathcal{L}_i$,
$\mathcal{N}_i$), \ref{limits-lemma-descend-separated-finite-presentation}
(to make $X_i \to B_i$ separated), \ref{limits-lemma-eventually-proper}
(to make $Y_i \to B_i$ proper), and \ref{limits-lemma-limit-ample}
(to make $\mathcal{L}_i$, $\mathcal{N}_i$ ample).
Because
$$
\mathit{Mor}_B(Y, X) = B \times_{B_i} \mathit{Mor}_{B_i}(Y_i, X_i)
$$
and similarly for $\mathit{Mor}^P_B(Y, X)$ we reduce
to the case discussed in the next paragraph.

\medskip\noindent
Assume $B$ is a Noetherian affine scheme. By
Properties, Lemma \ref{properties-lemma-ample-on-product}
we see that $\mathcal{M}$ is ample. By Lemma \ref{lemma-hilb-qc-over-base}
we see that $\Hilbfunctor^{P, \mathcal{M}}_{Y \times_B X/B}$ is of
finite presentation over $B$ and hence Noetherian.
By construction
$$
\mathit{Mor}^{P, \mathcal{M}}_B(Y, X) =
\mathit{Mor}_B(Y, X) \cap
\Hilbfunctor^{P, \mathcal{M}}_{Y \times_B X/B}
$$
is an open subspace of $\Hilbfunctor^{P, \mathcal{M}}_{Y \times_B X/B}$ and
hence quasi-compact (as an open of a Noetherian algebraic space
is quasi-compact).
\end{proof}

\begin{proposition}
\label{proposition-sections-Mor-quasi-projective}
\begin{reference}
\cite[Expos\'e 221, Section 4(c), pp. 267--268]{FGA}.
\end{reference}
Let $S$ be a Noetherian scheme.
\begin{enumerate}
\item Let $Y \to S$ be flat and proper and let $Z \to Y$ be an
$S$-morphism such that $Z \to S$ is quasi-projective. The functor
$\mathit{Sec}_{Y/S}(Z/Y)$ of sections is a scheme locally of finite type
over $S$. It is an open subscheme of $\Hilbfunctor_{Z/S}$, and each
fixed-Hilbert-polynomial component is quasi-projective over $S$.
\item Let $Y \to S$ be flat and projective and let $X \to S$ be
quasi-projective. Then $\mathit{Mor}_S(Y, X)$ is a scheme locally of finite
type over $S$. The graph construction identifies it with an open subscheme
of $\Hilbfunctor_{Y \times_S X/S}$, and every fixed-polynomial subfunctor
$\mathit{Mor}^{P,\mathcal{M}}_S(Y, X)$ is quasi-projective over $S$, where
$$
\mathcal{M} = \operatorname{pr}_1^*\mathcal{N} \otimes
\operatorname{pr}_2^*\mathcal{L}
$$
for relatively very ample invertible modules $\mathcal{N}$ on $Y/S$ and
$\mathcal{L}$ on $X/S$. The subfunctor of morphisms $Y_T \to X_T$ which
are immersions is open.
\item If $X \to S$ is projective as well, then
$\mathit{Isom}_S(Y, X)$ is an open subscheme of
$\mathit{Mor}_S(Y, X)$.
\end{enumerate}
\end{proposition}

\begin{proof}
A section of $Z_T \to Y_T$ has an image which is a closed subscheme of
$Z_T$, and a closed subscheme $\Gamma \subset Z_T$ is the image of a
section exactly when the projection $\Gamma \to Y_T$ is an isomorphism.
On the universal family over $\Hilbfunctor_{Z/S}$ this is an open condition
by More on Morphisms of Spaces, Lemma
\ref{spaces-more-morphisms-lemma-where-isomorphism}. Thus the section
functor is an open subspace of $\Hilbfunctor_{Z/S}$. Proposition
\ref{proposition-quot-quasi-projective-over-base} and Quot, Lemma
\ref{quot-lemma-hilb-is-quot} show that the fixed-polynomial Hilbert spaces
are quasi-projective schemes. Their open section loci are quasi-compact,
because these Hilbert spaces are Noetherian, and hence are quasi-projective.
Taking their disjoint union proves (1).

For (2), use
$$
\mathit{Mor}_S(Y, X) =
\mathit{Sec}_{Y/S}((Y \times_S X)/Y).
$$
The product $Y \times_S X$ is quasi-projective over $S$: its projection to
$X$ is a base change of the projective morphism $Y \to S$, and one applies
Morphisms, Lemmas \ref{morphisms-lemma-base-change-quasi-projective} and
\ref{morphisms-lemma-composition-quasi-projective}. The statement about
graphs also follows directly from Quot, Lemma
\ref{quot-lemma-Mor-into-Hilb-open}. The displayed invertible module is
relatively very ample: locally on $S$, take the projective embeddings
defined by $\mathcal{N}$ and $\mathcal{L}$ and compose their product with
the Segre embedding of Constructions, Lemma
\ref{constructions-lemma-segre-embedding}. The fixed-polynomial assertion
follows from (1) or from Remark \ref{remark-Mor-numerical}. Finally, every
immersion $Y_T \to X_T$ is closed: it is proper because $Y_T \to T$ is
proper and $X_T \to T$ is separated. The closed-immersion locus of the
universal morphism is open by More on Morphisms of Spaces, Lemma
\ref{spaces-more-morphisms-lemma-where-closed-immersion}.

If $X$ is projective, the isomorphism locus is open by Lemma
\ref{lemma-Isom-in-Mor}. This proves (3).
\end{proof}

\begin{remark}
\label{remark-automorphism-scheme-nilpotents}
\begin{reference}
\cite[Expos\'e 221, final paragraph of Section 4(c), p. 268]{FGA}.
\end{reference}
For a projective scheme over a field, the automorphism functor is represented
by the open subscheme $\mathit{Isom}(X, X)$ above. Its scheme structure is
essential: the reduction seen by classical constructions over an
algebraically closed field can discard infinitesimal automorphisms. Over an
imperfect field, the nilradical appearing after extending the ground field
need not descend, so replacing the automorphism scheme by that reduction is
not a base-field construction.
\end{remark}






\begin{proposition}[Sections over a proper flat source]
\label{proposition-sections-proper-flat}
\begin{reference}
\cite[Expos\'e 221, Section 7]{FGA}.
\end{reference}
Let $Y \to S$ be a flat proper morphism of schemes of finite presentation.
For a morphism $Z \to Y$, let $\mathit{Sec}_{Y/S}(Z/Y)$ be the functor
which assigns to $T \to S$ the set of sections of $Z_T \to Y_T$.
\begin{enumerate}
\item If $Z \to Y$ is separated and of finite presentation, then
$\mathit{Sec}_{Y/S}(Z/Y)$ is a separated algebraic space locally of finite
presentation over $S$.
\item If $Z \to Y$ is a closed immersion, then
$\mathit{Sec}_{Y/S}(Z/Y) \to S$ is represented by a closed immersion. It is
of finite presentation if $Z \to S$ is of finite presentation.
\item If $Z \to Y$ is affine, then $\mathit{Sec}_{Y/S}(Z/Y)$ is an affine
scheme over $S$. If in addition $Z \to Y$ is of finite presentation, then
this affine scheme is of finite presentation over $S$.
\item Let $\mathcal{E}$ be a finite locally free $\mathcal{O}_Y$-module. The
functor of sections of the vector bundle
$\mathbf{V}(\mathcal{E}) \to Y$ is a vector bundle over $S$ in the sense of
Constructions, Definition \ref{constructions-definition-vector-bundle}. Its
associated quasi-coherent module is of finite presentation.
\end{enumerate}
\end{proposition}

\begin{proof}
For (1), composition with $Z \to Y$ and the identity of $Y$ give a
cartesian diagram
$$
\xymatrix{
\mathit{Sec}_{Y/S}(Z/Y) \ar[r] \ar[d] &
\mathit{Mor}_S(Y, Z) \ar[d] \\
S \ar[r]^-{\operatorname{id}_Y} & \mathit{Mor}_S(Y, Y).
}
$$
Both relative morphism functors in the right column are algebraic spaces
locally of finite presentation by Quot, Proposition
\ref{quot-proposition-Mor}. The first is separated over $S$ by Lemma
\ref{lemma-Mor-s-lfp}. The assertions follow by base change.

For (2), a proper morphism is universally pure by More on Flatness, Lemma
\ref{flat-lemma-proper-pure}. The result is therefore exactly More on
Flatness, Lemma
\ref{flat-lemma-Weil-restriction-closed-subschemes}.

For (3), write
$$
Z = \underline{\Spec}_Y(\mathcal{A}).
$$
Choose a coequalizer presentation of the quasi-coherent algebra
$\mathcal{A}$ by free commutative algebras
$$
\xymatrix{
\operatorname{Sym}(\mathcal{F}_1)
\ar@<0.5ex>[r] \ar@<-0.5ex>[r] &
\operatorname{Sym}(\mathcal{F}_0) \ar[r] & \mathcal{A}.
}
$$
After arbitrary base change on $S$, algebra maps from $\mathcal{A}$ to the
structure sheaf are the equalizer of the corresponding pair of module-map
functors
$$
\xymatrix{
\mathit{Hom}(\mathcal{F}_0, \mathcal{O}_Y)
\ar@<0.5ex>[r] \ar@<-0.5ex>[r] &
\mathit{Hom}(\mathcal{F}_1, \mathcal{O}_Y).
}
$$
Each of these functors is affine over $S$ by Quot, Proposition
\ref{quot-proposition-hom}: the module $\mathcal{O}_Y$ is flat over $S$ and
has proper support. An equalizer of morphisms of affine $S$-schemes is
affine. Under the additional finite-presentation hypothesis, (1) shows that
the result is of finite presentation.

For (4), a section of $\mathbf{V}(\mathcal{E})_T \to Y_T$ is an
$\mathcal{O}_{Y_T}$-linear map
$\mathcal{E}_T \to \mathcal{O}_{Y_T}$. Apply Quot, Proposition
\ref{quot-proposition-hom}. The explicit construction in Quot, Lemma
\ref{quot-lemma-cohomology-perfect-complex} presents this affine functor,
locally on $S$, by homogeneous linear equations in a finite affine space.
It is consequently $\mathbf{V}(\mathcal{Q})$ for a finitely presented
quasi-coherent $\mathcal{O}_S$-module $\mathcal{Q}$.
\end{proof}

\begin{remark}[The classical field-valued refinement]
\label{remark-sections-proper-field-quasi-projective}
Let $S = \Spec(k)$, let $Y$ be a proper $k$-scheme, and let $Z \to Y$ be
quasi-projective. The argument recorded in
\cite[Expos\'e 221, Section 7]{FGA} uses Chow's lemma and factorization of a
finite morphism to show that $\mathit{Sec}_{Y/k}(Z/Y)$ is a scheme which is
a disjoint union of quasi-projective $k$-schemes. The projective cover in
Chow's lemma is automatically flat over the field. Proposition
\ref{proposition-sections-proper-flat} gives the base-independent
algebraic-space statement; representability by a scheme and
quasi-projectivity are additional conclusions, not part of the general
algebraic-space theorem.
\end{remark}

\begin{remark}[Scheme representability is essential]
\label{remark-Nagata-Weil-restriction-scheme}
\begin{reference}
\cite[Expos\'e 221, Section 7]{FGA}.
\end{reference}
There are a quadratic field extension $k'/k$ and a smooth proper
nonprojective threefold $X$ over $k'$ for which the restriction-of-scalars
functor $\operatorname{Res}_{k'/k}(X)$ is not representable by a scheme.
The corresponding finite \'etale locus in the degree-two Hilbert functor and
the symmetric square likewise need not be schemes.

In the terminology of the cited source, these objects ``do not exist''
because only schemes are admitted as representing objects. As fppf sheaves
they do exist as algebraic spaces: use Criteria for Representability,
Proposition
\ref{criteria-proposition-restriction-of-scalars-algebraic-space}, Quot,
Proposition \ref{quot-proposition-hilb}, and Theorem
\ref{theorem-divided-power-zero-cycles}. Thus the example is an obstruction
to scheme representability, not to modern algebraic-space representability.
\end{remark}

\begin{remark}[A closed graph does not produce a scheme quotient]
\label{remark-FGA-closed-graph-quotient}
The addendum to \cite[Expos\'e 221]{FGA} retracts the quotient conjecture in
Expos\'e 212, Section 8. Even for a smooth variety in characteristic zero,
having a closed action graph does not imply existence or quasi-projectivity
of a scheme quotient. A free action has an fppf quotient algebraic space by
Bootstrap, Lemma \ref{bootstrap-lemma-quotient-free-action}; a scheme or a
quasi-projective quotient requires additional hypotheses. This distinction
is not removed by the Keel--Mori theorem of Morphisms of Algebraic Stacks,
Theorem \ref{stacks-more-morphisms-theorem-keel-mori}.
\end{remark}





\section{Properties of the stack of polarized proper schemes}
\label{section-polarized}

\noindent
In this section we discuss properties of the moduli stack
$$
\Polarizedstack \longrightarrow \Spec(\mathbf{Z})
$$
whose category of sections over a scheme $S$ is the category of
proper, flat, finitely presented scheme over $S$ endowed
with a relatively ample invertible sheaf. This is an algebraic
stack by Quot, Theorem \ref{quot-theorem-polarized-algebraic}.

\begin{lemma}
\label{lemma-polarized-diagonal-separated-fp}
The diagonal of $\Polarizedstack$ is separated
and of finite presentation.
\end{lemma}

\begin{proof}
Recall that $\Polarizedstack$ is a limit preserving algebraic stack, see
Quot, Lemma \ref{quot-lemma-polarized-limits}.
By Limits of Stacks, Lemma \ref{stacks-limits-lemma-limit-preserving-diagonal}
this implies that
$\Delta : \Polarizedstack \to \Polarizedstack \times \Polarizedstack$
is limit preserving. Hence $\Delta$ is locally of finite presentation
by Limits of Stacks, Proposition
\ref{stacks-limits-proposition-characterize-locally-finite-presentation}.

\medskip\noindent
Let us prove that $\Delta$ is separated. To see this, it suffices to show
that given an affine scheme $U$ and two objects
$\upsilon = (Y, \mathcal{N})$ and $\chi = (X, \mathcal{L})$
of $\Polarizedstack$ over $U$, the algebraic
space
$$
\mathit{Isom}_{\Polarizedstack}(\upsilon, \chi)
$$
is separated. The rule which to an isomorphism $\upsilon_T \to \chi_T$
assigns the underlying isomorphism $Y_T \to X_T$ defines a morphism
$$
\mathit{Isom}_{\Polarizedstack}(\upsilon, \chi)
\longrightarrow
\mathit{Isom}_U(Y, X)
$$
Since we have seen in Lemmas \ref{lemma-Mor-s-lfp} and
\ref{lemma-Isom-in-Mor} that the target is
a separated algebraic space, it suffices to prove that this morphism
is separated. Given an isomorphism $f : Y_T \to X_T$
over some scheme $T/U$, then clearly
$$
\mathit{Isom}_{\Polarizedstack}(\upsilon, \chi)
\times_{\mathit{Isom}_U(Y, X), [f]} T
=
\mathit{Isom}(\mathcal{N}_T, f^*\mathcal{L}_T)
$$
Here $[f] : T \to \mathit{Isom}_U(Y, X)$ indicates the $T$-valued
point corresponding to $f$ and
$\mathit{Isom}(\mathcal{N}_T, f^*\mathcal{L}_T)$
is the algebraic space discussed in Section \ref{section-hom-isom}.
Since this algebraic space is affine over $U$, the claim implies
$\Delta$ is separated.

\medskip\noindent
To finish the proof we show that $\Delta$ is quasi-compact. Since
$\Delta$ is representable by algebraic spaces, it suffice to check
the base change of $\Delta$ by a surjective smooth morphism
$U \to \Polarizedstack \times \Polarizedstack$ is quasi-compact
(see for example Properties of Stacks, Lemma
\ref{stacks-properties-lemma-check-property-covering}).
We can assume $U = \coprod U_i$ is a disjoint union of affine opens.
Since $\Polarizedstack$ is limit preserving (see above), we
see that $\Polarizedstack \to \Spec(\mathbf{Z})$ is locally of
finite presentation, hence $U_i \to \Spec(\mathbf{Z})$ is
locally of finite presentation
(Limits of Stacks, Proposition
\ref{stacks-limits-proposition-characterize-locally-finite-presentation}
and Morphisms of Stacks, Lemmas
\ref{stacks-morphisms-lemma-composition-finite-presentation} and
\ref{stacks-morphisms-lemma-smooth-locally-finite-presentation}).
In particular, $U_i$ is Noetherian affine. This reduces us to the
case discussed in the next paragraph.

\medskip\noindent
In this paragraph, given a Noetherian affine scheme $U$ and two objects
$\upsilon = (Y, \mathcal{N})$ and $\chi = (X, \mathcal{L})$
of $\Polarizedstack$ over $U$, we show the algebraic space
$$
\mathit{Isom}_{\Polarizedstack}(\upsilon, \chi)
$$
is quasi-compact. Since the connected components of $U$ are open and closed
we may replace $U$ by these. Thus we may and do assume $U$ is connected.
Let $u \in U$ be a point. Let $P$ be the Hilbert polynomial
$n \mapsto \chi(Y_u, \mathcal{N}_u^{\otimes n})$, see
Varieties, Lemma \ref{varieties-lemma-numerical-polynomial-from-euler}.
Since $U$ is connected and since
the functions $u \mapsto \chi(Y_u, \mathcal{N}_u^{\otimes n})$
are locally constant (see 
Derived Categories of Schemes, Lemma
\ref{perfect-lemma-chi-locally-constant-geometric})
we see that we get the same Hilbert polynomial in every point of $U$.
Set
$\mathcal{M} = \text{pr}_1^*\mathcal{N}
\otimes_{\mathcal{O}_{Y \times_U X}} \text{pr}_2^*\mathcal{L}$
on $Y \times_U X$. Given
$(f, \varphi) \in \mathit{Isom}_{\Polarizedstack}(\upsilon, \chi)(T)$
for some scheme $T$ over $U$ then for every $t \in T$ we have
$$
\chi(Y_t, (\text{id} \times f)^*\mathcal{M}^{\otimes n}) =
\chi(Y_t,
\mathcal{N}_t^{\otimes n} \otimes_{\mathcal{O}_{Y_t}}
f_t^*\mathcal{L}_t^{\otimes n}) =
\chi(Y_t, \mathcal{N}_t^{\otimes 2n}) = P(2n)
$$
where in the middle equality we use the isomorphism
$\varphi : f^*\mathcal{L}_T \to \mathcal{N}_T$.
Setting $P'(t) = P(2t)$ we find that the morphism
$$
\mathit{Isom}_{\Polarizedstack}(\upsilon, \chi)
\longrightarrow
\mathit{Isom}_U(Y, X)
$$
(see earlier) has image contained in the intersection
$$
\mathit{Isom}_U(Y, X) \cap \mathit{Mor}^{P', \mathcal{M}}_U(Y, X)
$$
The intersection is an intersection of open subspaces of
$\mathit{Mor}_U(Y, X)$ (see Lemma \ref{lemma-Isom-in-Mor} and
Remark \ref{remark-Mor-numerical}).
Now $\mathit{Mor}^{P', \mathcal{M}}_U(Y, X)$
is a Noetherian algebraic space as it is of finite
presentation over $U$ by Lemma \ref{lemma-Mor-qc-over-base}.
Thus the intersection
is a Noetherian algebraic space too. Since the morphism
$$
\mathit{Isom}_{\Polarizedstack}(\upsilon, \chi)
\longrightarrow
\mathit{Isom}_U(Y, X) \cap \mathit{Mor}^{P', \mathcal{M}}_U(Y, X)
$$
is affine (see above) we conclude.
\end{proof}

\begin{lemma}
\label{lemma-polarized-qs-lfp}
The morphism $\Polarizedstack \to \Spec(\mathbf{Z})$ is quasi-separated and
locally of finite presentation.
\end{lemma}

\begin{proof}
To check $\Polarizedstack \to \Spec(\mathbf{Z})$ is quasi-separated we have to
show that its diagonal is quasi-compact and quasi-separated.
This is immediate from Lemma \ref{lemma-polarized-diagonal-separated-fp}.
To prove that $\Polarizedstack \to \Spec(\mathbf{Z})$ is locally of finite
presentation, it suffices to show that $\Polarizedstack$
is limit preserving, see Limits of Stacks, Proposition
\ref{stacks-limits-proposition-characterize-locally-finite-presentation}.
This is Quot, Lemma \ref{quot-lemma-polarized-limits}.
\end{proof}

\begin{lemma}
\label{lemma-bounded-polarized}
Let $n \geq 1$ be an integer and let $P$ be a numerical polynomial.
Let
$$
T \subset |\Polarizedstack|
$$
be a subset with the following property: for every $\xi \in T$
there exists a field $k$ and an object $(X, \mathcal{L})$
of $\Polarizedstack$ over $k$ representing $\xi$ such that
\begin{enumerate}
\item the Hilbert polynomial of $\mathcal{L}$ on $X$ is $P$, and
\item there exists a closed immersion $i : X \to \mathbf{P}^n_k$
such that $i^*\mathcal{O}_{\mathbf{P}^n}(1) \cong \mathcal{L}$.
\end{enumerate}
Then $T$ is a Noetherian topological space, in particular quasi-compact.
\end{lemma}

\begin{proof}
Observe that $|\Polarizedstack|$ is a locally Noetherian topological
space, see Morphisms of Stacks, Lemma
\ref{stacks-morphisms-lemma-Noetherian-topology}
(this also uses that $\Spec(\mathbf{Z})$ is Noetherian and
hence $\Polarizedstack$ is a locally Noetherian algebraic stack
by Lemma \ref{lemma-polarized-qs-lfp} and
Morphisms of Stacks, Lemma
\ref{stacks-morphisms-lemma-locally-finite-type-locally-noetherian}).
Thus any quasi-compact subset of $|\Polarizedstack|$ is
a Noetherian topological space and any subset of such is
also Noetherian, see
Topology, Lemmas \ref{topology-lemma-finite-union-Noetherian} and
\ref{topology-lemma-Noetherian}.
Thus all we have to do is a find a quasi-compact subset
containing $T$.

\medskip\noindent
By Lemma \ref{lemma-hilb-proper-over-base} the algebraic space
$$
H =
\Hilbfunctor^{P, \mathcal{O}(1)}_{\mathbf{P}^n_\mathbf{Z}/\Spec(\mathbf{Z})}
$$
is proper over $\Spec(\mathbf{Z})$. By
Quot, Lemma \ref{quot-lemma-extend-hilb-to-spaces}\footnote{We will see
later (insert future reference here) that $H$ is a scheme and hence the
use of this lemma and Quot, Lemma \ref{quot-lemma-extend-polarized-to-spaces}
isn't necessary.} the identity morphism of $H$ corresponds
to a closed subspace
$$
Z \subset \mathbf{P}^n_H
$$
which is proper, flat, and of finite presentation over $H$ and
such that the restriction $\mathcal{N} = \mathcal{O}(1)|_Z$
is relatively ample on $Z/H$ and has Hilbert polynomial $P$
on the fibres of $Z \to H$. In particular, the pair $(Z \to H, \mathcal{N})$
defines a morphism
$$
H \longrightarrow \Polarizedstack
$$
which sends a morphism of schemes $U \to H$ to the classifying morphism
of the family $(Z_U \to U, \mathcal{N}_U)$, see
Quot, Lemma \ref{quot-lemma-extend-polarized-to-spaces}.
Since $H$ is a Noetherian algebraic space
(as it is proper over $\mathbf{Z})$)
we see that $|H|$ is Noetherian and hence quasi-compact. The map
$$
|H| \longrightarrow |\Polarizedstack|
$$
is continuous, hence the image is quasi-compact.
Thus it suffices to prove $T$ is contained
in the image of $|H| \to |\Polarizedstack|$.
However, assumptions (1) and (2) exactly express the fact
that this is the case: any choice of a closed immersion
$i : X \to \mathbf{P}^n_k$ with
$i^*\mathcal{O}_{\mathbf{P}^n}(1) \cong \mathcal{L}$ we get a
$k$-valued point of $H$ by the moduli interpretation of $H$.
This finishes the proof of the lemma.
\end{proof}






\section{Properties of moduli of complexes on a proper morphism}
\label{section-complexes}

\noindent
Let $f : X \to B$ be a morphism of algebraic spaces which is proper,
flat, and of finite presentation. Then the stack
$\Complexesstack_{X/B}$ parametrizing relatively perfect complexes
with vanishing negative self-exts is algebraic. See
Quot, Theorem \ref{quot-theorem-complexes-algebraic}.

\begin{lemma}
\label{lemma-complexes-diagonal-affine-fp}
The diagonal of $\Complexesstack_{X/B}$ over $B$ is affine
and of finite presentation.
\end{lemma}

\begin{proof}
The representability of the diagonal by algebraic spaces
was shown in Quot, Lemma \ref{quot-lemma-complexes-diagonal}.
From the proof we find that we have to show:
given a scheme $T$ over $B$ and objects
$E, E' \in D(\mathcal{O}_{X_T})$ such that
$(T, E)$ and $(T, E')$ are objects of the fibre category
of $\Complexesstack_{X/B}$ over $T$, then
$\mathit{Isom}(E, E') \to T$
is affine and of finite presentation.
Here $\mathit{Isom}(E, E')$ is the functor
$$
(\Sch/T)^{opp} \to \textit{Sets},\quad
T' \mapsto \{\varphi : E_{T'} \to E'_{T'}
\text{ isomorphism in }D(\mathcal{O}_{X_{T'}})\}
$$
where $E_{T'}$ and $E'_{T'}$ are the derived pullbacks of $E$ and $E'$
to $X_{T'}$. Consider the functor $H = \SheafHom(E, E')$ defined
by the rule
$$
(\Sch/T)^{opp} \to \textit{Sets},\quad
T' \mapsto \Hom_{\mathcal{O}_{X_{T'}}}(E_T, E'_T)
$$
By Quot, Lemma \ref{quot-lemma-complexes-open-neg-exts-vanishing}
this is an algebraic space affine and of finite presentation over $T$.
The same is true for $H' = \SheafHom(E', E)$, $I = \SheafHom(E, E)$, and
$I' = \SheafHom(E', E')$. Therefore we see that
$$
\mathit{Isom}(E, E') = (H' \times_T H) \times_{c, I \times_T I', \sigma} T
$$
where $c(\varphi', \varphi) = (\varphi \circ \varphi', \varphi' \circ \varphi)$
and $\sigma = (\text{id}, \text{id})$ (compare with the proof of
Quot, Proposition \ref{quot-proposition-isom}). Thus
$\mathit{Isom}(E, E')$ is affine over $T$ as a fibre product of
schemes affine over $T$. Similarly, $\mathit{Isom}(E, E')$ is
of finite presentation over $T$.
\end{proof}

\begin{lemma}
\label{lemma-complexes-qs-lfp}
The morphism $\Complexesstack_{X/B} \to B$ is quasi-separated and
locally of finite presentation.
\end{lemma}

\begin{proof}
To check $\Complexesstack_{X/B} \to B$ is quasi-separated we have to
show that its diagonal is quasi-compact and quasi-separated.
This is immediate from Lemma \ref{lemma-complexes-diagonal-affine-fp}.
To prove that $\Complexesstack_{X/B} \to B$ is locally of finite
presentation, we have to show that $\Complexesstack_{X/B} \to B$
is limit preserving, see
Limits of Stacks, Proposition
\ref{stacks-limits-proposition-characterize-locally-finite-presentation}.
This follows from Quot, Lemma \ref{quot-lemma-complexes-limits}
(small detail omitted).
\end{proof}





\input{chapters}

\bibliography{my}
\bibliographystyle{amsalpha}

\end{document}
